% concatenated from main.tex
\documentclass[11pt, reqno, a4paper, twoside]{article}
\usepackage[top=2cm,bottom=1.8cm,left=1.3cm,right=1.3cm]{geometry}
\setlength{\parindent}{10pt}
\setlength{\parskip}{3pt}
\usepackage[colorlinks=true,breaklinks=true,linkcolor=lightblue,citecolor=lightblue,urlcolor=lightblue]{hyperref}
\usepackage{newtxtext}
\usepackage[final]{pdfpages}

\usepackage{amsmath,amssymb,amsfonts,mathtools,fancyhdr}
\usepackage{tabularx}
\usepackage{multirow}        
\usepackage{diagbox}
\usepackage{tikz}
\usepackage{algorithm}
\usepackage{algpseudocode}
\usepackage{graphicx,epsfig,scalerel,subfigure,booktabs}
\usepackage[super]{nth}
\usepackage{xparse}
\usepackage{xurl}
\usepackage{hyperref}
\usepackage{siunitx}
\usepackage{afterpage}
\usepackage{enumitem}
\usepackage{comment}
\usepackage{stmaryrd}
\usepackage[normalem]{ulem}
\usepackage{xcolor}

\definecolor{lightgreen}{rgb}{0.22,0.50,0.25}
\definecolor{lightblue}{rgb}{0.22,0.45,0.70}

\newcommand*\tnorm[1]{|\!|\!| #1 |\!|\!|}
\newcommand{\bn}{\boldsymbol{n}}
\newcommand{\br}{\boldsymbol{r}}
\newcommand{\bs}{\boldsymbol{s}}
\newcommand{\bu}{\boldsymbol{u}}
\newcommand{\bv}{\boldsymbol{v}}
\newcommand{\bw}{\boldsymbol{w}}
\newcommand{\bb}{\boldsymbol{b}}
\renewcommand{\bf}{\boldsymbol{f}}
\newcommand{\bg}{\boldsymbol{g}}
\newcommand{\bt}{\boldsymbol{t}}
\newcommand{\bx}{\boldsymbol{x}}
\newcommand{\by}{\boldsymbol{y}}
\newcommand{\bz}{\boldsymbol{z}}
\newcommand\beps{\boldsymbol{\varepsilon}}
\newcommand\bsigma{\boldsymbol{\sigma}}
\newcommand\bzeta{\boldsymbol{\zeta}}
\newcommand\bzero{\boldsymbol{0}}
\newcommand\bkappa{\boldsymbol{\kappa}}
\newcommand\bxi{\boldsymbol{\xi}}
\newcommand\btau{\boldsymbol{\tau}}
\newcommand\bnabla{\boldsymbol{\nabla}}

\newcommand{\rF}{\mathtt{F}}

\newcommand{\rH}{\mathrm{H}}
\newcommand{\rL}{\mathrm{L}}
\newcommand{\rW}{\mathrm{W}}
\newcommand{\rQ}{\mathrm{Q}}

\newcommand{\bB}{\mathbf{B}}
\newcommand{\bH}{\mathbf{H}}
\newcommand{\bX}{\mathbf{X}}
\newcommand{\bR}{\mathbf{R}}
\newcommand{\bL}{\mathbf{L}}
\newcommand{\bW}{\mathbf{W}}
\newcommand{\bV}{\mathbf{V}}
\newcommand{\bQ}{\mathbf{Q}}
\newcommand{\bZ}{\mathbf{Z}}

\newcommand{\bbA}{\mathbb{A}}
\newcommand{\bbD}{\mathbb{D}}
\newcommand{\bbB}{\mathbb{B}}
\newcommand{\bbC}{\mathbb{C}}

\newcommand{\bbJ}{\mathbb{J}}

\newcommand{\cA}{\mathcal{A}}
\newcommand{\cB}{\mathcal{B}}
\newcommand{\cC}{\mathcal{C}}
\newcommand{\cE}{\mathcal{E}}
\newcommand{\cL}{\mathcal{L}}
\newcommand{\cJ}{\mathcal{J}}
\newcommand{\cI}{\mathcal{I}}
\newcommand{\cT}{\mathcal{T}}
\newcommand{\cS}{\mathcal{S}}
\newcommand{\cR}{\mathcal{R}}

\newcommand\bdiv{\mathop{\mathbf{div}}\nolimits}
\newcommand\vdiv{\mathop{\mathrm{div}}\nolimits}
\newcommand\brot{\mathop{\mathbf{rot}}\nolimits}
\newcommand\vrot{\mathop{\mathrm{rot}}\nolimits}
\newcommand\bcurl{\mathop{\mathbf{curl}}\nolimits}
\newcommand\vcurl{\mathop{\mathrm{curl}}\nolimits}
\newcommand\tr{\mathop{\mathrm{tr}}\nolimits}


\newtheorem{remark}{Remark}[section]
\newtheorem{theorem}{Theorem}[section]
\newtheorem{lemma}[theorem]{Lemma}
\newtheorem{proposition}[theorem]{Proposition}
\newtheorem{corollary}[theorem]{Corollary}
\newenvironment{proof}{\noindent{\it Proof.}}{\hfill$\square$}
\allowdisplaybreaks
\numberwithin{equation}{section}
\numberwithin{remark}{section}
\numberwithin{figure}{section}
\numberwithin{table}{section}

%\usepackage[dvipsnames]{xcolor}
\newcommand{\rrbcom}[1]{{\leavevmode\color{lightgreen}{[\textbf{RRB}: #1]}}} 
\newcommand{\akpcom}[1]{{\leavevmode\color{orange}{[\textbf{AKP}: #1]}}}
\newcommand{\hhcom}[1]{{\leavevmode\color{violet}{[\textbfg{HH}: #1]}}}
\definecolor{darkteal}{RGB}{100, 0, 120}
\newcommand{\rdcom}[1]{{\leavevmode\color{BrickRed}{[\textbf{RD}: #1]}}}

\newcommand{\akp}[1]{{\leavevmode\color{orange}{#1}}}
\newcommand{\hh}[1]{{\leavevmode\color{violet}{#1}}}
\newcommand{\cgray}[1]{{\leavevmode\color{gray}{#1}}}
 \newcommand{\cblue}[1]{{\leavevmode\color{blue}{#1}}}
\newcommand{\rrb}[1]{{\leavevmode\color{lightgreen}{#1}}}
\newcommand{\rd}[1]{\textcolor{BrickRed}{#1}}

\begin{document}
    

\pagestyle{fancy}
\rhead{{\small Das, Hutridurga, Pani, Ruiz-Baier}}
\lhead{\textit{\small Robust methods for Darcy--Forchheimer/transport equations}}

\title{Parameter-robust well-posedness and discretisation for coupled Darcy--Forchheimer and advection-diffusion-reaction equations\thanks{\textbf{Updated:} \today. \textbf{Funding:} This work has been partially supported by  the Australian Research Council through the \textit{Future Fellowship} grant FT220100496, and by the Center of Advanced Study (CAS) at the Norwegian Academy of Science and Letters under the program \textit{Mathematical Challenges in Brain Mechanics}. RD acknowledges support from the Industrial Research and Consultancy Centre (IRCC), IIT Bombay (grant no. RI/0519-10001788-001). HH acknowledges financial support from the Anusandhan National Research Foundation (ANRF), Government of India (grant no. ANRF/ARG/2025/011989/MS).}}




\author{Rishi Das\thanks{
IITB–Monash Research Academy, Indian Institute of Technology Bombay, Powai, Mumbai, Maharashtra 400076, India. Email: rishi.das@monash.edu.} 
\and  Harsha Hutridurga\thanks{Department of Mathematics, Indian Institute of Technology Bombay, Powai, Mumbai, Maharashtra 400076, India. Email: hutridurga@iitb.ac.in.}
\and Amiya K. Pani\thanks{Department of Mathematics, BITS-Pilani, K.K. Birla Goa Campus, Zuarinagar, Goa-403726, India. Email: amiyap@goa.bits-pilani.ac.in.}
\and Ricardo Ruiz-Baier\thanks{School of Mathematics, Monash University, 9 Rainforest Walk, Melbourne VIC 3800, Australia. Email: ricardo.ruizbaier@monash.edu.}}
\date{\today}



\maketitle

\begin{abstract}
\noindent 
We adapt the recent theory for unique solvability of perturbed saddle-point problems in Banach spaces to the case of parameter-independent stability bounds. This also constitutes an extension to the Brezzi--Braess theory for parameter-robust stability of perturbed saddle-point problems from the Hilbert to the Banach setting. We apply the new abstract result to the mixed formulation of the advection-diffusion-reaction equation and tackle also its coupling with the Darcy--Forchheimer equations. The complete system is shown to be stable irrespective of the model parameters of permeability, Forchheimer coefficient,   and reaction modulation. We discretise the problem with mixed finite element methods and show convergence in appropriately parameter-weighted norms. We also design operator-based preconditioners for the full system, utilising weighted norms that provide robustness with respect to {some model parameters}. 
%except {for} the diffusivity coefficient,
%and
%noting that for the advection-diffusion-reaction equation, the Schur complement-based preconditioner involves the inverse of a primal DG advection-diffusion-reaction block. 
We also present a few numerical examples that show the {convergence and robustness of the proposed method} in some regimes of practical interest.

\medskip
\noindent\textbf{MSC (2010):} 65N30, 65N12, 65N15.

\medskip 
\noindent\textbf{Keywords:} Nonlinear flow in porous media, robust well-posedness in Banach spaces, mixed finite element methods, robust stability, fixed-point analysis, preconditioners, numerical experiments.
\end{abstract}

%%%%%%%%%%%%%%%%%%%%%%%%%%%%%%%%%%%%%
\section{Introduction} 
This paper investigates the robust parameter stability of a system of the Darcy--Forchheimer equations coupled with the advection-diffusion-reaction (ADR) equation as presented in \cite{sayah2021finite}.  
  Let $\Omega\subset \mathbb{R}^d$, $d\in \{2,3\}$, denote a  {bounded} domain with Lipschitz-continuous boundary $\partial\Omega$, on which its outward unit normal  is
denoted by $\bn$.  The boundary is decomposed into  two sub-boundaries $\Gamma_{\bu}$ and $\Gamma_p$, both with positive $(d-1)-$Hausdorff measure. The Darcy--Forchheimer equations are written in their discharge velocity and pressure head formulation $\bu,p$ featuring a nonlinear drag term, and the ADR transport using a conservative form in terms of total (diffusive plus advective) flux  $\bzeta$ and solute concentration  $\varphi$: 
\begin{subequations}\label{eq:strong}
\begin{align}
    \kappa^{-1}\bu + {\rF|\bu|\bu} + \nabla p & = \bf(\varphi) \quad \text{in $\Omega$},\\
    \vdiv \bu & = 0 \quad \text{in $\Omega$},\\
    \alpha^{-1}\bzeta & = \nabla \varphi - { \alpha^{-1}} \bu \varphi \quad \text{in $\Omega$},\\
    \varrho_0 \varphi - \vdiv \bzeta & = m \quad \text{in $\Omega$},\\
    \bu\cdot \bn = \varphi & = 0 \quad \text{on $\Gamma_{\bu}$},\\
    p  = \bzeta \cdot \bn  & = 0 \quad \text{on $\Gamma_p$}.
\end{align}\end{subequations}
Here, $\kappa$ is the hydraulic conductivity (permeability of the porous matrix divided by the fluid viscosity),
% , assuming that $\kappa$ is uniformly bounded and uniformly positive definite, i.e., there exist $\kappa_m, \kappa_M > 0$ such that
% \(
% 0 < \kappa_m \le \|\kappa^{-1}\|_{L^\infty(\Omega)} \le \kappa_M,
% \)
% Without loss of generality, we assume that
% \begin{align}\label{est:kappa}
% 1 \le \|\kappa^{-1}\|_{L^\infty(\Omega)}.
% \end{align}} 
$\bf$ is a forcing vector that depends on the concentration field which satisfies the following condition (cf. \cite[Assumption~2.1]{sayah2021finite}) 
 % \begin{align*}  \bf(\bx,\eta) = \bf_{\bzero}(\bx)+\bf_1(\eta),\end{align*} where  \(\bf_{\bzero}\in\bL^2(\Omega)\) and \(\bf_1\) is a Lipschitz-continuous function with positive constant \(C_{\bf_1}\) and \(\bf_1(0)={\boldsymbol{0}}\) satisfying the inequality 
\begin{align}\label{eq:lip.bound.f}\|\bf(\varphi)\|_{0,\Omega}\leq \|\bf(0)\|_{0,\Omega}+ \alpha_1\|\varphi\|_{0,6,\Omega},\end{align} 
where  \(\alpha_1\) is a positive constant.
 In addition,  $\rF>0$ is the Forchheimer coefficient, $\alpha>0$ is the diffusivity coefficient, $\varrho_0>0$ is the reaction coefficient  and $m$ is the concentration source/sink. A variety of discretisations with finite differences and finite element methods can be found in, e.g. \cite{deugoue2022numerical,huang2025numerical,bonetti2025conforming}, or in \cite{allendes2023numerical} for the case of singular sources (see also \cite{caucao2021fully,gatica2022p} for analysis of similar models in the fully mixed case). Such problems find applications in reservoir modeling and geothermal energy optimisation, for example.

\smallskip
\noindent\textbf{Main contribution.}  
Owing to the form of \eqref{eq:strong}, and  with the aim of preserving both robustness of the parameters--both at the continuous and discrete levels--the variational formulation of the problem is cast involving sums and intersections of reflexive Banach spaces. The main tool in the theoretical analysis of well-posedness (as well as error estimates) is a new abstract result for perturbed saddle-point problems in reflexive Banach spaces that stands an extension of \cite{correa2022continuous} (to the parameter-robust context) as well as of \cite{boon2021robust} (to its non-Hilbert counterpart). 
 A more closely related reference that deals with the design of preconditioning for flow-transport couplings, is the recent work \cite{cai2026parameter} addressing robust stability of a linearised model for coupled heat and Biot equations. The main differences here are that, due to the presence of the nonlinear terms and Banach spaces functional structure, we need to take into consideration sums of three non-Hilbert spaces and realise the action of duality maps that are non-standard when compared to the recent literature on operator preconditioning (see, e.g., \cite{baerland2020observation,cai2026parameter,hong2022robust}). It should be noted that, in the context of ADR problems, the presence of the advective term renders the underlying saddle-point structure non-symmetric. The treatment of this term within the design of preconditioners remains challenging for both preconditioning strategies and iterative solvers. 
 %The preconditioners  we explore here are  inspired by Schur complement and operator preconditioning frameworks for symmetric elliptic operators; however, their extension to the nonsymmetric and nonnormal setting lacks a firm theoretical basis in the literature (see \cite{wathen2015preconditioning,ernst2000residual} for a comprehensive discussion). Another crucial point regarding our robust analysis is that, while the boundedness and inf-sup stability of the individual subproblems~Darcy--Forchheimer and the reaction-diffusion systems are fully uniform with respect to their model   parameters and the mesh size, the  robustness with respect to $\alpha$ and $\rF$ is partially lost in the fully coupled problem due to the coupling of the nonlinear Darcy--Forchheimer velocity and the transport flux through the fixed-point analysis.  Although this loss of robustness is not reflected in the observed convergence rates, such discrepancy could be a consequence of   an inadequate  scaling in the weighted norms associated with the Darcy--Forchheimer velocity and transport flux coupling, as well as from limitations of the fixed-point framework.
 
{The preconditioners  we explore here are  inspired by Schur complement and operator preconditioning frameworks for symmetric elliptic operators; however, their extension to the nonsymmetric and nonnormal setting lacks a firm theoretical basis  \cite{wathen2015preconditioning,ernst2000residual}. Additionally, while the individual Darcy--Forchheimer and reaction-diffusion subproblems are uniformly bounded and inf-sup stable, robustness with respect to $\alpha$ and $\rF$ is partially lost in the coupled system due to the fixed-point coupling of the nonlinear velocity and transport flux. This analytical loss of robustness does not manifest in our observed numerical convergence rates, suggesting it may be an artifact of sub-optimal scaling in the weighted norms or limitations inherent to the fixed-point framework.}


 
\smallskip
\noindent\textbf{Common notation and preliminaries.} For $s\geq 0$ and $r\in[1,+\infty]$, we denote by $\rL^r(\Omega)$ and $\rW^{s,r}(\Omega)$ the usual Lebesgue and Sobolev spaces endowed with the norms $\|\bullet\|_{0,r,\Omega}$ and $\|\bullet\|_{s,r,\Omega}$, respectively.
Note that $\rW^{0,r}(\Omega)=\rL^r(\Omega)$. 
If $r = 2$, we write $\rH^{s}(\Omega)$ in place of $\rW^{s,2}(\Omega)$, and denote the corresponding norm by $\|\bullet\|_{s,\Omega}$. {Recall also the \(\rL^r(\Omega)\)-\(\rL^s(\Omega)\) embedding for finite measure domains, i.e.,   \(\|\bullet\|_{0,r,\Omega}\leq|\Omega|^{\frac1r-\frac1s}\|\bullet\|_{0,s,\Omega},\) for \(1<r\leq s<\infty.\)} Also, $C_{\mathsf{emb}}$ denotes the constants arising from a Sobolev embedding, which depends only on $|\Omega|.$ We 
denote by $\mathbf{H}$ the vectorial counterpart of a generic scalar functional space $\rH$. The $\rL^2(\Omega)$ inner product for scalar, vector, or tensor valued functions
is denoted by $(\bullet,\bullet)_{0,\Omega}$. 
% We also recall the Hilbert space $\bH(\vdiv,\Omega)=\{ \bv \in \bL^2(\Omega): \vdiv \bv \in \rL^2(\Omega)\}$ endowed with the norm $\|\bv\|_{\vdiv,\Omega}^2 = \|\bv\|^2_{0,\Omega} + \|\vdiv \bv\|^2_{0,\Omega}$. 
Moreover, we  use for $r,s \in [1,+\infty)$, the Banach space 
$\bH^r(\vdiv_s,\Omega)=\{ \bv \in \bL^r(\Omega): \vdiv \bv \in \rL^s(\Omega)\},$ endowed with the norm $\|\bv\|_{r,\vdiv_s,\Omega}^2 = \|\bv\|^r_{0,r,\Omega} + \|\vdiv \bv\|^s_{0,s,\Omega}$, {and in particular, for $r=2, s=2,$ we call it $\bH(\vdiv,\Omega)$.}
Often, we  use a subscript $\Gamma_i$ on the functional space to denote that the (appropriate) traces vanish on the part of the boundary $\Gamma_i\subset \partial \Omega$. 

For two Banach spaces $(X,\|\bullet\|_X)$ and $(Y,\|\bullet\|_Y),$ we denote by $\cL(X,Y)$ the space of bounded linear maps from $X$ to $Y$. Let $(Z,\|\bullet\|_Z)$ be another Banach space. The intersection and sum spaces are 
\begin{gather*} X \cap Y \cap Z = \{ v : v\in X \ \text{and} \ v\in Y \ \text{and}\ v\in Z\}, \\ X+ Y + Z = \{u+v+w: u\in X \ \text{and}\ v\in Y\ \text{and} \ w \in Z\},\end{gather*}
and their respective norms are (see \cite{baerland2020observation}, noting that the same norm characterisation done there for Hilbert spaces holds also for the Banach case)
\[ \| r\|_{X\cap Y \cap Z} = \|r\|_X + \|r\|_Y + \|r\|_Z, \quad \| z\|_{X + Y + Z} = \inf_{\substack{r = u+v+w, \\ u \in X, v\in Y, w\in Z}}\,\bigl[\|u\|_X + \|v\|_Y+ \|w\|_Z\bigr].\]
Furthermore, if $X \cap Y \cap Z$ is dense in $X$ and $Y$ and $Z$, then the following relation (understood as an identification under an isometry) holds 
\begin{equation}\label{eq:duality} (X+Y+Z)' = X'\cap Y' \cap Z'.\end{equation}
For a fixed $\alpha>0$, by $\alpha X$ we denote the Banach space whose elements coincide with those in $X$ but are measured with the norm 
\(\| \bullet \|^2_{\alpha X}= \alpha^2\|\bullet \|^2_X. \)
{For any scalar function \(\varphi,\) we recall that the sign function (denoted by \(\mathrm{sgn}(\varphi)\)) takes the value \(1,\) for non-negative \(\varphi,\) and \(-1\), otherwise. 
Additionally, \(\varphi\, \mathrm{sgn}(\varphi) = |\varphi|.\) 

\smallskip
\noindent\textbf{Outline.} The remainder of the paper is organised as follows. Section~\ref{sec:abstract} presents the weak formulation for \eqref{eq:strong}, specifying the required parameter-weighted spaces. Also we state and prove an abstract theory for robust stability of perturbed saddle-point problems in Banach spaces. In Section~\ref{sec:existence}, we formulate our problem in the context of abstract theory and show, using fixed-point arguments, that there exists a unique solution of the coupled system and that we have a parameter-independent continuous dependence on data. Section~\ref{sec:fe} deals with the finite element formulation, the robust discrete stability--via a discrete fixed-point argument--and the derivation of robust a priori error estimates. We construct variable block-diagonal operator preconditioners in Section \ref{sec:precond} and show their robustness with respect to all model parameters and the discretisation parameter,  except the diffusivity coefficient. Finally, Section~\ref{sec:numer} includes a series of numerical tests showing the  convergence of the method in 2D, also demonstrating the performance of the  preconditioners. 


%%%%%%%%%%%%%%%%%%%%%%%
\section{Mixed weak formulation and abstract setting}\label{sec:abstract}
We now consider the weighted spaces proposed in \cite[Section~3]{dhpr_sinum26} for Darcy velocity and pressure: 
\[\bV:=  \kappa^{-\frac{1}{2}}\bL^2(\Omega) \cap \bH_{\Gamma_{\bu}}(\vdiv,\Omega) \cap \rF^{\frac13} \bL^3(\Omega), \quad\text{and}\quad \rQ := \rL^2(\Omega) + \kappa^{\frac{1}{2}}\rH^1_{\Gamma_p}(\Omega) + \rF^{-\frac13}\rW^{1,\frac32}_{\Gamma_p}(\Omega),\]
endowed with the following norms 
\begin{subequations}\label{eq:norms}
    \begin{align}
   \|\bv\|_{\bV} &:= \kappa^{-\frac12}\|\bv\|_{0,\Omega} + \|\vdiv\bv\|_{0,\Omega} + \rF^{\frac13}\|\bv\|_{0,3,\Omega}, \quad\text{and}\\
   % \|q\|_{\rQ} &:= \!\!\inf_{\substack{q = q_1+q_2+q_3, \\ (q_1,q_2,q_3) \in \rL^2(\Omega) \times \rH^1_{\Gamma_{\!p}}\!(\Omega)\times \rW^{1,\frac32}_{\Gamma_p}(\Omega)}}\!\!\bigl[\|q_1\|_{0,\Omega} + \kappa^{\frac{1}{2}}\|q_2\|_{1,\Omega} + \rF^{-\frac13} \|q_3\|_{1,\frac32,\Omega}\bigr],\label{norm:Q}
   {\|q\|_{\rQ}}&:=\inf_{\substack{s\in\rH^1_{\Gamma_p}}}\bigl(\|q-s\|_{0,\Omega}+\kappa^{\frac12}\|\nabla q\|_{0,\Omega}+\rF^{-\frac13}\|\nabla q\|_{0,\frac32,\Omega}\bigr),\label{def:Qhat-norm}
    \end{align}
\end{subequations}
respectively. Denote as follows the trial and test function spaces for flux and concentration (following \cite[Section 3]{boon2021robust}) 
\[ \bW := \alpha^{-\frac12} \bL^2(\Omega) \cap \bH_{\Gamma_p}(\vdiv_{\frac65},\Omega), \quad \text{and}\quad\Psi := \alpha^{\frac12}\rH^1_{\Gamma_{\bu}}(\Omega) + \rL^6(\Omega), \]
equipped with the norms 
\begin{subequations}
\begin{align}\|\bw\|_{\bW} &:= \alpha^{-\frac12}\|\bw\|_{0,\Omega} + \| \vdiv \bw\|_{0,\frac65,\Omega},\quad\text{and} \\
\|\psi\|_{\Psi} &:= \inf_{s \in \alpha^{\frac12}\rH^1_{\Gamma_{\bu}}(\Omega)}\bigl(\|\psi-s\|_{0,6,\Omega} + \alpha^{\frac12}\|\nabla s\|_{0,\Omega} \bigr), \label{eq:Psi.norm}\end{align}\end{subequations}
respectively. The weak formulation of \eqref{eq:strong} consists in finding  $(\bu,p,\bzeta,\varphi)\in \bV\times \rQ\times \bW \times \Psi$ such that 
\begin{subequations}\label{eq:weak}
    \begin{align}
    a_1(\bu,\bv) + d_1(\bu; \bu,\bv) +  b_1(\bv,p) & = F^{\varphi}(\bv) \quad \forall \bv \in \bV,\label{eq:weak,1}\\
    b_1(\bu,q) & = 0 \quad \forall q \in \rQ,\label{eq:weak,2}\\
    a_2(\bzeta,\bw) + b_2(\bw,\varphi) + d_2^{\bu}(\bw, \varphi) & = 0 \quad \forall \bw \in \bW, \label{eq:weak,3}\\
    b_2(\bzeta, \psi) - c_2(\varphi,\psi) & = H(\psi) \quad \forall \psi \in \Psi, \label{eq:weak,4}
    \end{align}
\end{subequations}
where the  bilinear forms $a_1:\bV\times\bV\to\mathbb{R}$, $b_1:\bV\times\rQ \to \mathbb{R}$,  $a_2:\bW\times\bW\to\mathbb{R}$, $b_2:\bW\times\Psi \to \mathbb{R}$, $c_2: \Psi \times \Psi \to \mathbb{R}$,  $d_2^{\hat{\bu}}: \bW \times \Psi \to \mathbb{R}$ (for a fixed $\hat{\bu}$),  
linear functionals $F^{\hat{\varphi}}:\bV\to\mathbb{R}$ (for a fixed $\hat{\varphi}$), $H:\Psi \to \mathbb{R}$, and nonlinear form $d_1:\bV\times\bV\times\bV \to \mathbb{R}$, are defined as: 
\begin{gather*}
  a_1(\bu,\bv):= \int_\Omega \kappa^{-1}\bu\cdot\bv, \quad {b_1(\bv,q):= -\int_\Omega q\,\vdiv\bv},  \quad d_1(\bw;\bu,\bv):=\int_\Omega \rF |\bw|\bu\cdot\bv, \\ 
  a_2(\bzeta,\bw):= \int_\Omega \alpha^{-1}\bzeta\cdot\bw, \quad  b_2(\bw,\psi):= \int_\Omega \psi\,\vdiv\bw, \quad   c_2(\varphi,\psi):=\int_\Omega { \varrho_0} \varphi\,\psi, \\  
  d_2^{\hat{\bu}}(\bw,\psi):={\int_\Omega  \alpha^{-1}\hat{\bu}\cdot \bw \,\psi},\quad 
  F^{\hat{\varphi}}(\bv):= \int_\Omega \bf(\hat{\varphi})\cdot\bv, \quad H(\psi) := - \int_\Omega m\,\psi.   
\end{gather*}
 

%%%%%%%%%%%%%%%%%%%%%%%%%%%%%%%%%%%%%%%%%%%%%%%%%%%%
%\section{Abstract setting for robust well-posedness}
%Next, we adapt the theory of \cite{boon2021robust,braess1996stability} to the case of Banach spaces. 
Let $\bX$ and $Y$ be reflexive Banach spaces equipped with the norms $\|\bullet\|_{\bX}$ and $\|\bullet\|_Y$, respectively. Further, let \(Y_c\) be a dense linear subspace of \(Y.\) Consider the problem of finding $(\bsigma,u) \in \bX\times Y$ such that 
\begin{subequations}\label{eq:abstract}
    \begin{align}
    a(\bsigma,\btau) + b(\btau,u) & = F(\btau) \quad \forall \btau \in \bX, \label{eq:weak.1}\\
    b(\bsigma, v) - t^2c(u,v) & = G(v)\label{eq:weak.2} \quad \forall v\in {Y_c},
\end{align}\end{subequations}
where $0\leq t\leq 1$ is a fixed scaling parameter and  $a:\bX\times \bX \to \mathbb{R}$, $b:\bX\times Y \to \mathbb{R}$, and $c:{Y_c} \times {Y_c} \to \mathbb{R}$ are given bilinear forms. Here, also $(F,G) \in \bX'\times Y'$ are given linear forms.} Let $\bB:\bX\to Y'$ be a bounded linear map induced by $b(\bullet,\bullet)$ as $\bX \ni \btau\mapsto \bB\btau$ with $\langle \bB \btau, v\rangle := b(\btau,v)$ for all $v\in Y$.

We extend the Brezzi--Braess theorem (\cite[Theorem~2.1]{boon2021robust}) to the case of reflexive Banach spaces. For this we adapt  \cite[Theorem~49.15]{ern2021galerkin} to the case of perturbed saddle-point problems. The invertibility of $\cA$ will follow as a consequence of the following conditions
%\rrb{wording is misleading. The BNB theorem applies for problems written with a general $\cA$. They can be perturbed too. So I don't think we are extending BNB.}
%We need to verify the following conditions to achieve the unique solvability of \eqref{eq:abstract}
\begin{subequations}
   \begin{align}
\mathrm{BNB~1:}\qquad&\inf_{\substack{(\bsigma,u)\in\bX\times Y \\ (\bsigma,u)\neq(\bzero,0)}}
\ \sup_{\substack{(\btau,v)\in\bX\times Y \\ (\btau,v)\neq(\bzero,0)}}
\frac{\cA(\bsigma,u;\btau,v)}
     {\tnorm{(\bsigma,u)}_{\bX\times Y}\,
      \tnorm{(\btau,v)}_{\bX\times Y}}
 \geq C_{\natural} \quad \forall (\bsigma,u),(\btau,v)\in\bX\times Y,\\
\mathrm{BNB~2:}\qquad& \forall(\bsigma,u)\in\bX\times Y, \quad [\forall (\btau,v)\in\bX\times Y,\, \cA(\bsigma,u;\btau,v)=0]\implies[(\bsigma,u)={(\bzero, 0)}]\label{eq:BNB2},
\end{align}\end{subequations}
%In order to complete the proof we have to show that the two conditions~\(\mathrm{BNB~1}\) and \(\mathrm{BNB~2}\) of \cite[Theorem~25.9]{ern2021galerkin} 
where \(C_{\natural}\) is a positive constant independent of \(t\) and {\(\tnorm{(\btau,v)}_{\bX\times Y}:=\|\btau\|_{\bX}+\|v\|_{Y}+t|v|_c,\) and \(|\bullet|^2_c={c(\bullet,\bullet)}\)}.
%
The main result is discussed below. 

\begin{theorem}[Robust well-posedness in Banach spaces]\label{th:abstract}
 Assume that  
\begin{itemize}
\item $a(\bullet,\bullet)$, $b(\bullet,\bullet)$, and $c(\bullet,\bullet)$  are bounded with {continuity constants} \(\|a\|, \|b\|\) and \(\|c\|,\) respectively, independent of the scaling parameter $t$. 
    \item $a(\bullet,\bullet)$ and $c(\bullet,\bullet)$ are symmetric and positive semi-definite.
    \item Inf-sup condition in $\bZ = \mathrm{ker}(\bB)$: there exists $C_a>0$, independent of the parameter $t$, such that 
    \begin{align}\label{eq:inf-sup.a} \sup_{\btau \in \bZ\setminus\{\boldsymbol{0}\}} \frac{a(\bsigma,\btau)}{\|\btau\|_{\bX}} \geq C_a \| \bsigma\|_{\bX} \quad \forall \bsigma \in \bZ.\end{align}
    \item Inf-sup condition: there exists $C_b>0$, independent of $t$, such that 
    \begin{align}\label{eq:inf-sup.b} \sup_{\btau \in \bX\setminus\{\boldsymbol{0}\}} \frac{b(\btau,v)}{\|\btau\|_{\bX}} \geq C_b \| v\|_Y \quad \forall v \in Y.\end{align}    
    \item Braess condition: there exists $C_c>0$, independent of $t$, such that 
    \begin{align}\label{eq:Braess}\sup_{(\btau,v) \in \bX\times Y \setminus\{(\boldsymbol{0},0)\}} \frac{a(\bsigma,\btau)+b(\bsigma,v)}{\tnorm{(\btau,v)}_{\bX\times Y}} \geq C_c \| \bsigma \|_{\bX} \quad \forall \bsigma \in \bX.\end{align} 
    % where \(\tnorm{(\bzeta,v)}_{\bX\times Y}:=\|\bzeta\|_{\bX}+\|v\|_{Y}+t|v|_c,\) and \(|\bullet|^2_c={c(\bullet,\bullet)}\).
\end{itemize}
Then, for each pair $(F,G) \in \bX'\times Y'$, the problem \eqref{eq:abstract} is uniquely solvable. Moreover, there holds 
\[\tnorm{(\bsigma,u)}_{\bX\times Y} \leq C\bigl[ \|F\|_{\bX'} + \|G\|_{Y'}\bigr], \]
where \[C:=\max\bigl[\frac{1}{MC_c},6\beta_1^{-1}+9C_b^2\beta_1^{-2}\bigr],\]  
 does not depend on the scaling parameter \(\,t\), and \(M={\bigl({1+\frac{2\|a\|}{C_b}}\bigr)^{-1}} \text{ and } \beta_1=\min\{{C_b,\|a\|}\}\). 
\end{theorem}
\begin{proof}
We follow closely  the analysis in \cite[Theorem~49.15]{ern2021galerkin} and \cite[Lemma 3]{braess1996stability}. 
% Since \(\bX\) is a reflexive Banach space and \(\bZ\) is a closed subspace of \(\bX,\) it follows that \(\bX\) can be decomposed inequality as the direct sum of \(\bZ\) and \(\bZ^c\), i.e., \(\bX=\bZ\oplus\bZ^c.\) 
First, we define \begin{align}\label{eq:curl.A}\mathcal{A}(\bsigma,u;\btau,v):=a(\bsigma,\btau)+b(\btau,u)+b(\bsigma,v)-t^2c(u,v), \quad \forall (\bsigma,u),(\btau,v)\in \bX \times Y,\end{align}
% Since \[\sup_{(\bzero,0)\neq(\btau,v)\in{\bX\times Y}}\frac{\mathcal{A}(\bsigma,u;\btau,v)}{\tnorm{(\btau,v)}_{\bX\times Y}}\leq\sup_{(\bzero,0)\neq(\btau,v)\in{\bX\times Y}}\frac{F(\btau)+G(v)}{\|\btau\|_{\bX}+\|v\|_Y}\leq\|F\|_{\bX'}+\|G\|_{Y'},\]
%  for a proof of \(\tnorm{(\bsigma,u)}_{\bX\times Y} \leq  C\bigl[\|F\|_{\bX'} + \|G\|_{Y'}\bigr],\) it is enough to show that \begin{align}\label{eq:inf-sup.A}\tnorm{(\bsigma,u)}_{\bX\times Y}\leq C\sup_{(\bzero,0)\neq(\btau,v)\in{\bX\times Y}}\frac{\mathcal{A}(\bsigma,u;\btau,v)}{\tnorm{(\btau,v)}_{\bX\times Y}}.\end{align}
%  Also, note that \eqref{eq:inf-sup.A} provides the inf-sup stability of \(\mathcal{A}(\bullet,\bullet;\bullet,\bullet)\) on the space \(\bX\times Y,\) and since all the bilinear forms \(a(\bullet,\bullet), b(\bullet,\bullet)\) and \(c(\bullet,\bullet)\) are bounded, the boundedness of \(\mathcal{A}(\bullet,\bullet;\bullet,\bullet)\), i.e., \(\mathcal{A}(\bullet,\bullet;\bullet,\bullet)\lesssim\tnorm{(\bullet,\bullet)}_{\bX\times Y}\) trivially holds. Hence, once we establish \eqref{eq:inf-sup.A}, both existence-uniqueness and the continuous dependence of data are achieved. \akp{Give a proper reference }: \rrb{yes, we need to invoke the  Banach--Ne\v{c}as--Babu\v{s}ka theory (generalised Lax--Milgram)} 
and then, we focus on proving \(\mathrm{BNB~1}\), for which we  consider two cases. 
 
\smallskip 
\noindent\textbf{Case~1.} Let us first assume that \begin{align}\label{eq:case.1}\|\bsigma\|_{\bX}\leq \frac{C_b}{2\|a\|}\{\|u\|_Y+t|u|_c\}.\end{align}
Re-writing \(\tnorm{(\bullet,\bullet)}_{\bX\times Y} \) in the following form, with \(\beta_1=\min\{{C_b,\|a\|}\},\) and using \eqref{eq:case.1}, we obtain 
\begin{align}\label{eq:case.1.a}\tnorm{(\bsigma,u)}_{\bX\times Y}=\|\bsigma\|_{\bX}+\|u\|_{Y}+t|u|_c\nonumber&\leq \beta_1^{-1}\bigl(\|a\|\|\bsigma\|_{\bX}+C_b\|u\|_Y+C_bt|u|_c\bigr)\nonumber\\
&\leq {\frac32C_b\beta_1^{-1}}\{\|u\|_Y+t|u|_c\}.
\end{align}
Note that the positive semi-definiteness of \(a(\bullet,\bullet)\) yields \begin{align}\label{eq:case.1.b}
t|u|_c\leq\bigl[a(\bsigma,\bsigma)+t^2|u|_c^2\bigr]^{\frac12}=\big(\mathcal{A}(\bsigma,u;\bsigma,-u)\Big)^{\frac12}.\end{align}
We continue with the inf-sup condition of \(b(\bullet,\bullet)\)  and use \eqref{eq:curl.A}. Then, we use the boundedness of \(a(\bullet,\bullet)\) and \eqref{eq:case.1} to derive the following upper bound of \(\|u\|_Y\) 
\begin{align}\label{eq:case.1.c}
C_b\|u\|_Y\leq\sup_{\bzero\neq\btau\in\bX}\frac{b(\btau,u)}{\|\btau\|_{\bX}}&=\sup_{\bzero\neq\btau\in\bX}\frac{\mathcal{A}(\bsigma,u;\btau,0)-a(\bsigma,\btau)}{\|\btau\|_{\bX}}\nonumber\\
&\leq\sup_{(\bzero,0)\neq(\btau,v)\in{\bX\times Y}}\frac{\mathcal{A}(\bsigma,u;\btau,v)}{\tnorm{(\btau,v)}_{\bX\times Y}}+\|a\|\|\bsigma\|_{\bX}\nonumber\\
&\leq\sup_{(\bzero,0)\neq(\btau,v)\in{\bX\times Y}}\frac{\mathcal{A}(\bsigma,u;\btau,v)}{\tnorm{(\btau,v)}_{\bX\times Y}}+\frac{C_b}{2}\{\|u\|_Y+t|u|_c\}.
\end{align}
Note that, \eqref{eq:case.1.c} provides a bound on \({\|u\|_{Y}}\) once one subtracts \(\frac{C_b\|u\|_Y}{2}\) from both sides first, and then multiplies both sides by \({2}{C_b}^{-1}\).
%
Invoking \eqref{eq:case.1.b}-\eqref{eq:case.1.c} in \eqref{eq:case.1.a}, we obtain %the following 
\begin{align}\label{eq:case2.bound}\tnorm{(\bsigma,u)}_{\bX\times Y}&\leq 3\beta_1^{-1}\sup_{(\bzero,0)\neq(\btau,v)\in{\bX\times Y}}\frac{\mathcal{A}(\bsigma,u;\btau,v)}{\tnorm{(\btau,v)}_{\bX\times Y}}+3C_b\beta_1^{-1}t|u|_c\nonumber\\
&\leq 3\beta_1^{-1}\sup_{(\bzero,0)\neq(\btau,v)\in{\bX\times Y}}\frac{\mathcal{A}(\bsigma,u;\btau,v)}{\tnorm{(\btau,v)}_{\bX\times Y}}+3C_b\beta_1^{-1}\mathcal{A}(\bsigma,u;\bsigma,-u)^{\frac12}.
\end{align}
%Let us now use an application of the Young's inequality, which says that, 
We recall that for any three positive real numbers \(x, y, z\) such that \(x\leq y+z\), an application of the Young's inequality shows \(x\leq2y+\frac{z^2}{x}\) (see, \cite[Equation~3.4]{braess1996stability}). Choose \[x=\tnorm{(\bsigma,u)}_{\bX\times Y},\, y = 3\beta_1^{-1}\sup_{(\bzero,0)\neq(\btau,v)\in{\bX\times Y}}\frac{\mathcal{A}(\bsigma,u;\btau,v)}{\tnorm{(\btau,v)}_{\bX\times Y}} \quad \text{ and }\quad  z= 3C_b\beta_1^{-1}\mathcal{A}(\bsigma,u;\bsigma,-u)^{\frac12},\] in this context, %and we 
to obtain the following stability bound 
\begin{align*}
\tnorm{(\bsigma,u)}_{\bX\times Y}&\leq 6\beta_1^{-1}\sup_{(\bzero,0)\neq(\btau,v)\in{\bX\times Y}}\frac{\mathcal{A}(\bsigma,u;\btau,v)}{\tnorm{(\btau,v)}_{\bX\times Y}}+9C_b^2\beta_1^{-2}\frac{\mathcal{A}(\bsigma,u;\bsigma,-u)}{\tnorm{(\bsigma,u)}_{\bX\times Y}}\\
&\leq(6\beta_1^{-1}+9C_b^2\beta_1^{-2})\sup_{(\bzero,0)\neq(\btau,v)\in{\bX\times Y}}\frac{\mathcal{A}(\bsigma,u;\btau,v)}{\tnorm{(\btau,v)}_{\bX\times Y}}.
\end{align*}

\smallskip 
\noindent
\textbf{Case~2.} On the other hand, we, now, assume that 
\begin{align}\label{eq:case.2}\|\bsigma\|_{\bX}\geq \frac{C_b}{2\|a\|}\{\|u\|_Y+t|u|_c\}.\end{align}
% Let us consider the unique decomposition of \(\bsigma=\bsigma_{\bzero}+\widetilde{\bsigma},\) where \(\bsigma_{\bzero}\in \bZ\) and \(\widetilde{\bsigma}\in\bZ^c.\)
Then, adding \(\|\bsigma\|_{\bX}\) to both sides of \eqref{eq:case.2} shows \begin{align}\label{eq:case2.bound-x}
\|\bsigma\|_{\bX}\geq M\tnorm{(\bsigma,u)}_{\bX\times Y},\quad \text{where } M={\bigl({1+\frac{2\|a\|}{C_b}}\bigr)^{-1}}.\end{align} 
Multiplying by \(C_c\) to both sides of \eqref{eq:case2.bound-x} and applying the Braess condition, yields
%we derive the following estimate
\begin{align*}
    MC_c\tnorm{(\bsigma,u)}_{\bX\times Y}&\leq C_c\|\bsigma\|_{\bX}
    \leq\sup_{(\bzero,0)\neq(\btau,v)\in{\bX\times Y}}\frac{a(\bsigma,\btau)+b(\bsigma,v)}{\tnorm{(\btau,v)}_{\bX\times Y}}
    \\&\leq \sup_{(\bzero,0)\neq(\btau,v)\in{\bX\times Y}}\frac{\mathcal{A}(\bsigma,u;\btau,0)}{\tnorm{(\btau,v)}_{\bX\times Y}}
    \leq \sup_{(\bzero,0)\neq(\btau,v)\in{\bX\times Y}}\frac{\mathcal{A}(\bsigma,u;\btau,v)}{\tnorm{(\btau,v)}_{\bX\times Y}}.
\end{align*}
Combining the two cases, we arrive at the following stability bound 
\[\tnorm{(\bsigma,u)}_{\bX\times Y}\leq C  \sup_{(\bzero,0)\neq(\btau,v)\in{\bX\times Y}}\frac{\mathcal{A}(\bsigma,u;\btau,v)}{\tnorm{(\btau,v)}_{\bX\times Y}},\] where \begin{align}\label{eq:C}C:=\max\bigl[\frac{1}{MC_c},6\beta_1^{-1}+9C_b^2\beta_1^{-2}\bigr].\end{align}

{We now turn to showing \(\mathrm{BNB}~2\). We consider a fixed \((\bsigma,u)\in \bX\times Y\) for which \(\mathcal{A}(\bsigma,u;\btau,v)= 0\) {for all $(\btau,v) \in \bX\times Y$}, i.e., 
\begin{subequations}\label{eq:abstract.0}
    \begin{align}
    a(\bsigma,\btau) + b(\btau,u) & = 0 \quad \forall \btau \in \bX, \label{eq:weak.0.1}\\
    b(\bsigma, v) - t^2c(u,v) & = 0\label{eq:weak.0.2} \quad \forall v\in {Y_c}.
\end{align}\end{subequations}
Our goal is to show that \((\bsigma,u)=(\bzero,0).\) 
Substituting \(v=-u\) in \eqref{eq:weak.0.2} and then using the positive semi-definiteness of \(c(\bullet,\bullet),\) we conclude that \(b(\bsigma,u)=0\) (note that, the inf-sup condition of \(b(\bullet,\bullet)\) guarantees that, \(b(\bsigma,u)\geq 0\) for all \((\bsigma, u) \in \bX \times Y\)), which in particular implies that \(\bsigma\in\ker(\bB)\).
Taking \(\btau\in\ker(\bB)\) in \eqref{eq:weak.0.1}, and using \eqref{eq:inf-sup.a}, we infer that \(a(\bsigma,\btau)=0,\,\forall\tau\in\bZ,\) which further implies that \(\bsigma=\bzero\) (using \eqref{eq:inf-sup.a}). Finally, \eqref{eq:weak.0.1} yields \(b(\btau,u)=0,\,\forall \btau\in \bX,\)  and using \eqref{eq:inf-sup.b}, we conclude that \(u=0\).} This completes the rest of the proof.
\end{proof}
\begin{remark}
Unfortunately, we cannot readily derive   \(t\)-robust stability   such as  Theorem~\ref{th:abstract} by directly applying \cite[Theorem~3.4]{correa2022continuous} with the \(t\)-dependent norm \(\tnorm{(\bullet,\bullet)}_{\bX\times Y}:=\|\bullet\|_{\bX}+\|\bullet\|_Y+t|\bullet|_{c}\). This is  because the bounds \cite[Eq.~(3.56)-(3.58)]{correa2022continuous} depend explicitly on the scaling parameter \(t\) as follows
   \begin{align*}
       \tnorm{(\bsigma,u)}_{\bX\times Y}\leq& \bigl[\frac{1}{C_b}+(1+\frac{\|a\|}{C_b})\bigl(\frac{1}{C_a}+\frac{(2\max\bigl(C^{\sharp}_1,C^{\sharp}_2\bigr)^{\frac12})}{t}\bigr)+\\&\bigl[\frac{1}{C_b}+t(1+\frac{\|a\|}{C_b})\bigl(\frac{1}{C_a}+\frac{(2\max\bigl(C^{\sharp}_1,C^{\sharp}_2\bigr)^{\frac12})}{t}\bigr)\bigr] \|F\|_{\bX'}+ \bigl[(\frac{\|a\|}{C_b}+(1+\frac{\|a\|}{C_b})\frac{1}{C_b}+\\&\bigl(\frac{1}{C_a}+\frac{(2\max\bigl(C^{\sharp}_1,C^{\sharp}_2\bigr)^{\frac12})}{t}\bigr)+\bigl[\frac{1}{C_b}+t(1+\frac{\|a\|}{C_b})\bigl(\frac{1}{C_a}+\frac{(2\max\bigl(C^{\sharp}_1,C^{\sharp}_2\bigr)^{\frac12})}{t}\bigr)\bigr]\|G\|_{Y'},
   \end{align*}
   where \begin{align*}C^{\sharp}_1&:=(1+\frac{\|a\|}{C_a})\frac{1}{C_a}\frac{1}{C_b}+t^2\frac{\|c\|}{C_b^2}\bigl(1+\frac{\|a\|}{C_a}\bigr)^2,\qquad \text{and }\\
   C^{\sharp}_2&:=\frac{1}{C_b}(1+\frac{\|a\|}{C_a})\bigl[1+\frac{\|a\|}{C_b}+t^2\frac{\|a\|^2\|c\|}{C_b^3}\bigl(1+\frac{\|a\|}{C_a}\bigr)\bigr],
   \end{align*}
   and the stability  {deteriorates} with $t\to 0$. In addition, note that the proof of Theorem~\ref{th:abstract} differs from \cite[Lemma~3.1]{braess1996stability} as the Braess condition \eqref{eq:Braess} is different from \cite[(3.1)]{braess1996stability}.
   % Also, \(\mathrm{BNB~2}\) \eqref{eq:BNB2} needs to be verified in the Banach case, whereas it is not required  in the Hilbert space context. 
\end{remark}

%%%%%%%%%%%%%%%%%%%%%%%%%%%%%%%%%%%%%%%%%%%%%%%%%%%
\section{Existence and uniqueness of weak solution}\label{sec:existence}
This section  employs a fixed-point framework to split the Darcy--Forchheimer equations from the ADR system. After showing that the two uncoupled systems are well-posed, we  construct a fixed-point operator and then prove that it maps a ball into itself, and it satisfies the contraction mapping theorem. The treatment is similar to the analysis in, e.g.,  \cite{caucao2021fully,correa2023new,gatica2022p}.  


\subsection{Unique solvability of the decoupled sub-problems}
Now, for a fixed concentration $\hat{\varphi}\in \Psi$, consider the Darcy--Forchheimer sub-system in weak form as follows
\begin{subequations}\label{eq:weak.fixed.u.hat}
    \begin{align}
    a_1(\bu,\bv) + d_1(\bu; \bu,\bv) +  b_1(\bv,p) & = F^{\hat\varphi}(\bv) \quad \forall \bv \in \bV,\label{eq:weak,1.hat.u}\\
    b_1(\bu,q) & = 0 \quad \forall q \in \rQ.\label{eq:weak,2.hat.u}
    \end{align}
\end{subequations}
We recall from \cite[Theorem~3.3]{dhpr_sinum26} the following result for the well-posedness of the system \eqref{eq:weak,1.hat.u}-\eqref{eq:weak,2.hat.u}.

\begin{theorem}[Well-posedness of  Darcy--Forchheimer equations]\label{th:continuous-DF}
Assume that   $\bf(\hat{\varphi})\in \bL^2(\Omega)$ { for a fixed  $\hat{\varphi}\in \Psi.$}  Then there exists a unique  solution $(\bu,p)\in \bV\times\rQ$ to the problem \eqref{eq:weak,1.hat.u}-\eqref{eq:weak,2.hat.u}. Furthermore, there is a positive constant $C^{\mathsf{DF}}_{c},$ independent of the model parameters $\kappa, \rF$, such that 
\begin{equation*}
 \|(\bu, p )\|_{\bV\times\rQ}  \leq C^{\mathsf{DF}}_{c}\max \bigl\{\|\bf(\hat{\varphi})\|_{0,\Omega} ,\|\bf(\hat{\varphi})\|_{0,\Omega}^2\bigr\}.
    \end{equation*}
\end{theorem}
In particular, there is a positive constant ${C^{\mathsf{DF}}_{c,\bu}}$ (see, \cite[Equation~3.10]{dhpr_sinum26}), independent of   $\bu, \kappa,$ and $ \rF$, such that 
 \begin{align}\label{eq:u.estimate}
    \|\bu\|_{\bV}\leq {C^{\mathsf{DF}}_{c,\bu}}\,\|\bf(\hat{\varphi})\|_{0,\Omega} .\end{align}
    In what follows, we state and prove the unique-solvability result for the ADR equations. Prior to this, we establish the following {lemma}, which  later assists in proving the parameter-robust inf-sup stability. 
{\begin{lemma}\label{lem:inf-sup:ARD}
   There exists an operator \(\cS_c\) satisfying the following properties:
    \begin{gather}\nonumber
    \cS_c \in \cL(\Psi', \bW),\quad 
    \|\cS_c\|_{\cL(\Psi', \bW)}\ \text{is independent of } \alpha,{\varrho_0},\quad \text{and}\\
    {{\langle{\vdiv}_{\frac65} [\cS_c(\ell)], \psi\rangle}} = \langle \ell, \psi \rangle \quad \text{for all}\ \ell \in \Psi',\psi \in \Psi.\label{S:prop}
\end{gather}

\end{lemma}
\begin{proof}
We   define $\cS_c$ as a bounded linear operator from two different spaces, i.e. \(\cS_c\in\cL\bigl(\alpha^{-\frac12}(\rH^1_{\Gamma_{\bu}}(\Omega))',\alpha^{-\frac12}\bL^2(\Omega)\bigr),\) and \(\cS_c\in\cL\bigl((\rL^6(\Omega))',\bH_{\Gamma_p}(\vdiv_{\frac65},\Omega)\bigr),\) and then call upon the duality relation between sum and intersection between Banach spaces, i.e. \(\Psi'=[\alpha^{\frac12}\rH^1_{\Gamma_{\bu}}(\Omega)+\rL^6(\Omega)]'=\alpha^{-\frac12}(\rH^1_{\Gamma_{\bu}}(\Omega))'\cap(\rL^6(\Omega))'\) (see, for instance,  \eqref{eq:duality}, with \(X=\alpha^{\frac12}\rH^1_{\Gamma_{\bu}}(\Omega), Y=\rL^6(\Omega),\) and \(Z\) being the null space), to conclude that \( \cS_c \in {\cL}(\Psi', \bW) \).

The existence of \(\cS_c\in\cL\bigl((\rH^1_{\Gamma_{\bu}}(\Omega))',\bL^2(\Omega)\bigr)\)
follows directly from \cite[Lemma~3.2]{dhpr_sinum26}. Moreover, a scaling by \(\alpha^{-\frac12}\) does not compromise the \(\alpha\)-robustness of the boundedness constant of \(\cS_c.\) 

For the existence of \(\cS_c\in\cL\bigl((\rL^6(\Omega))',\bH_{\Gamma_p}(\vdiv_{\frac65},\Omega)\bigr)\) with \(\ell\in\bigl(\rL^6(\Omega)\bigr)',\) let \(\vartheta \in\rH^{1}_{\Gamma_p}(\Omega)\) be the unique solution of the following elliptic problem with mixed boundary conditions {and homogeneous data} (owing to, e.g., \cite[Chapter~7]{maz_i_a2010elliptic})
\begin{align*}
    \Delta \vartheta=\ell \quad \text{in }\Omega, \quad \nabla \vartheta \cdot \bn = 0\quad \text{in }\Gamma_{\bu}, \quad \vartheta=0 \quad \text{in } \Gamma_p.
\end{align*}
Next, setting \(\tilde{\boldsymbol{\vartheta}} = -\nabla \vartheta\in\rL^{2}(\Omega),\) it is evident that \[\vdiv(\tilde{\boldsymbol{\vartheta}})=\Delta \vartheta=\ell\in\rL^{\frac65}(\Omega), \quad \tilde{\boldsymbol{\vartheta}}\cdot \bn = 0, \quad \text{and}\quad\|\tilde{\boldsymbol{\vartheta}}\|_{\vdiv_{\frac65},\Omega}\leq \|\ell\|_{0,\frac65,\Omega}.\]
Then it suffices to define \(\cS_c(\ell) =\tilde{\boldsymbol{\vartheta}},\) which further implies that \(\cS_c\in\cL\bigl((\rL^6(\Omega))',\bH_{\Gamma_p}(\vdiv_{\frac65},\Omega)\bigr)\). 
\end{proof}
% \rd{\begin{lemma}\label{est:L3-L2+Hdiv6/5}
%    For $\btau\in \bL^2(\Omega)\cap\bH_{\Gamma_p}(\vdiv_{\frac65},\Omega),$ the following inequality holds
%     \begin{align*}
%         \|\btau\|_{0,3,\Omega}\lesssim (\|\btau\|_{0,\Omega}+\|\vdiv\btau\|_{0,\frac65,\Omega}),
%     \end{align*}
%     where the absorbed constant is dependent on the Sobolev embedding constants, and the Poincar\'e inequality.
% \end{lemma}
% \begin{proof}
%     Let us start by recalling that $\bigl(\bL^3(\Omega)\bigr)^*=\bL^{\frac32}(\Omega),$ and $\|\bullet\|_{0,3,\Omega}$ is given by the following identity \begin{align}\label{def:L3-L3/2.duality.norm}
% \|\btau\|_{0,3,\Omega}=\sup_{\bz\in\bL^{\frac32}(\Omega),\|\bz\|_{0,\frac32,\Omega}=1}\int_{\Omega}\btau\cdot\bz.\end{align}
% Now for $\bz\in\bL^{\frac32}(\Omega),$ suppose  $w\in\rW_{\Gamma_p}^{1,6}(\Omega)$ be the unique solution of the following pde
% \begin{align*}
%     -\Delta w&=\vdiv\bz\quad\text{in }\Omega,\\
%     w=0 \quad\text{on } \Gamma_p, &\,\nabla w\cdot\bn=0\quad\text{on }
% \Gamma_{\bu}.\end{align*}
% Then from elliptic regularity, it follows that $w\in \rW_{\Gamma_p}^{1,6}(\Omega),$ and  $\|w\|_{1,6,\Omega}\leq C_{\mathrm{ell.reg}}\|\bz\|_{0,\frac32,\Omega},$ and in particular, 
% $$\|w\|_{0,6,\Omega}+\|\nabla w\|_{0,6,\Omega}\leq C\|\bz\|_{0,\frac32,\Omega}.$$
% Also, defining $\bz_0=\bz-\nabla w,$ it is not difficult to observe that $\vdiv \bz_0=0.$ Now using the  integration-by-parts formula, we have $$\int_{\Omega}\btau\cdot\bz=\int_{\Omega}\btau\cdot\nabla w + \int_{\Omega}\btau\cdot\bz_0.$$
% Using H\"older's inequality, and elliptic estimate, the first term can be bounded by the following 
% $$\int_{\Omega}\btau\cdot\nabla w\leq \|\vdiv\btau\|_{0,\frac65,\Omega}\|w\|_{0,6,\Omega}\leq C\|\vdiv\btau\|_{0,\frac65,\Omega}\|\bz\|_{0,\frac32,\Omega}.$$ For the second term, we get the following bound : $$\int_{\Omega}\btau\cdot\bz_0\leq\|\btau\|_{0,\Omega}\|\bz_0\|_{0,\Omega}\leq C\cdots$$
% \end{proof}
% }
%\noindent

Next, for a fixed advecting velocity \(\hat{\bu}\in\bV\), let us consider the following ADR system in its weak form
\begin{subequations}\label{eq:weak.hat.phi}
    \begin{align}
    a_2(\bzeta,\bw) + b_2(\bw,\varphi) + d_2^{\hat\bu}(\bw, \varphi) & = 0 \quad \forall \bw \in \bW, \label{eq:weak,3.hat.phi}\\
    b_2(\bzeta, \psi) - c_2(\varphi,\psi) & = H(\psi) \quad \forall \psi \in \Psi, \label{eq:weak,4.hat.phi}
    \end{align}
\end{subequations}
and we state the following continuous stability result.
\begin{theorem}[Well-posedness of decoupled advection-diffusion-reaction equations]\label{th:continuous-ARD}
 Assume that   
 $m \in \rL^{\frac65}(\Omega)$. Then, for a fixed $\hat{\bu}\in \bV,$ there {exist} positive constants \(C_c, C^c_{\alpha}\) (depending on $\alpha$), such that, whenever {\(\|\hat\bu\|_{{0,3,\Omega}}\leq  \frac{C_c}{2C^c_{\alpha}},\)} there exists a unique  solution $(\bzeta,\varphi)\in \bW\times\Psi$ to the problem \eqref{eq:weak,3.hat.phi}-\eqref{eq:weak,4.hat.phi}.  Furthermore, there is a positive constant $C^{\mathsf{ADR}}_{c}$, independent of $\varrho_0,$ such that 
\begin{equation}\label{eq:bound.phi.psi}
 \tnorm{(\bzeta, \varphi)}_{\bW\times\Psi}  \leq C^{\mathsf{ADR}}_{c} \| m \|_{0,\frac65,\Omega} .
    \end{equation}
\end{theorem}


\begin{proof} 
We first use Theorem~\ref{th:abstract} to show the well-posedness of the following reaction-diffusion system, which is identical to \eqref{eq:weak.hat.phi} without the advection part, i.e., one seeks \((\bzeta,\varphi)\in\bW\times \Psi\) such that
\begin{subequations}\label{eq:weak.form.reac-diff}
\begin{align}
    a_2(\bzeta,\bw) + b_2(\bw,\varphi) & = 0 \quad \forall \bw \in \bW, \label{eq:weak,3_1}\\
    b_2(\bzeta, \psi) - c_2(\varphi,\psi) & = H(\psi) \quad \forall \psi \in \Psi. \label{eq:weak,4_1}
\end{align}
\end{subequations}


\smallskip
\noindent\textbf{Boundedness of \(a_2(\bullet,\bullet),\,b_2(\bullet,\bullet)\) and \(c_2(\bullet,\bullet)\).} Observe that a use of  the {Cauchy--Schwarz} inequality gives the
 boundedness of the bilinear forms \(a_2(\bullet,\bullet)\) and \(c_2(\bullet,\bullet)\) as
 \begin{align*}
a_2(\bzeta,\btau)&=\int_{\Omega}\alpha^{-1}\bzeta\cdot\btau\leq\alpha^{-\frac12}\|\bzeta\|_{0,\Omega}\,\alpha^{-\frac12}\|\btau\|_{0,\Omega} \leq\|\bzeta\|_{\bW}\|\btau\|_{\bW}\quad \forall \bzeta,\btau\in \bW,\\
c_2(\varphi,\psi)&=\int_{\Omega}{\varrho_0}\varphi\psi\leq\|\varphi\|_{0,\Omega}\|\psi\|_{0,\Omega}\leq \|\varphi\|_{\Psi_c}\|\psi\|_{\Psi_c}\qquad \forall \varphi, \psi \in \Psi_c,
 % \leq \rrb{\|\varrho_0 - g \|^{1/2}_{0,\Omega}\|\varphi\|_{0,4,\Omega}\,\|\varrho_0 - g \|^{1/2}_{0,\Omega}\|\psi\|_{0,4,\Omega} } \\
 % & \lesssim \rrb{\|\varphi\|_{0,6,\Omega}\|\psi\|_{0,6,\Omega} }\qquad \forall \varphi, \psi \in \Psi,
%\varrho_0^{\frac12}\|\varphi\|_{0,\Omega}\,\varrho_0^{\frac12}\|\psi\|_{0,\Omega}\leq\|\varphi\|_{c}\|\psi\|_{c} 
\end{align*}
In order to show the boundedness of \(b_2(\bullet,\bullet)\), we follow a decomposition argument over \(\Psi.\) Let \(\psi=\psi_0+\psi_1,\) with \(\psi_0\in\rL^6(\Omega)\) and \(\psi_1\in\rH^1_{\Gamma_{\bu}}(\Omega).\) Then, 
\(b_2(\btau,\psi)=b_2(\btau,\psi_0+\psi_1)=\langle\vdiv\btau,\psi_0\rangle-(\alpha^{-\frac12}\btau,\alpha^{\frac12}\nabla\psi_1).\)
    Next, a use of the H\"{o}lder's  with  the {Cauchy--Schwarz} inequality, {the fact that $a^2+b^2\leq(a+b)^2,$\, for any two  {nonnegative} real numbers $a,b,$} and then taking the infimum over all decompositions of \(\psi\in\Psi\) give the desired boundedness as follows
    \begin{align*}
   b_2&(\btau,\psi) \leq \|\vdiv\btau\|_{0,\frac65,\Omega}\|\psi_0\|_{0,6,\Omega}+\|\alpha^{-\frac12}\btau\|_{0,\Omega}\|\alpha^{\frac12}\nabla\psi_1\|_{0,\Omega}\\
    &\leq\bigl(\|\vdiv\btau\|^2_{0,\frac65,\Omega}+\|\alpha^{-\frac12}\btau\|^2_{0,\Omega}\bigr)^{\frac12}\bigl(\|\psi_0\|^2_{0,6,\Omega}+\|\alpha^{\frac12}\nabla\psi_1\|^2_{0,\Omega}\bigr)^{\frac12} \\
    &\leq\bigl(\|\vdiv\btau\|_{0,\frac65,\Omega}+\|\alpha^{-\frac12}\btau\|_{0,\Omega}\bigr)\bigl(\|\psi_0\|_{0,6,\Omega}+\|\alpha^{\frac12}\nabla\psi_1\|_{0,\Omega}\bigr)=\|\btau\|_{\bW}\|\psi\|_{\Psi}\quad {\forall \btau \in \bW, \psi \in \Psi.}
\end{align*}

\smallskip
\noindent\textbf{Symmetry and positive-semidefiniteness of \(a_2(\bullet,\bullet)\) and \(c_2(\bullet,\bullet)\).} Clearly, \(a_2(\bullet,\bullet)\) and \(c_2(\bullet,\bullet)\) are symmetric, as well as positive-semidefinte bilinear forms, as  \[a_2(\bzeta,\bzeta)={\alpha^{-1}\int_\Omega \bzeta\cdot\bzeta} \geq0 \quad \text{and}\quad  c_2(\varphi,\varphi)= \varrho_0\int_{\Omega}\varphi^2\geq0 \qquad \forall \bzeta\in\bW, \varphi\in\Psi_c.\] 
%which together imply the positive semidefiniteness of 
%\(a(\bullet,\bullet)\) and \(c(\bullet,\bullet).\) 

\smallskip
\noindent\textbf{Inf-sup of \(a_2(\bullet,\bullet)\) on $\bZ$.} Choosing $\btau = \bzeta$ and using the kernel characterisation with definition of the $\bW$-norm, we arrive at  
\[ \sup_{\btau \in \bZ\setminus\{\boldsymbol{0}\}} \frac{a_2(\bzeta,\btau)}{\|\btau\|_{\bW}} {{\geq \alpha^{-1}\frac{\int_\Omega \bzeta\cdot\bzeta}{\|\bzeta\|_{\bW}}} =}  C_{a_2}\| \bzeta\|_{\bW} \quad \forall \bzeta \in \bZ.\]
Then, this readily 
satisfies the second hypothesis of Theorem~\ref{th:abstract} with the constant~\(C_{a_2}=1\).

\smallskip
\noindent\textbf{Inf-sup stability of \(b_2(\bullet,\bullet)\).} 
From the existence of $\cS_c$ defined in Lemma~\ref{lem:inf-sup:ARD}, it immediately follows that
\begin{align*}
\sup_{\bzero\neq\btau \in \bW} \!\!\frac{\int_\Omega \psi\,\vdiv \btau}{\|\btau\|_\bW} \geq \sup_{0\neq\ell \in \Psi'} \frac{{\langle{\vdiv}_{\frac65} [\cS_c(\ell)], \psi\rangle}}{\|\cS_c(\ell)\|_\bW} = \sup_{0\neq\ell \in \Psi'} \frac{\langle \ell, \psi \rangle}{\|\cS_c(\ell)\|_\bW} \geq \|\cS_c\|^{-1}_{\cL(\Psi', \bW)} \sup_{0\neq\ell \in \Psi'} \frac{\langle \ell, \psi \rangle}{\|\ell\|_{\Psi'}} = C_{b_2} \|\psi\|_{\Psi},
\end{align*}
where $C_{b_2} = \|\cS_c\|^{-1}_{\cL(\Psi', \bW)}$ is independent of \(\alpha,\varrho_0\), and hence, the inf-sup condition is satisfied.

\smallskip
\noindent\textbf{Braess condition.} Now, fixing  \(\btau=\bzeta,\psi=0\), and observing that \(\tnorm{(\bzeta,0)}_{\bW\times \Psi}=\|\bzeta\|_{\bW}\), we arrive at  
\begin{align}\label{eq:Braess.1st}
\sup_{(\bzero,0)\neq(\btau,\psi)\in\bW\times\Psi}\frac{a_2(\bzeta,\btau)+b_2(\bzeta,\psi)}{\tnorm{(\btau,\psi)}_{\bW\times\Psi}} \geq \alpha^{-1}\frac{\int_\Omega \bzeta\cdot \bzeta}{\tnorm{(\bzeta,0)}_{\bW\times\Psi}}=\alpha^{-1}\frac{\|\bzeta\|^2_{0,\Omega}}{\|\bzeta\|_{\bW}}.
\end{align}
 Next, taking \(\btau = \bzero,\) and \(\psi=\mathrm{sgn}(\vdiv\bzeta)|\vdiv\bzeta|^{\frac15}{\in\rL^{\frac65}(\Omega)}\), 
we evaluate the following quantities 
\begin{align*}
   b_2(\boldsymbol{\zeta}, \psi)
  &= \int_{\Omega} |\operatorname{div} \boldsymbol{\zeta}|^{\frac{6}{5}}
  = \|\operatorname{div} \boldsymbol{\zeta}\|_{0,\frac{6}{5},\Omega}^{\frac{6}{5}}, 
  \quad \quad \|\psi\|^6_{0,\Omega}
  = \int_\Omega|\vdiv\bzeta|^{\frac65}=\|\operatorname{div} \boldsymbol{\zeta}\|^{\frac65}_{0,\frac65,\Omega},
  \\
    \tnorm{(\btau, \psi)}_{\boldsymbol{W} \times \Psi}
  &= \tnorm{(\boldsymbol{0}, \psi)}_{\boldsymbol{W} \times \Psi}
  {\leq} \|\psi\|_{\Psi} + {\varrho_0}\|\psi\|_{0,6,\Omega}
  \leq {2}\|\psi\|_{0,6,\Omega},  
\end{align*}
where {the last inequality follows from \eqref{eq:Psi.norm}, and by using $\varrho_0\leq 1$}. 
 Furthermore using {the {above} expressions} for \(b_2(\boldsymbol{\zeta}, \psi),\,\|\psi\|_{0,\Omega}\) and \(\tnorm{(\btau, \psi)}_{\boldsymbol{W} \times \Psi}\), we immediately obtain 
 %the following inequality 
\begin{align}\label{eq:Braess.2nd} \sup_{(\bzero,0)\neq(\btau,\psi)\in\bW\times\Psi}\frac{a_2(\bzeta,\btau)+b_2(\bzeta,\psi)}{\tnorm{(\btau,\psi)}_{\bW\times\Psi}} \geq{\frac12} \frac{\|\vdiv \bzeta\|^{\frac65}_{0,\frac65,\Omega}}{\|\vdiv \bzeta\|^{\frac15}_{0,\frac65,\Omega}}={\frac12}\|\vdiv \bzeta\|_{0,\frac65,\Omega}.
\end{align}
Now, combining \eqref{eq:Braess.1st} and \eqref{eq:Braess.2nd}, we arrive at 
\begin{align}\label{eq:Braess.3rd}
\sup_{(\bzero,0)\neq(\btau,\psi)\in\bW\times\Psi}\frac{a_2(\bzeta,\btau)+b_2(\bzeta,\psi)}{\tnorm{(\btau,\psi)}_{\bW\times\Psi}}\geq\max\bigl(\alpha^{-1}{\frac{\|\bzeta\|^2_{0,\Omega}}{\|\bzeta\|_{\bW}},{\frac12}\|\vdiv \bzeta\|_{0,\frac65,\Omega}}\bigr).
\end{align}
Set \(s:=\alpha^{-\frac12}\frac{\|\bzeta\|_{0,\Omega}}{\|\bzeta\|_{\bW}}\in[0,1)\) and use \eqref{eq:Braess.3rd} to derive the following  relation 
%\rrb{is this still the same with $g=0$?} [\rd{yes!}]
\begin{align*}\max\bigl({\alpha^{-1}\frac{\|\bzeta\|^2_{0,\Omega}}{\|\bzeta\|_{\bW}},\|\vdiv \bzeta\|_{0,\frac65,\Omega}}\bigr)=\bigl(\max_{s\in[0,1)}\bigl[s^2,{\frac12}\sqrt{1-s^2}\bigr]\bigr) \|\bzeta\|_{\bW}.\end{align*}
After denoting \[\gamma_c:=\bigl(\max_{s\in[0,1)}\bigl[s^2,{\frac12}\sqrt{1-s^2}\bigr]\bigr)<1,\] we can conclude from \eqref{eq:Braess.3rd} that \[\sup_{(\bzero,0)\neq(\btau,\psi)\in\bW\times\Psi}\frac{a_2(\bzeta,\btau)+b_2(\bzeta,\psi)}{\tnorm{(\btau,\psi)}_{\bW\times\Psi}}\geq \gamma_c \|\bzeta\|_{\bW}.\]
An appeal to Theorem~\ref{th:abstract} now yields the well-posedness of \eqref{eq:weak.form.reac-diff}, and additionally, we obtain the following inf-sup type estimate for $\cA$  
\begin{align}\label{eq:A.stab}
\sup_{(\bzero,0)\neq(\tau,\psi)\in\bW\times\Psi}\frac{\cA(\bzeta,\varphi;\btau,\psi)}{\tnorm{(\btau,\psi)}_{\bW\times \Psi}}\geq C_c\tnorm{(\bzeta,\varphi)}_{\bW\times \Psi},
\end{align}
where \(C_c\) depends on \(\|a_2\|,C_{b_2},\gamma_c\) {(see \eqref{eq:C} for the explicit expression of \(C_c\))}.

Next, in order to show the well-posedness of the uncoupled ADR equations, for a fixed \(\hat{\bu}\in\bV,\) we consider another bilinear form \(\mathcal{A}^{\mathsf{ADR}}:\bW\times\Psi\to \mathbb{R}\) defined by 
\begin{align}\label{eq:A:ARD}\mathcal{A}^{\mathsf{ADR}}(\bzeta,\varphi;\btau,\psi):=\mathcal{A}(\bzeta,\varphi;\btau,\psi)+d_2^{\hat{\bu}}(\btau,\varphi),\quad \forall (\bzeta,\varphi), (\btau,\psi)\in\bW\times \Psi.\end{align}
Thus,  showing the well-posedness of the system~\eqref{eq:weak,3}-\eqref{eq:weak,4} boils down to analysing the problem of finding \((\bzeta,\varphi)\in\bW\times \Psi\) satisfying \[\mathcal{A}^{\mathsf{ADR}}(\bzeta,\varphi;\btau,\psi)=H(\psi),\quad \forall \psi\in \Psi.\] 


Note that, by the generalised Hölder's inequality with {1/3+1/2+1/6=1}, and the uniform positivity and boundedness of $\kappa,$ the following estimate readily follows:\begin{align}\label{eq:d2.bound}
d_2^{\hat{\bu}}(\btau,\varphi)\leq C^c_{\alpha}\|\hat{\bu}\|_{0,3,\Omega}\,\tnorm{(\btau,\psi)}_{\bW\times\Psi}\,\tnorm{(\bzeta,\varphi)}_{\bW\times\Psi},\end{align} where {$C^c_{\alpha}:={{\alpha^{-\frac12}}}(1+C_{\mathsf{emb}}\alpha^{-\frac12}).$}  
Invoking the bound of \(d_2^{\hat{\bu}}(\bullet,\bullet)\) (cf. \eqref{eq:d2.bound}), it follows from \eqref{eq:A.stab} and \eqref{eq:A:ARD} that
\[\sup_{(\bzero,0)\neq(\btau,\psi)\in {\bW\times\Psi}}\frac{\mathcal{A}^{\mathsf{ADR}}(\bzeta,\varphi;\btau,\psi)}{\tnorm{(\btau,\psi)}_{\bW\times\Psi}}\geq\bigl[C_c-C^c_{\alpha}{\|\hat{\bu}\|_{0,3,\Omega}}\bigr]\tnorm{(\bzeta,\varphi)}_ {\bW\times\Psi}, \qquad  \forall(\bzeta,\varphi)\in \bW\times \Psi.\]
Hence, under the assumption that { for each $\hat{\bu} \in \bV$ such that  {$\|\hat{\bu}\|_{{0,3,\Omega}}\leq \frac{C_c}{2C^c_{\alpha}}$}, we arrive at 
\begin{equation}\label{eq:ARD}
\sup_{(\bzero,0)\neq(\btau,\psi)\in {\bW\times\Psi}}\frac{\mathcal{A}^{\mathsf{ADR}}(\bzeta,\varphi;\btau,\psi)}{\tnorm{(\btau,\psi)}_ {\bW\times\Psi}}\geq\frac{C_c}{2} \tnorm{(\bzeta,\varphi)}_ {\bW\times\Psi}.
\end{equation}
Therefore,} the ADR system (cf. \eqref{eq:weak,3}-\eqref{eq:weak,4}) is well-posed. 
Additionally, using \eqref{eq:A:ARD} in \eqref{eq:ARD}, and using the H\"{o}lder's inequality to see that \(H(\psi)\leq\|m\|_{0,\frac65,\Omega}\|\psi\|_{0,6,\Omega},\) we obtain the continuous dependence on the data as follows: \begin{align}\label{eq:C.c.ard}
\tnorm{(\bzeta,\varphi)}_{\bW\times \Psi}\leq C^{\mathsf{ADR}}_{c}\|m\|_{0,\frac65,\Omega},
\end{align} with {$C^{\mathsf{ADR}}_{c}:={2}{C_c}^{-1}(1+C_{\mathsf{emb}}\alpha^{-\frac12})$} being independent of~\(\varrho_0\). This concludes the rest of the proof.
\end{proof}

\subsection{Fixed-point analysis}
Now, define the following  map
\[ \cS^{\mathsf{DF}}:\Psi\to \bV \times \rQ, \quad \hat{\varphi} \mapsto \cS^{\mathsf{DF}}(\hat{\varphi})= (\cS^{\mathsf{DF}}_1(\hat{\varphi}), \cS^{\mathsf{DF}}_2(\hat{\varphi})):= (\bu, p), \]
where $( \bu,p) \in \bV \times \rQ$ is the unique solution of the Darcy--Forchheimer equations {for fixed $\hat\varphi\in \ \Psi$}, as stated in Theorem~\ref{th:continuous-DF}, {and $\cS_1^{\mathsf{DF}}:\Psi\to\bV$,  $\cS_2^{\mathsf{DF}}:\Psi\to\rQ.$}
Further, {for a fixed $\hat\bu\in\bV,$} set the solution operator associated with the mixed ADR equations as
\[ \cS^{\mathsf{ADR}}: \bV \to  \bW\times \Psi, \quad \hat{\bu} \mapsto \cS^{\mathsf{ADR}}(\hat{\bu})=(\cS^{\mathsf{ADR}}_1(\hat{\bu}),\cS^{\mathsf{ADR}}_2(\hat{\bu})):= ( \bzeta, \varphi),\]
where $(\bzeta, \varphi)$ is the unique solution of the ADR equations, as stated in Theorem~\ref{th:continuous-ARD}, {and $\cS^{\mathsf{ADR}}_1:\bV\to\bW, \cS^{\mathsf{ADR}}_2:\bV\to\Psi$}. These maps are well-defined; the same holds for the following one
\begin{equation}\label{Operator-S.V->V}
\cS: \bV  \to  \bV, \quad \hat{\bu} \mapsto \cS(\hat{\bu}):= \cS^{\mathsf{DF}}_1(\cS^{\mathsf{ADR}}_2(\hat{\bu})).
\end{equation}
Finding a fixed point $\bu$ of $\cS$ is  therefore  equivalent to find a solution to   \eqref{eq:weak}. We use the Banach contraction mapping theorem, and consider, for a generic $r>0$, the following closed ball  
\[ \mathrm{B}_{\bzero}^r : = \{ \hat{\bu} \in \bV:
  \|\hat{\bu}\|_{\bV}  \leq r \},\]
and proceed next to show that $\cS$ maps it to itself and that $\cS$ is a contraction. For \(\bu\in\mathrm{B}_{\bzero}^r\), using  Theorem~\ref{th:continuous-DF}, \eqref{eq:u.estimate}, \eqref{eq:lip.bound.f}, and Theorem~\ref{th:continuous-ARD} shows that \begin{align}\label{eq:S.bound}
\|\cS(\bu)\|_{\bV}=\|\cS_1^{\mathsf{DF}}(\cS_2^{\mathsf{ADR}}(\bu))\|_{\bV}\leq {C^{\mathsf{DF}}_{c,\bu}}\bigl[\|\bf(\cS_2^{\mathsf{ADR}}(\bu))\|_{0,\Omega}\bigr]\leq {C^{\mathsf{DF}}_{c,\bu}}\bigl[{\|\bf(0)\|_{0,\Omega}}+C^{\mathsf{ADR}}_{c}\alpha_1\|m\|_{0,\frac65,\Omega}\bigr].
\end{align}
 
\begin{lemma}[Ball mapping property]\label{lemma:ball}
Under the small data assumption 
\begin{equation}\label{data-bounded-map-closed}
{C^{\mathsf{DF}}_{c,\bu}}\bigl[{\|\bf(0)\|_{0,\Omega}}+C^{\mathsf{ADR}}_{c}\alpha_1\|m\|_{0,\frac65,\Omega}\bigr] \leq r<\frac12,
\end{equation}
we have that $\cS\big(\mathrm{B}_{\bzero}^r\big)\subseteq\mathrm{B}_{\bzero}^r$.
\end{lemma}
\begin{proof}
  Given $\hat{\bu} \in \mathrm{B}_{\bzero}^r$, a use of  \eqref{Operator-S.V->V} with  \eqref{eq:S.bound} and \eqref{data-bounded-map-closed}  yields
     \(
         \|\cS(\hat{\bu})\|_{\bV}= \|  \cS^{\mathsf{DF}}_1(\cS^{\mathsf{ADR}}_2(\hat{\bu}))\|_{\bV}
         % \lesssim \|m  \|_{0,\frac65,\Omega}  
         \leq r<\frac12,
     \)  
     and hence, $\cS\big(\mathrm{B}_{\bzero}^r\big)\subseteq\mathrm{B}_{\bzero}^r,$ which concludes the proof.
\end{proof}
 \begin{lemma}[Lipschitz-continuity]\label{lemma:continuity-S} 
 There exists a positive constant $L_{\cS}$ such that 
 \begin{equation}\label{continuity-S}
  \|\cS(\bu_{1})-\cS(\bu_{2})\|_{\bV}
\leq L_{\cS}\|\bu_{1}-\bu_{2}\|_{\bV} \qquad \forall \bu_1,\bu_2 \in \bV.
 \end{equation} 
 \end{lemma}
\begin{proof}
It is enough to show the Lipschitz-continuity of the maps \(\cS_1^{\mathsf{DF}}\) and \(\cS_2^{\mathsf{ADR}}\), as in that case, using the Lipschitzness of the individual maps, one simply obtains the following 
\begin{align*}\|\cS(\bu_{1})-\cS(\bu_{2})\|_{\bV}&=\|\cS_1^{\mathsf{DF}}(\cS_2^{\mathsf{ADR}}(\bu_1))-\cS_1^{\mathsf{DF}}(\cS_2^{\mathsf{ADR}}(\bu_2))\|_{\bV}\\
&\leq L_{\cS_1}\|\cS_2^{\mathsf{ADR}}(\bu_1)-\cS_2^{\mathsf{ADR}}(\bu_2)\|_{\bV}\leq L_{\cS_1}L_{\cS_2}\|\bu_1-\bu_2\|_{\bV},
\end{align*}
with respective Lipschitz constants, \(L_{\cS_1}\) and \(L_{\cS_2}\) (to be specified below) associated {with} {the maps} \(\cS^{\mathsf{DF}}_1\) and \(\cS^{\mathsf{ADR}}_2\). We now proceed to show the Lipschitzness of \(\cS_1^{\mathsf{DF}}\) and \(\cS_2^{\mathsf{ADR}}\).

\smallskip
\noindent\textbf{Lipschitz-continuity of \(\cS^{\mathsf{DF}}_1\).}
Let \(\varphi_1, \varphi_2\in\Psi\) be two  elements in \(\Psi.\) Fixing \(\varphi_1\) and \(\varphi_2\) results in two  solutions \((\bu_1,p_1), (\bu_2,p_2) \in \bV\times \rQ\) corresponding to   \eqref{eq:weak,1}-\eqref{eq:weak,2}. Now, denote \(\cS^{\mathsf{DF}}_1(\varphi_1):=\bu_1\) and \(\cS^{\mathsf{DF}}_1(\varphi_2):=\bu_2.\) Substituting \(\bv=\bu_1-\bu_2,q=p_1-p_2\) into \eqref{eq:weak,1}, \eqref{eq:weak,2}, {using \eqref{eq:lip.bound.f}, and} an application of  the uniformly strong monotonicity of \(a_1(\bullet,\bullet)+d_1(\bullet;\bullet,\bullet)\) in the new equation obtained by subtracting \eqref{eq:weak,2} from \eqref{eq:weak,1} implies
\begin{align}\label{eq:Lipschitz.S_DF}
    \|\cS_1^{\mathsf{DF}}({\varphi_1})-\cS_1^{\mathsf{DF}}({\varphi_2})\|_{\bV}=\|\bu_1-\bu_2\|_{\bV}\leq L_{\cS_1}\|\varphi_1-\varphi_2\|_{\Psi},
\end{align}
where {\(L_{\cS_1}:=\alpha^{-\frac12}_{\mathsf{DF}}{\alpha_1}C^2_{\mathsf{emb}}\), \(\alpha_{\mathsf{DF}}\)} being the strong monotonicity constant (see, \cite[Equation~3.8]{dhpr_sinum26}), and \(\alpha_1\) is the local Lipschitz continuity constant for  the Darcy--Forchheimer system (see, \cite[Equation~3.7]{dhpr_sinum26}), which can be further bounded by \({\|\bf(0)\|_{0,\Omega}},\alpha_1, C^{\mathsf{ADR}}_{c}, \) and \(\|m\|_{0,\frac65,\Omega}\) (cf.  \eqref{eq:u.estimate} and \eqref{eq:lip.bound.f}).

\smallskip
\noindent\textbf{Lipschitz-continuity of \(\cS^{\mathsf{ADR}}_2\).}
{Let $\bu_1, \bu_2\in\bV$ be two elements in $\bV.$ Fixing $\bu_1, \bu_2$ in \eqref{eq:weak,3},\eqref{eq:weak,4} produces the solutions $(\bzeta_1,\varphi_1), (\bzeta_2,\varphi_2)\in\bW\times\Psi$ corresponding to  \eqref{eq:weak,3}-\eqref{eq:weak,4}. Let us denote $\varphi:=\varphi_1-\varphi_2.$}  Subtracting \eqref{eq:weak,4} from \eqref{eq:weak,3}, adding {$\int_{\Omega}\bu_1\cdot\bw\varphi_2, \int_{\Omega}\bu_2\cdot\bw\varphi_1$,} and subtracting {$2\int_{\Omega}\bu_2\cdot\bw\varphi_2$}, and using Theorem~\ref{th:continuous-ARD}  we obtain
\begin{align*}
   \frac1{C^{\mathsf{ADR}}_{c}} \tnorm{(\bzeta,\varphi)}_{\bW\times\Psi}&\leq\sup_{\substack{(\bw,\psi)\in\bW\times \Psi\\(\bw,\psi)\neq (\bzero,0)}}\frac{{-\int_{\Omega}((\bu_1-\bu_2)\cdot\bw)\varphi_2 + \int_{\Omega}(\bu_2\cdot\bw)\varphi}}{\tnorm{(\bw,\psi)}_{\bW\times \Psi}}.
    \end{align*}
    Finally, using H\"{o}lder's inequality and \eqref{eq:bound.phi.psi}, we conclude the following 
  {  \begin{align}\label{eq:Lip.ARD}
   \frac1{C^{\mathsf{ADR}}_{c}}\tnorm{(\bzeta,\varphi)}_{\bW\times\Psi}&\leq\sup_{\substack{(\bw,\psi)\in\bW\times \Psi\\(\bw,\psi)\neq (\bzero,0)}}\frac{{\rF^{-\frac13}}\|\bu_1-\bu_2\|_{\bV} \tnorm{(\bw,\psi)}_{\bW\times\Psi}\|m\|_{0,\frac65,\Omega}}{{\tnorm{(\bw,\psi)}_{\bW\times \Psi}}}\nonumber\\&\quad +\sup_{\substack{(\bw,\psi)\in\bW\times \Psi\\(\bw,\psi)\neq (\bzero,0)}}\frac{{\rF^{-\frac13}}\|\bu_2\|_{\bV}\tnorm{(\bw,\psi)}_{\bW\times\Psi}\tnorm{(\bzeta,\varphi)}_{\bW\times\Psi}}{{\tnorm{(\bw,\psi)}_{\bW\times \Psi}}},\nonumber
\end{align}
which further implies that 
\begin{align}
    \tnorm{(\bzeta,\varphi)}_{\bW\times\Psi} 
    \leq {\frac{2C^c_{\alpha}}{C_c(C^c_{\alpha}-\rF^{-\frac13})}}\|m\|_{0,\frac65,\Omega}\; \|\bu_1-\bu_2\|_{\bV},
\end{align}
where the above inequality makes sense when $C^c_{\alpha}\rF^{\frac13} \,>\, 1.$ 
This concludes the Lipschitz-continuity 
of \(\cS\) by invoking the bounds from \eqref{eq:Lipschitz.S_DF} and \eqref{eq:Lip.ARD} as 
{\begin{align}\label{eq:Lipschitz.S_ARD} \|\cS(\bu_1)-\cS(\bu_2)\|_{\bV}&=\|\cS_1^{\mathsf{DF}}(\cS_2^{\mathsf{ADR}}(\bu_1))-\cS_1^{\mathsf{DF}}(\cS_2^{\mathsf{ADR}}(\bu_2))\|_{\bV}\\&\leq\alpha_{\mathsf{DF}}^{-1}\alpha_1\|\cS_2^{\mathsf{ADR}}(\bu_1)-\cS_2^{\mathsf{ADR}}(\bu_2)\|_{\Psi}\nonumber\leq \frac{2C^c_{\alpha}}{C_c(C^c_{\alpha}-1)}\|m\|_{0,\frac65,\Omega}\alpha_{\mathsf{DF}}^{-\frac12}\alpha_1\|\bu_1-\bu_2\|_{\bV},
\end{align}}
with the Lipschitz constant {\(L_{\cS}:={\frac{2C^c_{\alpha}}{C_c(C^c_{\alpha}-\rF^{-\frac13})}}\|m\|_{0,\frac65,\Omega}\alpha_{\mathsf{DF}}^{-\frac12}\alpha_1.\)}}
This completes the rest of the proof.
\end{proof}

As a consequence of the previous analysis, we can state the following main result. 
\begin{theorem}[Well-posedness of the fully-coupled continuous problem]\label{th:wellposedness}
 Suppose that \begin{align*}{C^{\mathsf{DF}}_{c,\bu}}\bigl[{\|\bf(0)\|_{0,\Omega}}+C^{\mathsf{ADR}}_{c}\alpha_1\|m\|_{0,\frac65,\Omega}\bigr]  \leq r<\frac12,\quad \text{and
that } \quad L_{\cS}<1. \end{align*}
Then the coupled problem \eqref{eq:weak} has a unique solution $(\bu,p,\bzeta,\varphi) \in \bV \times \rQ\times \bW \times \Psi$. Moreover, there holds 
% \rrb{fit into margin}
\begin{align*}
    \|\bu\|_{\bV} + \|p\|_{\rQ} + \|\bzeta\|_{\bW} + \|\varphi\|_{\Psi} 
  \leq  C^{\mathsf{ADR}}_{c}\|m\|_{0,\frac65,\Omega}+&{C^{\mathsf{DF}}_{c,\bu}}\max \biggl\{{C^{\mathsf{DF}}_{c,\bu}}\bigl[{\|\bf(0)\|_{0,\Omega}}+C^{\mathsf{ADR}}_{c}\alpha_1\|m\|_{0,\frac65,\Omega}\bigr],\\&\bigl({C^{\mathsf{DF}}_{c,\bu}}\bigl[{\|\bf(0)\|_{0,\Omega}}+C^{\mathsf{ADR}}_{c}\alpha_1\|m\|_{0,\frac65,\Omega}\bigr]\bigr)^2\biggr\}, \label{estimate su}
\end{align*}
where the constants do not depend on the model parameters $\kappa,\rF,\varrho_0$. %,\varrho_0$. 
\end{theorem}
\begin{proof}
We first recall that Lemma \ref{lemma:ball} guaranties that $\cS$ maps $\rW$ to itself. Then, bearing in mind the Lipschitz-continuity of $\cS : \rW \to \rW$ given by Lemma \ref{lemma:continuity-S} along with the fact that $L_{\cS}<1$, a direct application of the Banach fixed-point theorem yields the existence of a unique fixed point $\bu \in \bV.$ Hence, a unique solution of the system \eqref{eq:weak}. In addition, the {\it a priori} estimate is a direct consequence of the continuous dependence on data provided by Theorems \ref{th:continuous-DF} and \ref{th:continuous-ARD}.
\end{proof}
\begin{remark}
   {Let us briefly remark that, due to the coupling between the Darcy--Forchheimer velocity and the transport {equation, our  fixed-point framework does not allow us to readily} obtain $\alpha$-independent estimates for $\|\hat{\bu}\|_{\bV}$. {This manifests itself} in the hypothesis of Theorem~\ref{th:continuous-ARD} {as well as in}  the Lipschitz constant $L_{\mathcal{S}}$ (cf.  \eqref{eq:Lipschitz.S_ARD}). }
\end{remark}
%%%%%%%%%%%%%%%%%%%%%%%%%%%%%%%%%%%%%%%
\section{Finite element discretisation}\label{sec:fe}
Let \( {\cT}_h \) be a regular simplicial mesh defined in a bounded Lipschitz domain \( \Omega \), consisting of tetrahedral elements (or triangular elements in two dimensions) \( K \) with diameter \( h_K \). We define the mesh size as \( h := \max\bigl(h_K : K \in {\cT}_h\bigr) \). Given an integer \( k \geq 0 \) and a generic element \( K \in {\cT}_h \), we denote by \( \mathbb{P}_k(K) \) the space of polynomials defined locally in \( K \) with a degree at most \( k \), and denote by \( \mathbb{P}_k(\cT_h) \) its global counterpart. In addition, let ${\mathbb{RT}_k(\cT_h)}$ denote the \( \mathbf{H}_{\Gamma_{\bullet}}({\vdiv}, \Omega) \)-discrete conforming subspace defined by the Raviart--Thomas  elements of order \( k \), i.e.,  
\[{\mathbb{RT}_k(\cT_h)} := \{ \bu \in \mathbf{H}_{\Gamma_{\bullet}}({\vdiv_s}, \Omega):\bu|_{K} = [\mathbb{P}_k(K)]^d \oplus \widetilde{\mathbb{P}}_k(K) \bx, \ \bx \in \Omega \subset \mathbb{R}^d, K \in {\cT}_h \},\]
where \( \widetilde{\mathbb{P}}_k(K) \) is the space of all polynomials over \( K \) of degree exactly equal to \( k.\)
%%%%%%%%%%%%%%%%%%%%%%%%%%%%%%%%%%%%%%%%%%%%%%%%%%%%%%%%%%%%%%%%%%%%%%%%%%%%%%%%%%%%%%%%%%%%%%%%%%%%%%%%%%%%%%%%%%%%%%%%%%%%%%%%%%%%%%%%%%%%%%%%%%%%%%%%%%%%%%%%%%%%%%%%%%%%%%%%%%%%%%%
{We follow \cite[Section~4]{dhpr_sinum26} for constructing discrete spaces for the system~\eqref{eq:weak}. Define conforming discrete weighted subspaces, \(\bV_h\subseteq\bV\) and \(\bW_h\subseteq\bW\) as follows 
\begin{subequations}
\begin{align}
 \bV_h&:=  \kappa^{-\frac{1}{2}}\bL_h^2(\Omega) \cap \bH_{h,\Gamma_{\bu}}(\vdiv,\Omega) \cap \rF^{\frac13} \bL_h^3(\Omega),\label{eq:discrete.Vh}\\
  \bW_h&:= \alpha^{-\frac12}\bL^2_h(\Omega)\cap\bH_{h,\Gamma_p}(\vdiv_{\frac65},\Omega), \label{eq:discrete.Wh}
 \end{align}\end{subequations}
 where \( {\bL}_h^r(\Omega), \,\,{\bH}_{{h,\Gamma_{\bullet}}}({\vdiv}, \Omega), \,{\bH}_{{h,\Gamma_{\bullet}}}({\vdiv}_s, \Omega) \)-norms correspond to  functions in ${\mathbb{RT}_k(\cT_h)}$ equipped with  ${\bL}^r(\Omega)$, for \(r=2,3,\) {$\bH_{\Gamma_{\bullet}}(\vdiv,\Omega),$ and $\bH_{\Gamma_{\bullet}}(\vdiv_{\frac65},\Omega)$}.
Next, define the discrete gradient operator \( \nabla_h: {\mathbb{P}_k(\cT_h)} \to {\mathbb{RT}_k(\cT_h)} \) (cf. \cite[Equation~3.1]{arnold1997preconditioning}) by 
\begin{equation}\label{eq:disc.grad}
\sum_{K\in\cT_h}\int_K(\nabla_h q_h)\cdot\bv_h = -\sum_{K\in\cT_h}\int_K(\vdiv\bv_h)\,q_h \quad \forall q_h \in {\mathbb{P}_k(\cT_h)}, \text{ and } \forall \bv_h \in {\mathbb{RT}_k(\cT_h)}.
\end{equation}
Additionally, define the following broken Sobolev spaces 
\begin{align*}
    {\rW^{1,r}}(\mathcal{T}_h)&:=\{p_h\in\rL^2(\Omega):p_h|_K\in\rW^{1,r}(K) \quad \forall K\in \mathcal{T}_h\}, \quad \text{for } r=2,\frac32,
\end{align*}
and whenever needed, we add a subscript $\Gamma_{\bullet}$ denoting a part on the boundary, where the value of the function is prescribed to zero.  
These spaces are equipped with DG (broken Sobolev) seminorms 
\(|\bullet|_{{1,\mathcal{T}_h}}\) and 
\(|\bullet|_{{1, r, \mathcal{T}_h}}\),
defined as 
\begin{align*}
{\rH^1}(\mathcal{T}_h) \ni \zeta & \mapsto  |\zeta|^2_{{1,\mathcal{T}_h}} := \biggl(
\sum_{K\in\mathcal{T}_h}\| {\nabla} \zeta\|^{2}_{0,K} 
+ \sum_{F\in \mathcal{F}_h} h_F^{-1}\|[\![\zeta]\!]\|^2_{0,F}
\biggr)^{\frac12}, 
\\
\rW^{1,r}(\mathcal{T}_h) \ni \xi & \mapsto 
{|\xi|_{{1,r,\mathcal{T}_h}} := \biggl(\sum_{K\in\mathcal{T}_h}\|{\nabla} \xi\|^{r}_{0,r,K}
+ \sum_{F\in \mathcal{F}_h} h_F^{1-r}\|[\![\xi]\!]\|^{r}_{0,r,F}
\biggr)^{\frac{1}{r}}},\end{align*}
respectively, where $\mathcal{F}_h$ denotes the set of facets in the mesh $\cT_h$. 
%
Further, define the following weighted spaces for the pressure, as the sum of the infinite-dimensional (but broken) spaces $\rL^2(\Omega)$, \(\rH^1_{\Gamma_p}(\mathcal{T}_h)\) and \(\rW^{1,\frac32}_{\Gamma_p}(\mathcal{T}_h),\) and the concentration, as the sum of \(\rL^6(\Omega)\) and \(\rH^1_{\Gamma_{\bu}}(\cT_h):\)  
\begin{subequations}
\begin{align}\label{eq:broken.Qhat}
{\widehat{\rQ}}(\cT_h)&:=\rL^2(\Omega)+\kappa^\frac12\rH^1_{\Gamma_p}(\mathcal{T}_h)+\rF^{-\frac13}\rW^{1,\frac32}_{\Gamma_p}(\mathcal{T}_h),\\
\widehat{\Psi}(\cT_h)&:=\rL^6(\Omega)+\alpha^{\frac12}\rH^1_{\Gamma_{\bu}}(\cT_h).
\end{align}\end{subequations}
Define broken weighted norms on \({\widehat{\rQ}}(\cT_h)\) and \(\widehat{\Psi}(\cT_h)\) (see, \cite[Section~4]{dhpr_sinum26}) as  
\begin{align}
\|q\|_{{\widehat{\rQ}(\cT_h)}}&:= 
\inf_{\substack{
    q = q_{1} + q_{2} + q_{3} \\
    (q_{1},q_{2},q_{3}) \in \\ \rL^2(\Omega)
     \times\rH_{\Gamma_p}^1(\mathcal{T}_h)
    \times \rW_{\Gamma_p}^{1,\frac32}(\mathcal{T}_h)
}}
\left(
    \|q_{1}\|_{0,\Omega} 
    +\kappa^{\frac12}\, |q_{2}|_{1,\mathcal{T}_h} 
    + \rF^{-\frac13} |q_{3}|_{1,\frac32,\mathcal{T}_h}
\right),\label{norm:Q.hat}\\
\|\psi\|_{\widehat{\Psi}(\cT_h)}&:=\inf_{\substack{
    \psi = \psi_{1} + \psi_{2} \\
    (\psi_{1},\psi_{2}) \in \\ \rL^6(\Omega)
     \times\rH_{\Gamma_{\bu}}^1(\mathcal{T}_h)
}}
\left(
    \|\psi_{1}\|_{0,6,\Omega} 
    +\alpha^{\frac12}\, |\psi_{2}|_{1,\mathcal{T}_h}  
\right).\label{norm:Psi.hat}
\end{align}
Finally, we set  discrete spaces \(\rQ_h\) and \(\Psi_h\) (non-conforming subspaces of   \(\rQ\) and \(\Psi\),  respectively) as    
\begin{subequations}
\begin{align}
\rQ_h &:=  \rL_h^2(\Omega) +\kappa^{\frac12} \rH^1_{h,{\Gamma_p}}(\Omega)+\rF^{-\frac13}\rW^{1,\frac32}_{h,\Gamma_p}(\Omega),\label{eq:disc.Ph}\\
\Psi_h&:=\rL^6_h(\Omega)+\alpha^{\frac12}\rH^1_{h,\Gamma_{\bu}}(\Omega),\label{eq:disc.Psih}
\end{align}
\end{subequations}
where \( \rL_h^r(\Omega), \rH^1_h(\Omega), \) and \(\rW^{1,\frac32}_h(\Omega)\) are polynomial spaces having discrete functions from  \( {\mathbb{P}_k(\cT_h)} \), and equipped with the \( \rL^r(\Omega)\) norm (for \(r=2,6\))  and the broken \(\rH^1(\cT_h),\,\rW^{1,\frac32}(\cT_h)\) seminorms, respectively. 
We equip \(\rQ_h\) and \(\Psi_h\) with the $\|\bullet\|_{\widehat{\rQ}(\cT_h)}$ and $\|\bullet\|_{\widehat{\Psi}(\cT_h)}$ norms, respectively. 
Next, set discrete bilinear forms  
$b_{1,h}: \bV_h\times \rQ_h \to \mathbb{R}, b_{2,h}: \bW_h\times \Psi_h \to \mathbb{R}$ 
 and the discrete linear functional \(H_h:\Psi_h\to\mathbb{R}\) as follows
\begin{align*}
b_{1,h}(\bv_h,q_h)&:=\int_{\Omega}\vdiv\bv_h\,q_h,\quad
b_{2,h}(\bzeta_h,\psi_h):=\int_{\Omega}\vdiv\bzeta_h\psi_h, \quad H_h(\psi_h):=-\int_{\Omega}m\psi_h, 
\end{align*}
for all \( (\bv_h\,  q_h, \bzeta_h,\psi_h)\in \bV_h\times {\rQ_h} \times \bW_h\times\Psi_h.\)
Then, with the  FE spaces \eqref{eq:discrete.Vh}, \eqref{eq:discrete.Wh}, \eqref{eq:disc.Ph} and \eqref{eq:disc.Psih} 
 the mixed FE formulation of \eqref{eq:weak} consists in  finding $(\bu_h,p_h,\bzeta_h,\varphi_h)\in \bV_h\times \rQ_h\times\bW_h\times\Psi_h$ such that 
\begin{subequations}\label{eq:weak_Galerkin}
    \begin{align}
     a_1(\bu_h,\bv_h) + d_1(\bu_h; \bu_h,\bv_h) +  b_{1,h}(\bv_h,p_h) & = F^{\varphi_h}(\bv_h) \qquad \forall \bv_h \in \bV_h,\label{eq:weak,1h}\\
    b_{1,h}(\bu_h,q_h) & = 0 \qquad\qquad \forall q_h \in \rQ_h,\label{eq:weak,2h}\\
    a_2(\bzeta_h,\bw_h) + b_{2,h}(\bw_h,\varphi_h) + d_2^{\bu_h}(\bw_h, \varphi_h) & = 0\qquad \qquad \forall \bw_h \in \bW_h, \label{eq:weak,3h}\\
    b_{2,h}(\bzeta_h, \psi_h) - c_{2,h}(\varphi_h,\psi_h) & = H_h(\psi_h) \qquad \forall \psi_h \in \Psi_h. \label{eq:weak,4h}
    \end{align}
\end{subequations}

\noindent
{A direct discrete analysis based on
 \(\bV_h:=(\mathbb{RT}_k(\cT_h),\|\bullet\|_{\vdiv,\Omega}),\) \(\bW_h:=(\mathbb{RT}_k(\cT_h),\|\bullet\|_{\vdiv_{\frac65},\Omega}),\) \(\rQ_h:=(\mathbb{P}_k(\cT_h),\|\bullet\|_{0,\Omega}),\) and \(\Psi_h:=(\mathbb{P}_k(\cT_h),\|\bullet\|_{0,6,\Omega})\), 
endowed with the standard norms, does not lead to parameter-robust inf-sup and Braess conditions.
The parameter robustness is instead obtained by adopting the weighted discrete spaces \eqref{eq:discrete.Vh}-\eqref{eq:discrete.Wh}, \eqref{eq:disc.Ph}-\eqref{eq:disc.Psih}.}  

\subsection{Robust discrete solvability}
 Consider the following system for a fixed $\hat{\varphi}_h\in \Psi_h$ 
\begin{subequations}\label{eq:weak.fixed.u.hat.disc}
    \begin{align}
    a_1(\bu_h,\bv_h) + d_1(\bu_h; \bu_h,\bv_h) +  b_{1,h}(\bv_h,p_h) & = F^{\hat\varphi_h}(\bv_h) \quad \forall \bv_h \in \bV_h,\label{eq:weak,1.hat.u.disc}\\
    b_{1,h}(\bu_h,q_h) & = 0 \quad \forall q_h \in \rQ_h,\label{eq:weak,2.hat.u.disc}
    \end{align}
\end{subequations}
and recall from \cite[Theorem~4.3]{dhpr_sinum26} the following result.
\begin{theorem}[Well-posedness of discrete  Darcy--Forchheimer equations]\label{th:discrete-DF}
 Assume that 
 $\bf(\hat{\varphi}_h)\in \bL^2(\Omega).$ Then there exists a unique  solution $(\bu_h,p_h)\in \bV_h\times\rQ_h$ to   \eqref{eq:weak,1h}-\eqref{eq:weak,2h}. Furthermore, there is a positive constant ${C^{\mathsf{DF}}_{\mathsf{d}}},$ independent of the parameters $\kappa, \rF,h$,  such that 
\begin{equation*}
 \|(\bu_h, p_h )\|_{\bV\times\widehat{\rQ}(\cT_h)}  \leq {C^{\mathsf{DF}}_{\mathsf{d}}}\max \bigl(\|\bf(\hat{\varphi}_h)\|_{0,\Omega},\|\bf(\hat{\varphi}_h)\|_{0,\Omega} ^2\bigr).
    \end{equation*}
Additionally, there exists a positive constant \(C^{\mathsf{DF}}_{\mathsf{d},\bu_h},\) independent of \(\kappa,\rF, \bu_h,\) and \(h\), such that 
    \begin{align}\label{est:disc.uh.bound}
        \|\bu_h\|_{\bV}\leq C^{\mathsf{DF}}_{\mathsf{d},\bu_h}\,\|\bf(\hat{\varphi}_h)\|_{0,\Omega}.
    \end{align}
\end{theorem}

\noindent We next prove the following lemma, which is crucial in deducing the discrete inf-sup stability of \(b_{2,h}(\bullet,\bullet)\).
% (\cite{colmenares2020banach}, \cite{benavides2022new}, \cite{caucao2021fully})
\begin{lemma}\label{lem:disc.inf-sup:ARD}
   There exists an operator \(\cS_{\mathsf{d}}\) satisfying the following properties:
    \begin{gather}\nonumber
    \cS_{\mathsf{d}} \in \cL(\Psi_h', \bW_h),\quad 
    \|\cS_{\mathsf{d}}\|_{\cL(\Psi_h', \bW_h)}\ \text{is independent of } \alpha,\varrho_0, h\quad \text{and}\\
    \langle{{\vdiv}_{\frac65}} [\cS_{\mathsf{d}}(\ell)], \psi_h\rangle = \langle \ell, \psi_h \rangle \quad \text{for all}\ \ell \in \Psi_h',\psi_h \in \Psi_h.\label{S:disc.prop}
\end{gather}
\end{lemma}
\begin{proof}
{First, construct \(\cS_{\mathsf{d}}\) in such a way that it satisfies the following properties 
% \rrb{are all these duals of discrete spaces well-defined? and, are they discrete spaces themselves? I ask because that the maps need to be surjective between discrete spaces. This is not clear from the proof below. No interpolation/projection is taken \dots} 
\[\cS_{\mathsf{d}}\in\cL\bigl((\rH^1_{h,\Gamma_{\bu}}(\Omega))', (\bL_h^2(\Omega)\bigr),\quad \text{and} \quad  \cS_{\mathsf{d}}\in\cL\bigl((\rL_h^6(\Omega))', \bH_{h,\Gamma_p}(\vdiv_{\frac65},\Omega)\bigr),\]respectively, and then using the fact that (cf. \eqref{eq:duality}) \(\Psi_h'=(\rH^1_{h,\Gamma_{\bu}}(\Omega)+\rL_h^6(\Omega))'=(\rH^1_{h,\Gamma_{\bu}}(\Omega))'\cap(\rL_h^6(\Omega))',\) we conclude that \( \cS_{\mathsf{d}} \in {\cL}(\Psi_h', \bW_h) \).
The existence of \(\cS_{\mathsf{d}}\in\cL\bigl((\rH^1_{h,\Gamma_{\bu}}(\Omega))',\bL_h^2(\Omega)\bigr)\)
follows directly from \cite[Lemma~4.1]{dhpr_sinum26}.} 

Next, we aim to construct \(\cS_{\mathsf{d}}\in\cL\bigl((\rL_h^6(\Omega))',\bH_{h,\Gamma_p}(\vdiv_{\frac65},\Omega)\bigr)\).
{Let \(q_h\in\rL^6_h(\Omega),\) and \(\mathcal{O}\) be a convex bounded domain containing \(\overline{\Omega}.\) Now, define 
\begin{align}\label{def:ell}
    \ell :=
\begin{cases}
|q_h|^4\, q_h, & \text{in } \Omega,\\[2pt]
0, & \text{in } \mathcal{O} \setminus \overline{\Omega}.
\end{cases}
\end{align}
Note that, \(\ell\in\rL^{\frac65}(\mathcal{O}),\) satisfying the following 
\begin{align}\label{est:ell.disc}
\|\ell\|_{0,\frac65,\Omega}=\||q_h|^4q_h\|_{0,\frac65,\Omega}=\|q_h\|^5_{0,6,\Omega}.
\end{align}
It now follows from \cite[Corollary~1]{fromm1993potential} that, there exists a unique \(z\in \rW_{\Gamma_{\bu}}^{1,\frac65}(\mathcal{O})\cap\rW^{2,\frac65}(\mathcal{O})\) satisfying the following boundary value problem
\begin{align}\label{eq:disc.inf-sup.welp}
    \Delta z = \ell \quad \text{in } \Omega,\quad z=0\quad \text{on } \Gamma_{\bu},\quad \nabla z\cdot\bn=0\quad \text{on } \Gamma_p,
\end{align}
and additionally, there exists a positive constant \(C_{\mathsf{reg}}\) (depending only on \(d\) and \(\Omega\)) such that \(z\) satisfies the following  stability estimate (the last equality is due to \eqref{est:ell.disc})
\begin{align}\label{est:z}
\|z\|_{2,\frac65,\mathcal{O}}\leq C_{\mathsf{reg}}\|\ell\|_{0,\frac65,\Omega}=C_{\mathsf{reg}}\|q_h\|^5_{0,6,\Omega}.
\end{align}
Next, set \(\varsigma:=(\nabla z)|_{\Omega}\in\rW^{1,\frac65}(\Omega),\) and note from \eqref{eq:disc.inf-sup.welp} and \eqref{est:z} that 
\begin{align}
    \vdiv(\varsigma)&=\ell=|q_h|^4q_h \quad \text{in } \Omega, \label{est:div.varsigma}\\ \|\varsigma\|_{1,\frac65,\Omega}&\leq\|z\|_{2,\frac65,\Omega}\leq C_{\mathsf{reg}}\|q_h\|^5_{0,6,\Omega}.\label{est:varsigm} 
\end{align}
In turn, let us now define \(\varsigma_h:=\Pi^k_h(\varsigma)\) (\( \Pi_h^k\) being the Raviart--Thomas projection operator satisfying \cite[Equation~5.38, Equation~5.39]{colmenares2020banach}). Further, using the boundedness of the projection operator,  \eqref{est:varsigm}, and finally the continuous injection of \(\rW^{1,\frac65}(\Omega)\) into \(\rL^2(\Omega),\)  we obtain the following estimate (the absorbed constant being independent of \(h\))
\begin{align}
\|\varsigma_h\|_{0,\Omega}=\|\Pi_h^k(\varsigma)\|_{0,\Omega}\le C_{\mathsf{reg}}\|q_h\|^5_{0,6,\Omega}.
\end{align}
In addition, from \cite[Equation~(5.35)]{colmenares2020banach} and \eqref{est:div.varsigma}, we obtain the following commutativity of the divergence operator  and the projection operator
\begin{align}
    \vdiv(\varsigma_h)=\vdiv(\Pi^k_h(\varsigma))=\Pi_h^k(\vdiv(\varsigma))=\Pi_h^k(|q_h|^4q_h),
\end{align}
and further, using \eqref{est:div.varsigma}, and the boundedness of \(\Pi_h^k,\) it follows that
\begin{align}\label{est:div.triang}
\|\vdiv(\varsigma_h)\|_{0,\frac65,\Omega}\leq\|q_h\|^5_{0,6,\Omega}.
\end{align}
Then, it suffices to define \(\mathcal{S}_d(\ell):=\varsigma_h,\) which further implies that \(\cS_{\mathsf{d}}\in\cL\bigl((\rL_h^6(\Omega))',\bH_{h,\Gamma_p}(\vdiv_{\frac65},\Omega)\bigr).\) This completes the rest of the proof.
}
\end{proof}

\noindent
We now consider the following {discrete ADR} system for a fixed $\hat{\bu}_h \in \bV_h$
\begin{subequations}\label{eq:weak.hat.phi.disc}
    \begin{align}
    a_2(\bzeta_h,\bw_h) + b_{2,h}(\bw_h,\varphi_h) + d_2^{\hat\bu_h}(\bw_h, \varphi_h) & = 0 \quad \forall \bw_h \in \bW_h, \label{eq:weak,3.hat.phi.disc}\\
    b_{2,h}(\bzeta_h, \psi_h) - c_{2,h}(\varphi_h,\psi_h) & = H_h(\psi_h) \quad \forall \psi_h \in \Psi_h, \label{eq:weak,4.hat.phi.disc}
    \end{align}
\end{subequations}
and employ the discrete version of Theorem~\ref{th:abstract} to discuss its parameter-robust unique solvability.
\begin{theorem}[Well-posedness of discrete  advection-diffusion-reaction equations]\label{th:discrete-ARD}
Assume that 
 $m \in \rL^{\frac65}(\Omega)$. {Then, for a fixed $\hat{\bu}_h\in \bV_h,$ there {exist}  positive constants \(C_d, C^d_{\alpha}\) (depending on $\alpha$), such that, whenever {\(\|\hat\bu_h\|_{{0,3,\Omega}}\leq 
 \frac{C_d}{2C^d_{\alpha}},\)}} there exists a unique  solution $(\bzeta_h,\varphi_h)\in \bW_h\times\Psi_h$ to the problem \eqref{eq:weak.hat.phi.disc}. 
 % \rrb{[should it be of (5.9)?]}
 Furthermore, there is a positive constant $C^{\mathsf{ADR}}_{\mathsf{d}}$ independent of \(h,\) and  $\varrho_0$, such that 
\begin{equation*}
 \|(\bzeta_h, \varphi_h)\|_{\bW\times\widehat{\Psi}(\cT_h)}  \leq C^{\mathsf{ADR}}_{\mathsf{d}} \| m \|_{0,\frac65,\Omega} .
    \end{equation*}
\end{theorem}
\begin{proof}
     The boundedness of \(a_2(\bullet,\bullet),b_{2,h}(\bullet,\bullet),\,c_{2,h}(\bullet,\bullet),\) and  symmetry and positive definiteness of \(a_2(\bullet,\bullet)\) and \(c_{2,h}(\bullet,\bullet),\) and the inf-sup stability of \(a_2(\bullet,\bullet)\) in \(\bZ_h:=\{\btau_h\in\bW_h:b_{2,h}(\btau_h,\psi_h)=0,\,\forall \psi_h\in\Psi_h\}\) follow analogously as in Theorem~\ref{th:continuous-ARD}. Therefore, to complete the rest of the proof, it is enough to prove the discrete inf-sup stability of \(b_{2,h}(\bullet,\bullet)\) and the discrete Braess condition.

\smallskip
\noindent\textbf{Discrete inf-sup condition.} %This follows directly from Lemma~\ref{lem:disc.inf-sup:ARD}. 
From the existence of $\cS_{\mathsf{d}}$ defined in Lemma~\ref{lem:disc.inf-sup:ARD}, it %straightaway 
    follows that
\begin{align*}
\sup_{\bzero\neq\btau_h \in \bW_h} \!\!\frac{\langle{\vdiv} \btau_h, \psi_h\rangle}{\|\btau_h\|_\bW} &\geq \sup_{0\neq\ell \in \Psi_h'} \frac{\langle{{\vdiv}_{\frac65}} \cS_{\mathsf{d}}(\ell), \psi_h\rangle}{\|\cS_{\mathsf{d}}(\ell)\|_\bW} = \sup_{0\neq\ell \in \Psi_h'} \frac{\langle \ell, \psi_h \rangle}{\|\cS_{\mathsf{d}}(\ell)\|_\bW} \\&\geq \|\cS_{\mathsf{d}}\|^{-1}_{\cL(\Psi_h', \bW_h)} \sup_{0\neq\ell \in \Psi_h'} \frac{\langle \ell, \psi_h \rangle}{\|\ell\|_{\Psi_h'}} = \tilde{C}_b \|\psi_h\|_{\widehat{\Psi}(\cT_h)},
\end{align*}
where $\tilde{C}_b = \|\cS_{\mathsf{d}}\|^{-1}_{\cL(\Psi_h', \bW_h)},$ is independent of \(\alpha,\varrho_0,h\), and hence the discrete stability of inf-sup is achieved.

\smallskip
\noindent\textbf{Discrete Braess condition.} Setting \(\bw_h=\bzeta_h\), \(\psi_h=0,\) and exploiting \(\|\bzeta_h\|_{\bW}=\tnorm{(\bzeta_h,0)}_{\bW\times\widehat{\Psi}(\cT_h)},\) one obtains 
\begin{align}\label{eq:Braess.disc.1st} \sup_{(\bzero,0)\neq(\bw_h,\psi_h)\in\bW_h\times\Psi_h}\frac{a_2(\bzeta_h,\bw_h)+b_{2,h}(\bzeta_h,\psi_h)}{\tnorm{(\bw_h,\psi_h)}_{\bW\times\widehat{\Psi}(\cT_h)}}\geq\frac{\alpha^{-1}\int_\Omega\bzeta_h\cdot\bzeta_h}{\tnorm{(\bzeta_h,0)}_{\bW\times\widehat{\Psi}(\cT_h)}}=\frac{\alpha^{-1}\|\bzeta_h\|_{0,\Omega}^2}{\|\bzeta_h\|_{\bW}}.
    \end{align}
     
Next, {taking the particular values} \(\bw_h=\bzero,\) and \(\psi_h=\mathrm{sgn}(\vdiv\bzeta_h)|\vdiv\bzeta_h|^{\frac15}\), %in the left-hand side of \eqref{eq:Braess.disc.1st}, 
one obtains the following 
\begin{align}\label{eq:Braess.disc.2nd}  
(\mathrm{LHS})_{\eqref{eq:Braess.disc.1st}}\geq \frac{a_2(\bzeta_h,\bzero)+b_{2,h}(\bzeta_h,\mathrm{sgn}(\vdiv\bzeta_h)|\vdiv\bzeta_h|^{\frac15})}{\tnorm{(\bzero,\mathrm{sgn}(\vdiv\bzeta_h)|\vdiv\bzeta_h|^{\frac15})}_{\bW\times\widehat{\Psi}(\cT_h)}}\geq\frac12 \frac{\|\vdiv \bzeta_h\|^{\frac65}_{0,\frac65,\Omega}}{\|\vdiv \bzeta_h\|^{\frac15}_{0,\frac65,\Omega}}=\frac12\|\vdiv \bzeta_h\|_{0,\frac65,\Omega}.
\end{align}
Now, let us combine \eqref{eq:Braess.disc.1st} and \eqref{eq:Braess.disc.2nd} to obtain the following
\begin{align}\label{eq:Braess.disc.3rd}
\sup_{(\bzero,0)\neq(\bw_h,\psi_h)\in\bW_h\times\Psi_h}\frac{a_2(\bzeta_h,\bw_h)+b_{2,h}(\bzeta_h,\psi_h)}{\tnorm{(\bw_h,\psi_h)}_{\bW\times\widehat{\Psi}(\cT_h)}}\geq\max\biggl({\frac{\alpha^{-1}\|\bzeta_h\|^2_{0,\Omega}}{\|\bzeta_h\|_{\bW}},\frac12\|\vdiv \bzeta_h\|_{0,\frac65,\Omega}}\biggr).
\end{align}
Denoting \(\gamma:=\frac{\alpha^{-\frac12}\|\bzeta_h\|_{0,\Omega}}{\|\bzeta_h\|_{\bW}}\in[0,1)\), and using \eqref{eq:Braess.disc.3rd}, we establish the following inequality
\begin{align*}\max\biggl({\frac{\alpha^{-1}\|\bzeta_h\|^2_{0,\Omega}}{\|\bzeta_h\|_{\bW}},\|\vdiv \bzeta_h\|_{0,\frac65,\Omega}}\biggr)\geq \|\bzeta_h\|_{\bW}\biggl(\max_{t\in[0,1)}\bigl[\gamma^2,\frac12\sqrt{1-\gamma^2}\bigr]\biggr). \end{align*}
On the other hand, if we  denote \(\gamma_{\mathsf{d}}:=\biggl(\max_{t\in[0,1)}\bigl[\gamma^2,\frac12\sqrt{1-\gamma^2}\bigr]\biggr),\) we can readily conclude that \[\sup_{(\bzero,0)\neq(\bw_h,\psi_h)\in\bW_h\times\Psi_h}\frac{a_2(\bzeta_h,\bw_h)+b_{2,h}(\bzeta_h,\psi_h)}{\tnorm{(\bw_h,\psi_h)}_{\bW\times\widehat{\Psi}(\cT_h)}}\geq \gamma_{\mathsf{d}} \|\bzeta_h\|_{\bW}.\]
An appeal to  Theorem~\ref{th:abstract} now yields the well-posedness of the discrete reaction-diffusion system~\eqref{eq:weak,3h} and \eqref{eq:weak,4h} without the advection part. Additionally, the following estimate holds 
\begin{align}\label{eq:A.stab.disc}
\sup_{(\bzero,0)\neq(\tau_h,\psi_h)\in\bW_h\times\Psi_h}\frac{\cA(\bzeta_h,\varphi_h;\btau_h,\psi_h)}{\tnorm{(\btau_h,\psi_h)}_{\bW\times \widehat{\Psi}(\cT_h)}}\geq C_d\tnorm{(\bzeta_h,\varphi_h)}_{\bW\times \widehat{\Psi}(\cT_h)},
\end{align}
where \(C_d\) is the positive constant from 
\eqref{eq:C},   independent of the model and   discretisation parameters.

Finally, for the well-posedness of the uncoupled discrete ADR equations, for a fixed \(\hat{\bu}_h\in\bV_h,\) we consider another bilinear form \(\mathcal{A}^{\mathsf{ADR}}_{\mathsf{d}}:\bW_h\times\Psi_h\to \mathbb{R}\) defined by 
\begin{align}\label{eq:A:ARD.disc}\mathcal{A}_{\mathsf{d}}^{\mathsf{ADR}}(\bzeta_h,\varphi_h;\btau_h,\psi_h):=\mathcal{A}(\bzeta_h,\varphi_h;\btau_h,\psi_h)+d_2^{\hat{\bu}_h}(\btau_h,\varphi_h),\quad \forall (\bzeta_h,\varphi_h), (\btau_h,\psi_h)\in\bW_h\times \Psi_h.\end{align}
Thus,  to show the well-posedness of \eqref{eq:weak,3h}-\eqref{eq:weak,4h}, we now analyse  the problem: Find \((\bzeta_h,\varphi_h)\in\bW_h\times \Psi_h\) such that \[\mathcal{A}_{\mathsf{d}}^{\mathsf{ADR}}(\bzeta_h,\varphi_h;\btau_h,\psi_h)=H_h(\psi_h),\quad \forall \psi_h\in \Psi_h.\] 
From the boundedness of \(d_2^{\hat{\bu}_h}(\bullet,\bullet)\),{\begin{align}\label{eq:d2.bound.disc}d_2^{\hat{\bu}_h}(\btau_h,\varphi_h)\leq C^d_{\alpha}\|\hat{\bu}_h\|_{0,3,\Omega}\,\tnorm{(\btau_h,\psi_h)}_{\bW\times\widehat{\Psi}(\cT_h)}\tnorm{(\bzeta_h,\varphi_h)}_{\bW\times\widehat{\Psi}(\cT_h)},\end{align}} 
{where $C^d_{\alpha}=\alpha^{-\frac12}(1+{C_{\mathsf{emb}}}\alpha^{-\frac12})$}, combined with \eqref{eq:A.stab.disc} and \eqref{eq:A:ARD.disc}, we can assert that
\[\sup_{(\bzero,0)\neq(\btau_h,\psi_h)\in {\bW_h\times\Psi_h}}\frac{\mathcal{A}_{\mathsf{d}}^{\mathsf{ADR}}(\bzeta_h,\varphi_h;\btau_h,\psi_h)}{\tnorm{(\btau_h,{\psi_h})}_{\bW\times\widehat{\Psi}(\cT_h)}}\geq\bigl[C_d-C^d_{\alpha}\|\hat{\bu}_h\|_{0,3,\Omega}\bigr]\tnorm{(\bzeta_h,\varphi_h)}_ {\bW\times\widehat{\Psi}(\cT_h)}, \]
for all \((\bzeta_h,\varphi_h)\in \bW_h\times \Psi_h.\) Hence, {under the assumption that  
$\|\hat{\bu}_h\|_{0,3,\Omega}\leq \frac{C_d}{2C^d_{\alpha}}$, we} arrive at 
\begin{equation}\label{eq:ARD.disc}
\sup_{(\bzero,0)\neq(\btau_h,\psi_h)\in {\bW_h\times\Psi_h}}\frac{\mathcal{A}_{\mathsf{d}}^{\mathsf{ADR}}(\bzeta_h,\varphi_h;\btau_h,\psi_h)}{\tnorm{(\btau_h,\psi_h)}_ {\bW\times\widehat{\Psi}(\cT_h)}}\geq\frac{C_d}{2} \tnorm{(\bzeta_h,\varphi_h)}_ {\bW\times \widehat{\Psi}(\cT_h)}.
\end{equation}
Therefore, the discrete $\mathsf{ADR}$ system~\eqref{eq:weak,3h}-\eqref{eq:weak,4h} is well-posed. 
Additionally, invoking \eqref{eq:A:ARD.disc} in \eqref{eq:ARD.disc}, and  applying H\"{o}lder's inequality, 
% to note that \(H(\psi_h)\leq\|m_h\|_{0,\frac65,\Omega}\|\psi_h\|_{0,6,\Omega},\) 
gives the dependence on data {\(\tnorm{(\bzeta_h,\varphi_h)}_{\bW\times\widehat{\Psi}(\cT_h)}\leq C^{\mathsf{ADR}}_{\mathsf{d}}\|m\|_{0,\frac65,\Omega},\)} where {\(C^{\mathsf{ADR}}_{\mathsf{d}}:={2}{C_d}^{-1}(1+C_{\mathsf{emb}}\alpha^{-\frac12})\)}, which concludes the proof. 
\end{proof}
}
\subsection{Fixed-point analysis in discrete framework}
{Analogous to the continuous case (cf. \eqref{Operator-S.V->V}), we now define the fixed-point operator in the discrete framework as follows  
\begin{equation}\label{Operator-S_h.Vh->Vh}
\cS_h: \bV_h  \to  \bV_h, \quad \hat{\bu}_h \mapsto \cS_h(\hat{\bu}_h):= \cS_{1,h}^{\mathsf{DF}}(\cS_{2,h}^{\mathsf{ADR}}(\hat{\bu}_h)),\quad \hat{\bu}_h\in\bV_h,
\end{equation}
where the maps~\(\cS_{1,h}^{\mathsf{DF}_1}\) and \( \cS_{2,h}^{\mathsf{ADR}_2}\) are defined over the respective discrete spaces as follows
\[ \cS_h^{\mathsf{DF}}:\Psi_h\to \bV_h \times \rQ_h, \quad \hat{\varphi}_h \mapsto \cS_h^{\mathsf{DF}}(\hat{\varphi}_h)= (\cS^{\mathsf{DF}}_{1,h}(\hat{\varphi}_h), \cS^{\mathsf{DF}}_{2,h}(\hat{\varphi}_h)):= (\bu_h, p_h), \]
where $( \bu_h,p_h) \in \bV_h \times \rQ_h$ is the unique solution of the Darcy--Forchheimer equations as stated in Theorem~\ref{th:discrete-DF}.
Further, we define the solution operator associated with the mixed ADR equations as
\[ \cS_h^{\mathsf{ADR}}: \bV_h \to  \bW_h\times \Psi_h, \quad \hat{\bu}_h \mapsto \cS_h^{\mathsf{ADR}}(\hat{\bu}_h)=(\cS^{\mathsf{ADR}}_{1,h}(\hat{\bu}_h),\cS^{\mathsf{ADR}}_{2,h}(\hat{\bu}_h)):= ( \bzeta_h, \varphi_h),\]
where $(\bzeta_h, \varphi_h)$ is the unique solution of the ADR equations, as stated in Theorem~\ref{th:discrete-ARD}.

A parallel set of arguments can be applied here likewise the continuous case to show that the discrete fixed-point map~\(\cS_h\) follows the subsequent properties
\begin{itemize}
    \item For a positive real number \(0<\tilde r<\frac12,\) the map \(\cS_h\) preserves the ball-mapping property, that is, \(\cS_h(\mathrm{B}^{\tilde r}_{\bzero})\subseteq\mathrm{B}^{\tilde r}_{\bzero}\), under the {small data assumption}   $${C^{\mathsf{DF}}_{\mathsf{d},\bu_h}}
    %^{\eqref{est:disc.uh.bound}}
    \bigl({\|\bf(0)\|_{0,\Omega}}+{C^{\mathsf{ADR}}_{\mathsf{d}}\alpha_1}\|m\|_{0,\frac65,\Omega}\bigr)  \leq \tilde r<\frac12,$$
where $C^{\mathsf{DF}}_{\mathsf{d},\bu_h}$ comes from \eqref{est:disc.uh.bound}.
    \item \(\cS_h\) is Lipschitz continuous with the Lipschitz constant $L^{\mathsf{d}}_{\cS}:=\frac{2C^d_{\alpha}}{C_d(C^d_{\alpha}-\rF^{-\frac13})}\|m\|_{0,\frac65,\Omega}{\alpha}_{\mathsf{DF}}^{-\frac12}\alpha_1.$    
\end{itemize}
\noindent
We next state the discrete stability theorem for the fully-coupled discrete $\mathsf{DF}$-$\mathsf{ADR}$ system similar to the continuous case as in Theorem~\ref{th:wellposedness}.}
\begin{theorem}[Well-posedness of the fully-coupled discrete problem]\label{th:wellposedness.disc}
 Suppose that $${C^{\mathsf{DF}}_{\mathsf{d},\bu_h}}\bigl({\|\bf(0)\|_{0,\Omega}}+{C^{\mathsf{ADR}}_{\mathsf{d}}\alpha_1}\|m\|_{0,\frac65,\Omega}\bigr)  <\frac12,$$ and
that $\displaystyle L^{\mathsf{d}}_{\cS}<1$.  Then, the coupled problem \eqref{eq:weak_Galerkin} has a unique solution $(\bu_h,p_h,\bzeta_h,\varphi_h) \in \bV_h \times \rQ_h\times \bW_h \times \Psi_h$. Moreover, there are 
{{\begin{align*}
   \|\bu_h\|_{\bV_h} + \|p_h\|_{\widehat{\rQ}(\cT_h)} + \|\bzeta_h\|_{\bW_h} + &\|\varphi_h\|_{\widehat{\Psi}(\cT_h)} 
  \leq  C^{\mathsf{ADR}}_{\mathsf{d}}\|m\|_{0,\frac65,\Omega}+C^{\mathsf{DF}}_{\mathsf{d},\bu_h}\max \bigl[{C^{\mathsf{DF}}_{\mathsf{d},\bu_h}}\bigl[{\|\bf(0)\|_{0,\Omega}}\\&\quad + {C^{\mathsf{ADR}}_{\mathsf{d}}\alpha_1}\|m\|_{0,\frac65,\Omega} \bigr],\bigl({{C^{\mathsf{DF}}_{\mathsf{d},\bu_h}}}\bigl[{\|\bf(0)\|_{0,\Omega}}+{C^{\mathsf{ADR}}_{\mathsf{d}}\alpha_1}\|m\|_{0,\frac65,\Omega}\bigr]\bigr)^2\bigr], 
\end{align*}}}
where the constants do not depend on the model parameters $\kappa,\rF,\varrho_0,$ nor the discretisation parameter \(h.\) 
\end{theorem}
}
In the next section, we estimate the error between the exact and approximate solutions in the proposed weighted-norms, showing that the convergence is robust with respect to the material parameters.
\subsection{ {\it{A priori}} error analysis}\label{sec:apriori}
We start this section by proving the following parameter-robust \textit{a priori} error estimate for the coupled {system \eqref{eq:strong}. % with respect to \(\kappa, \rF, \alpha, \varrho_0,\) and \(h.\)
% \noindent
% Let \((\bu,p,\bzeta,\varphi)\in \bV\times \rQ\times \bW\times \Psi\) uniquely solve the continuous system~\eqref{eq:weak} and \((\bu_h,p_h,\bzeta_h,\varphi_h)\in \bV_h\times \rQ_h\times \bW_h\times \Psi_h\) uniquely solve the discrete system~\eqref{eq:weak_Galerkin}. 
% % To simplify the subsequent analysis, we denote by \(e_{\bu},e_p,e_{\bzeta},e_{\varphi}\) the following differences \[e_{\bu}:=\bu-\bu_h,\,e_p:=p-p_h,\,e_{\bzeta}:=\bzeta-\bzeta_h,\,e_{\varphi}:=\varphi-\varphi_h.\]
\begin{theorem}[Quasi-optimality]\label{th:Cea}
    Let \((\bu,p,\bzeta,\varphi)\) uniquely solve the system~\eqref{eq:weak} and \((\bu_h,p_h,\bzeta_h,\varphi_h)\) uniquely solve the system~\eqref{eq:weak_Galerkin}. Additionally, suppose that the assumptions of Theorem~\ref{th:wellposedness} and Theorem~\ref{th:wellposedness.disc} hold. Also, suppose that the following relation holds
    \[
       \rF^{-\frac13}(1+C_{\mathsf{emb}}\alpha^{-\frac12})C_c^{\mathsf{ADR}}C_{\mathsf{d}}^{\mathsf{ADR}}\|m\|_{0,\frac65,\Omega}\leq\frac12,
    \]
     with \(\alpha_{\mathsf{DF}}\) being the uniform strong monotonicity constant (cf. \cite[Equation~(4.9)]{dhpr_sinum26}). 
    Then there exists a positive constant \(C_{\mathsf{C\acute{e}a}}\) such that the following quasi-optimal error estimate holds %for the velocity, pressure, flux, and concentration 
    \begin{align}\label{eq:Cea}
        \|\bu-\bu_h\|_{\bV}&+\|p-p_h\|_{\widehat{\rQ}(\cT_h)}+\|\bzeta-\bzeta_h\|_{\bW}+\|\varphi-\varphi_h\|_{\widehat{\Psi}(\cT_h)}\nonumber\\
        &\qquad \leq C_{\mathsf{C\acute{e}a}}\bigl(\mathrm{dist}(\bu,\bV_h)+{(\mathrm{dist}(\bu,\bV_h))^2}+\mathrm{dist}(p,\rQ_h)+\mathrm{dist}(\bzeta,\bW_h)+\mathrm{dist}(\varphi,\Psi_h) \bigr),
    \end{align}
    where \(C_{\mathsf{C\acute{e}a}}\) is independent of \(\kappa,\rF,\varrho_0,\) and \(h.\) 
\end{theorem}


\begin{proof}
    Let \(\varphi\in\Psi\) solve the system~\eqref{eq:weak.hat.phi} for \(\hat\bu=\bu\in\bV,\) and \(\varphi_h\in\Psi_h\) solve  \eqref{eq:weak.hat.phi.disc} for \(\hat\bu_h=\bu_h\in\bV_h.\) Subtracting \eqref{eq:weak.fixed.u.hat.disc} from \eqref{eq:weak.fixed.u.hat}, with \(\hat\varphi=\varphi\) and \(\hat\varphi_h=\varphi_h,\) we obtain the following system of equations 
    % \rrb{why do we get $G_h$?}
    \begin{subequations}
        {\begin{align}
         a_1(\bu-\bu_h,\bv_h) + d_1(\bu; \bu,\bv_h) - d_1(\bu_h; \bu_h,\bv_h) +  b_{1,h}(\bv_h,p-p_h) & = F^{\varphi}(\bv_h)-F^{\varphi_h}(\bv_h) ,\label{eq:EquationCea}\\
    b_{1,h}(\bu-\bu_h,q_h) & = 0,
    \end{align}}
    \end{subequations}
    for all \(\bv_h\in\bV_h\) and \(q_h\in\rQ_h.\)
{  
{Let $\bu_h^*\in\bV_h$ such that $\bu_h-\bu_h^*\in\{\bv_h\in\bV_h:b_{1,h}(\bv_h,q_h)=0,\, \forall q_h\in\rQ_h\}.$}
    From the uniform strong monotonicity of \(a_1(\bullet,\bullet)+d_1(\bullet;\bullet,\bullet)\) on the kernel of \(b_{1,h}\) (see \cite[Equation~3.8]{dhpr_sinum26}),  we deduce 
    \begin{align}
        \alpha_{\mathsf{DF}}\|\bu_h-\bu_h^*\|^2_{\bV}\leq a_1(\bu_h,\bu_h-\bu_h^*)-a_1(\bu_h^*,\bu_h-\bu_h^*)+d_1(\bu_h;\bu_h,\bu_h-\bu_h^*)-d_1(\bu_h^*;\bu_h^*,\bu_h-\bu_h^*).
    \end{align}
    Further, adding and subtracting \(a_1(\bu,\bu_h-\bu_h^*)+d_1(\bu;\bu,\bu_h-\bu_h^*)\) to the above inequality, using \eqref{eq:EquationCea}, \eqref{eq:lip.bound.f}, and the local Lipschitz continuity of \(a_1(\bullet,\bullet)+d_1(\bullet;\bullet,\bullet)\) (see \cite[Equation~3.7]{dhpr_sinum26}), {and lastly using the triangle inequality} the following bound is obtained
\begin{align*}\alpha_{\mathsf{DF}}\|\bu_h-\bu_h^*\|_{\bV}&\leq\|p-p_h\|_{\widehat{\rQ}(\cT_h)}+\|F^{{\varphi}}-F^{{\varphi}_h}\|_{\bV'}+(1+\|\bu\|_{\bV}+\|\bu_h^*\|_{\bV})\|\bu-\bu_h^*\|_{\bV}\\
    &\leq \|p-p_h\|_{\widehat{\rQ}(\cT_h)} + (1+C_{\mathsf{emb}}\alpha^{-\frac12})\alpha_1\|{\varphi}-{\varphi}_h\|_{\widehat{\Psi}(\cT_h)}+(1+\|\bu\|_{\bV}+\|\bu_h^*\|_{\bV})\|\bu-\bu_h^*\|_{\bV}\\
    &\leq{\|p-p_h\|_{\widehat{\rQ}(\cT_h)} + (1+C_{\mathsf{emb}}\alpha^{-\frac12})\alpha_1\|{\varphi}-{\varphi}_h\|_{\widehat{\Psi}(\cT_h)}+(1+2\|\bu\|_{\bV})\|\bu-\bu_h^*\|_{\bV}+\|\bu-\bu_h^*\|_{\bV}^2}.
    \end{align*}
    Finally, using triangle inequality and the fact that \({\|\bu-\bu_h^*\|_{\bV}}\leq(1+\frac1{\beta_{\mathsf{d}}})\,\mathrm{dist}(\bu,\bV_h)\) (cf. \cite[Theorem~2.6]{gatica2014simple}), where \(\beta_{\mathsf{d}}\) is the continuous inf-sup constant corresponding to the Darcy--Forchheimer model, we arrive at the following error estimate for the velocity
\begin{align}\label{eq:error.est.u-uh}\|\bu-\bu_h\|_{\bV}&\leq\frac{1}{\alpha_{\mathsf{DF}}}\mathrm{dist}(p,\rQ_h)+(1+C_{\mathsf{emb}}\alpha^{-\frac12})\frac{\alpha_1}{\alpha_{\mathsf{DF}}}\|{\varphi}-{\varphi_h}\|_{\widehat{\Psi}(\cT_h)}\nonumber\\ \nonumber&+[1+\frac1{\alpha_{\mathsf{DF}}}(1+\frac1{\beta_{\mathsf{d}}})(1+\|\bu\|_{\bV})]\mathrm{dist}(\bu,\bV_h)+{\frac1{\alpha_{\mathsf{DF}}}(1+\frac1{\beta_{\mathsf{d}}})\mathrm{dist}(\bu,\bV_h)}\\
    &\leq C_G\bigl(\mathrm{dist}(\bu,\bV_h)+{\mathrm{dist}(\bu,\bV_h)^2}+\mathrm{dist}(p,\rQ_h)\bigr)+(1+C_{\mathsf{emb}}\alpha^{-\frac12})\frac{\alpha_1}{\alpha_{\mathsf{DF}}}\|\varphi-{\varphi}_h\|_{\widehat{\Psi}_h},\end{align}
    where {\(C_G:=
    \max\bigl({1+\frac1{\alpha_{\mathsf{DF}}}(1+\frac1{\beta_{\mathsf{d}}})(1+\|\bu\|_{\bV})},{1+\frac1{\alpha_{\mathsf{DF}}}(1+\frac1{\beta_{\mathsf{d}}})},\frac1{\alpha_{\mathsf{DF}}}\bigr).\)}
    Next, for the pressure estimate, let us consider an arbitrary element \(p_h^*\in\rQ_h,\) use \eqref{eq:EquationCea}, and the discrete inf-sup stability of the pressure, and lastly, Lipschitz-continuity of \(a_1(\bullet,\bullet)+d_1(\bullet;\bullet,\bullet)\) to conclude the following
    \begin{align*}
        \beta_{\mathsf{d}}\|p_h-p_h^*\|_{\widehat{\rQ}(\cT_h)}&\leq\sup_{\bzero\neq\bv_h\in\bV_h}\frac{b_{1,h}(\bv_h,p_h-p_h^*)}{\|\bv_h\|_{\bV}}\\
        &\leq \sup_{\bzero\neq\bv_h\in\bV_h}\frac{b_{1,h}(\bv_h,p-p_h^*)+b_{1,h}(\bv_h,p_h-p)} {\|\bv_h\|_{\bV}}\leq \sup_{\bzero\neq\bv_h\in\bV_h}\frac{b_{1,h}(\bv_h,p-p_h^*)}{\|\bv_h\|_{\bV}}\\
        &+\sup_{\bzero\neq\bv_h\in\bV_h}\frac{F^{{\varphi}}(\bv_h)-F^{{\varphi}_h}(\bv_h)-a_1(\bu_h-\bu_h^*,\bv_h)-d_1(\bu_h;\bu_h,\bv_h)+d_1(\bu_h^*;\bu_h^*,\bv_h)}{\|\bv_h\|_{\bV}}\\
        &\leq\|p-p_h^*\|_{\widehat{\rQ}(\cT_h)}+(1+C_{\mathsf{emb}}\alpha^{-\frac12})\alpha_1\|\varphi-\varphi_h\|_{\widehat{\Psi}(\cT_h)}+(1+2\|\bu\|_{\bV})\|\bu-\bu_h^*\|_{\bV} +{\|\bu-\bu_h^*\|_{\bV}^2}.
    \end{align*}
    Finally, a use of the triangle inequality yields the error estimate for pressure 
    % \rrb{fixed parentheses. please check}
\begin{align}\label{eq:error.est.p-ph}
        \|p-p_h\|_{\widehat{\rQ}(\cT_h)}&\leq(1+\frac1{\beta_{\mathsf{d}}}+C_G)\,\mathrm{dist}(p,\rQ_h)+(1+C_{\mathsf{emb}}\alpha^{-\frac12})(\frac{\alpha_1}{\alpha_{\mathsf{DF}}}+\frac{\alpha_1}{\beta_{\mathsf{d}}})\|\varphi-\varphi_h\|_{\widehat{\Psi}(\cT_h)}\nonumber\\&\qquad\qquad+C_G(\mathrm{dist}(\bu,\bV_h)+\mathrm{dist}(\bu,\bV_h)^2)\nonumber\\
        &\leq \tilde C_G \bigl(\mathrm{dist}(\bu,\bV_h)+\mathrm{dist}(\bu,\bV_h)^2+\mathrm{dist}(p,\rQ_h)+\|\varphi-\varphi_h\|_{\widehat{\Psi}(\cT_h)}\bigr),
    \end{align}
    where {\(\tilde C_G:=\max\bigl((1+\frac1{\beta_{\mathsf{d}}}+C_G), (1+C_{\mathsf{emb}}\alpha^{-\frac12})(\frac{\alpha_1}{\alpha_{\mathsf{DF}}}+\frac{\alpha_1}{\beta_{\mathsf{d}}})
    %, C_G(1+\|\bu\|_{\bV}+\|\bu_h^*\|_{\bV})
    \bigr).\)}
} 

    Next, subtracting \eqref{eq:weak.hat.phi.disc} from \eqref{eq:weak.hat.phi}, it is not difficult to observe that \((\bzeta^*-\bzeta_h,\varphi^*-\varphi_h)\in\bW\times \widehat{\Psi}(\cT_h)\) solves the following system of equations for any arbitrary element \((\bzeta_h^*,\varphi_h^*)\in\bW_h\times \Psi_h\) \((\bu-\bu_h) \) (by substituting, in particular, \(\hat\bu=\bu,\hat\bu_h=\bu_h\) in \eqref{eq:weak.fixed.u.hat}, 
    \eqref{eq:weak.fixed.u.hat.disc}, repectively). Hence,
    \begin{align*}
        a_2(\bzeta^*_h-\bzeta_h,\bw_h) + b_{2,h}(\bw_h,\varphi^*_h-\varphi_h) + d_2^{\bu_h}(\bw_h, \varphi^*_h-\varphi_h) & = \mathcal{H}_{1,h}(\bw_h) \quad \forall \bw_h \in \bW_h, \\
    b_{2,h}(\bzeta_h^*-\bzeta_h, \psi_h) - c_{2,h}(\varphi^*_h-\varphi_h,\psi_h) & = \mathcal{H}_{2,h}(\psi_h) \quad \forall \psi_h \in \Psi_h,
    \end{align*}
    where \begin{align*}
\mathcal{H}_{1,h}(\bw_h)&:=a_2(\bzeta_h^*-\bzeta,\bw_h) + b_{2,h}(\bw_h,\varphi_h^*-\varphi) + d_2^{\bu_h}(\bw_h, \varphi_h^*-\varphi)-d_2^{\bu-\bu_h}(\bw_h,\varphi),\\
\text{and}\quad \mathcal{H}_{2,h}(\psi_h)&:=b_{2,h}(\bzeta_h^*-\bzeta,\psi_h)+c_{2,h}(\varphi-\varphi_h^*,\psi_h).
\end{align*} 
    {Finally, invoking \eqref{eq:ARD.disc}, together with the discrete stability estimate for \((\bzeta_h-\bzeta_h^*,\varphi_h-\varphi_h^*),\) and applying the triangle inequality, we obtain} 
    % \rrb{phrase it better}
   \begin{align}\label{eq:stab.ard.cea}\|(\bzeta-\bzeta_h,\varphi-\varphi_h)\|_{\bW\times\widehat{\Psi}(\cT_h)}\leq C_{\mathsf{C\acute{e}a}}^{\mathsf{ADR}}&\bigl(\mathrm{dist}(\bzeta,\bW_h)+\mathrm{dist}(\varphi,\Psi_h)\bigr)\nonumber\\&+\rF^{-\frac13}(1+C_{\mathsf{emb}}\alpha^{-\frac12})C_c^{\mathsf{ADR}}C_{\mathsf{d}}^{\mathsf{ADR}}\|m\|_{0,\frac65,\Omega}\|\bu-\bu_h\|_{\bV},\end{align}
   where 
   % \rrb{fix lines}
\begin{align}\label{eq:Cea.ARD}
C_{\mathsf{C\acute{e}a}}^{\mathsf{ADR}}&:=C_{\mathsf{d}}^{\mathsf{ADR}}\bigl(3+C_{\mathsf{emb}}(1+C_{\mathsf{emb}}\alpha^{-\frac12})^2+\bigl(\rF^{-\frac13}(1+C_{\mathsf{emb}}\alpha^{-\frac12})\nonumber\\& \qquad\qquad \times C_{\mathsf{d},\bu_h}^{\mathsf{DF}}(\|\bf(\bzero)\|_{0,\Omega}+\alpha_1\alpha^{-\frac12}(1+C_{\mathsf{emb}}))C_{\mathsf{d}}^{\mathsf{ADR}\|m\|_{0,\frac65,\Omega}}\bigr)\bigr).
\end{align}
Next, invoking \eqref{eq:stab.ard.cea} in \eqref{eq:error.est.u-uh}, we arrive at the following error estimate for velocity
\begin{align}\label{eq:cea.u}
       \|\bu-&\bu_h\|_{\bV}\leq C_G\bigl(\mathrm{dist}(\bu,\bV_h)+{\mathrm{dist}(\bu,\bV_h)^2}+\mathrm{dist}(p,\rQ_h)\bigr)\nonumber\\&+C_{\mathsf{C\acute{e}a}}^{\mathsf{ADR}}(1+C_{\mathsf{emb}}\alpha^{-\frac12}){\frac{\alpha_1}{\alpha_{\mathsf{DF}}}}\bigl(\mathrm{dist}(\bzeta,\bW_h)+\mathrm{dist}(\varphi,\Psi_h)\bigr)\nonumber\\&+{(1+C_{\mathsf{emb}}\alpha^{-\frac12})}{\frac{\alpha_1}{\alpha_{\mathsf{DF}}}}\rF^{-\frac13}(1+C_{\mathsf{emb}}\alpha^{-\frac12})C_c^{\mathsf{ADR}}C_{\mathsf{d}}^{\mathsf{ADR}}\|m\|_{0,\frac65,\Omega}\|\bu-\bu_h\|_{\bV},  
   \end{align}
   and hence,   \begin{align}\label{eq:final.err.est.for.u}
       \|\bu-\bu_h\|_{\bV}\leq C^{\mathsf{C\acute{e}a}}_{\bu}\bigl(\mathrm{dist}(\bu,\bV_h)+{\mathrm{dist}(\bu,\bV_h)^2}+\mathrm{dist}(p,\rQ_h)+\mathrm{dist}(\bzeta,\bW_h)+\mathrm{dist}(\varphi,\Psi_h)\bigr),
   \end{align}
   where \(C^{\mathsf{C\acute{e}a}}_{\bu}:={\max\bigl(C_G,{\frac{\alpha_1}{\alpha_{\mathsf{DF}}}}{(1+C_{\mathsf{emb}}\alpha^{-\frac12})}C_{\mathsf{C\acute{e}a}}^{\mathsf{ADR}}\bigr)}(1-{\rF^{-\frac13}(1+C_{\mathsf{emb}}\alpha^{-\frac12})C_c^{\mathsf{ADR}}C_{\mathsf{d}}^{\mathsf{ADR}}}\|m\|_{0,\frac65,\Omega})^{-1}.\) Remember that, from the hypothesis of Theorem~\ref{th:Cea},  \((1-{\rF^{-\frac13}(1+C_{\mathsf{emb}}\alpha^{-\frac12})C_c^{\mathsf{ADR}}C_{\mathsf{d}}^{\mathsf{ADR}}}\|m\|_{0,\frac65,\Omega})> 0.\)
   
   {Invoking the bound of \eqref{eq:final.err.est.for.u} and \eqref{eq:stab.ard.cea} in \eqref{eq:error.est.p-ph} we obtain the error estimate for pressure as follows\begin{align}\label{est:final.error.est.for.p}
       \|p-p_h\|_{\widehat\rQ(\cT_h)}&\leq
C^{\mathsf{C\acute{e}a}}_{p}\bigl(\mathrm{dist}(\bu,\bV_h)+{\mathrm{dist}(\bu,\bV_h)^2}+\mathrm{dist}(p,\rQ_h)+\mathrm{dist}(\bzeta,\bW_h)+\mathrm{dist}(\varphi,\Psi_h)\bigr),
   \end{align}
   where $C^{\mathsf{C\acute{e}a}}_{p}:=\max\bigl(\tilde{C}_G,C_{C\acute{e}a}^{\mathsf{ADR}},C_{\mathsf d}^{\mathsf{ADR}}C_{\bu}^{C\acute{e}a}\|m\|_{0,\frac65,\Omega}\bigr).$}
 Finally, using \eqref{eq:final.err.est.for.u} and \eqref{est:final.error.est.for.p} in \eqref{eq:stab.ard.cea}, we obtain
  \begin{align*}%\label{eq:Cea}
        \|\bu-&\bu_h\|_{\bV}+\|p-p_h\|_{\widehat{\rQ}(\cT_h)}+\|\bzeta-\bzeta_h\|_{\bW}+\|\varphi-\varphi_h\|_{\widehat{\Psi}(\cT_h)}\nonumber\\
        & \leq C_{\mathsf{C\acute{e}a}}\bigl(\mathrm{dist}(\bu,\bV_h)+{\mathrm{dist}(\bu,\bV_h)^2}+\mathrm{dist}(p,\rQ_h)+\mathrm{dist}(\bzeta,\bW_h)+\mathrm{dist}(\varphi,\Psi_h) \bigr),
    \end{align*}
    where  {\(C_{\mathsf{C\acute{e}a}}:=\max\bigl({\rF^{-\frac13}(1+C_{\mathsf{emb}}\alpha^{-\frac12})C_c^{\mathsf{ADR}}C_{\mathsf{d}}^{\mathsf{ADR}}}\|m\|_{0,\frac65,\Omega}C_{\bu}^{C\acute{e}a}, C_{p}^{C\acute{e}a}, C_{C\acute{e}a}^{\mathsf{ADR}}\bigr),\)}
    is independent of the model parameters \(\kappa, \varrho_0,\) and the discretisation parameter~\(h.\) This concludes the rest of the proof.
\end{proof}
\begin{remark}\label{rem:pressure,flux.robust}
    {Let us remark that, similarly as in \cite[Corollary~4.5]{dhpr_sinum26}, it is possible to derive pressure-robust and concentration-robust error estimates for velocity and flux (because of the discrete velocity and flux kernels being subspaces of the continuous ones) as follows
    \begin{subequations}
        \begin{align}
            \|\bu-\bu_h\|_{\bV}\leq C_{\mathsf{V}}(\mathrm{dist}(\bu,\bV_h)+{\mathrm{dist}(\bu,\bV_h)^2}),\quad \text{and}\quad
            \|\bzeta-\bzeta_h\|_{\bW}\leq C_{\mathsf{F}}\,\mathrm{dist}(\bzeta,\bW_h),
        \end{align}
    \end{subequations}
    where the constants \(C_{\mathsf{V}}\) and \(C_{\mathsf{F}}\) are independent of \(\kappa,\varrho_0,\) and \(h\).}
\end{remark}

\noindent
In order to provide a theoretical convergence rate, we recall some useful approximation properties. 
Let $t \in (1,\infty)$ and denote by  \(\Pi_h:\bW^{l,t}(\Omega)\cap\bH(\vdiv_{1,t},\Omega)\to\mathbb{RT}_k(\cT_h)\)  the Fortin operator and by \(P_h:\rW^{l,t}(K)\to \mathbb{P}_k(K)\),  the local orthogonal \(\rL^2\) projection. Then, as a standard consequence of the Bramble--Hilbert lemma, optimal approximation, and inverse inequalities on polynomial spaces, we can assert that  
\begin{subequations}
     \begin{align}
\label{eq:approx-v}        
& \star {\|\bv- \Pi_h \bv\|_{m,t,\Omega}\leq C_1 h^{l-m}|\bv|_{l,t,\Omega}\quad \forall \bv\in \mathbf{W}^{l,t}(\Omega), \ 1\leq l \leq k+1,}\\
  \label{eq:approx-div}  & \star { \|\vdiv(\bv-\Pi_h\bv)\|_{m,t,\Omega}\leq C_2h^{l-m}|\vdiv\bv|_{l,t,\Omega} \quad \forall \bv \in \mathbf{W}^{1,t}(\Omega),\ \text{with}}\\
    \nonumber           & \qquad \qquad \qquad \qquad \qquad \qquad \qquad \qquad \qquad  {\vdiv(\bv)\in \rW^{1,t}(\Omega)\ \text{and} \ 0\leq l \leq k+1,}\\
  \label{eq:approx-p}  & \star{\|q-P_hq\|_{m,t,\Omega}\leq C_2 h^{l-m}|p|_{l,t,\Omega}\quad \forall q\in \rW^{l,t}(\Omega),\ 0\leq l \leq k+1,} \\
    \label{eq:approx-p-K}  & \star {\|\nabla(q-P_hq)\|_{0,t,{K}}\leq C_3 h_K^{l}|p|_{l+1,t,K}\quad \forall q\in \rW^{l+1,t}({K}),\  K\in\mathcal{T}_h,\ l \leq k,} 
     \end{align}
     \end{subequations}
     {for $ 0\leq m\leq l$, where $C_1,C_2>0$ are  independent of \(h\) and $C_3>0$ is independent of $h_K$ 
(see \cite[Sections 4.1 and 4.5]{gatica2022p} for the first three, and \cite{ern2021finite} for the fourth  one).}

To conclude this section, the following theorem provides the theoretical rate of convergence of the mixed finite element scheme \eqref{eq:weak}, under suitable regularity on the exact solution. 
 \begin{theorem}[Convergence rates]\label{th:convergence}
        Let $(\bu,p,\bzeta,\varphi) \in \bV \times \rQ \times \bW\times \Psi$ be the solution of the continuous problem \eqref{eq:weak}  with the following regularity conditions
{\begin{gather*}\bu\in \kappa^{-\frac12}{\bH^{l}}(\Omega)\cap{\rF^{\frac13}\mathbf{W}^{l,3}(\Omega)},  \quad \vdiv\bu\in {\rH^l}(\Omega), \quad \text{and} \quad {p=p_1+p_2+p_3}, \\
\text{with} \quad {p_1\in \rH^{l}(\Omega), \quad p_2\in \kappa^{\frac12}\rH^{l+1}(\Omega), \quad \text{and}\quad  p_3 \in \rF^{-\frac13}\rW^{l+1,\frac32}(\Omega)},\\
\bzeta\in \alpha^{-\frac12}\bH^{l}(\Omega)\cap\bH(\vdiv_{\frac65},\Omega), \quad \vdiv\bzeta\in {\rW}^{l,\frac65}(\Omega), \quad \text{and} \quad \varphi=\varphi_1+\varphi_2,\\
\text{with}\quad\varphi_1\in \alpha^{\frac12}\rH^{l+1}(\Omega),\quad \varphi_2\in\rW^{l+1,6}(\Omega),
\end{gather*}
where \(1\leq l\leq k+1\).  }
Then, there exist positive constants \(C^i_\sharp,\,i=1,\cdots4\), independent of \(h, \kappa, \rF, \alpha\), and \(\varrho_0,\) such that 
\begin{align*}
 \|\bu-\bu_h\|_{\bV}+\|p-p_h\|_{\widehat{\rQ}(\cT_h)}&+\|\bzeta-\bzeta_h\|_{\bW}+\|\varphi-\varphi_h\|_{\widehat{\Psi}(\cT_h)}\leq C^1_\sharp h^{l}\bigl[\kappa^{-\frac12}\|\bu\|_{l,\Omega}+\|\vdiv\bu\|_{l,\Omega}\\&+\rF^{\frac13}\|\bu\|_{l,3,\Omega}\bigr]+{(C^1_\sharp)^2 h^{2l}\bigl[\kappa^{-\frac12}\|\bu\|_{l,\Omega}+\|\vdiv\bu\|_{l,\Omega}+\rF^{\frac13}\|\bu\|_{l,3,\Omega}\bigr]^2}\\& +{C^2_\sharp h^{l}\bigl[ |p_1|_{l,\Omega}+\kappa^{\frac12}|p_2|_{l+1,\Omega}+\rF^{-\frac13}|p_3|_{l+1,\frac32,\Omega}\bigr]}\\&+{C^3_\sharp h^{l}\{\alpha^{-\frac12}\|\bzeta\|_{l,\Omega}+\|\vdiv\bzeta\|_{l,\frac65,\Omega}\}} +{C^4_\sharp h^{l}\{\alpha^{\frac12}|\varphi_1|_{l+1,\Omega}+|\varphi_2|_{l+1,6,\Omega}\}}.
\end{align*}
    \end{theorem}
    \begin{proof}  
        {Using the approximation properties \eqref{eq:approx-v}, \eqref{eq:approx-div}, \eqref{eq:approx-p},  and \eqref{eq:approx-p-K} in \eqref{eq:final.err.est.for.u} and \eqref{est:final.error.est.for.p}, we obtain the following convergence rates for velocity and pressure}
\begin{align*} 
 \|\bu - \bu_h\|_{\bV}   \leq C_{\mathsf{V}}\, \mathrm{dist}(\bu,\bV_h) &\leq C_{\mathsf{V}}\max(2C_1,C_2) h^l (\kappa^{-\frac12}\|\bu\|_{l,\Omega} + \|\vdiv\bu\|_{l,\Omega} + \rF^{\frac13}\|\bu\|_{l,3,\Omega})\\&+{\bigl(\max(2C_1,C_2) h^l (\kappa^{-\frac12}\|\bu\|_{l,\Omega} + \|\vdiv\bu\|_{l,\Omega} + \rF^{\frac13}\|\bu\|_{l,3,\Omega})\bigr)^2}, \,\text{and}\\
 \|p -  p_h\|_{\widehat{\rQ}(\cT_h)}  \leq C_p^{C\acute{e}a}(\| p_1 - P_hp_1\|_{0,\Omega}& + \kappa^{\frac12}|p_2-P_hp_2|_{1,\cT_h} + \rF^{-\frac13}|p_3-P_hp_3|_{1,\frac32,\cT_h}) \\
&\leq  C_p^{C\acute{e}a}\max\bigl(C_2,C_3,C_4\bigr)h^{l}\bigl(|p_1|_{l,\Omega} +  \kappa^{\frac12}|p_2|_{l+1,\Omega} +  \rF^{-\frac13}|p_3|_{l+1,\frac32,\Omega} \bigr),
\end{align*} 
where $C_4$ is same as in \cite[Theorem~4.6]{dhpr_sinum26}.
{Next, for the error estimate for flux, at first take \(m=0\) and \(t=0\) in \eqref{eq:approx-v}, and then, \(m=0\) and \(t=\frac65\) in \eqref{eq:approx-div}, which lead to
\begin{align*}
    \|\bzeta-\Pi_h\bzeta\|_{0,\Omega}\leq C_1C_{\mathsf{V}}h^l\,|\bzeta|_{l,\Omega},\quad\text{and}\quad 
    \|\vdiv(\bzeta-\Pi_h\bzeta)\|_{0,\frac65,\Omega} \leq C_2C_{\mathsf{V}}h^l|\vdiv\bu|_{l,\frac65,\Omega}.
\end{align*}}
{
In turn, owing to Remark~\ref{rem:pressure,flux.robust}, we derive the error bound on the flux as follows
\begin{align*}
    \|\bzeta-\bzeta_h\|_{\bW}\leq C_{\mathsf{F}}\,\mathrm{dist}(\bzeta,\bW_h) \leq C_{\mathsf{F}}\max(C_1,C_2)h^l(\alpha^{-\frac12}\|\bzeta\|_{l,\Omega}+\|\vdiv\bzeta\|_{0,\frac65,\Omega}).
\end{align*}
Lastly, for the concentration we employ the  decomposition \(\varphi=\varphi_1+\varphi_2,\) and the properties of the sum norm  \eqref{norm:Psi.hat}, and  apply  \eqref{eq:approx-p} with $m=0$ and  $t=2$, and \eqref{eq:approx-p-K} with $t=0$ and \(t=6,\) summed over $K\in \cT_h$. We combine that with the following trace inequality to control the jump term in the broken norms
\[h_F^{1-t} \|[\![\varphi_j - P_h \varphi_j]\!] \|^t_{0,t,F} =h_F^{1-t} \| [\![P_h \varphi_j]\!]\|^t_{0,t,F} \leq C_5 h^{lt}|\varphi_j|^t_{l+1,r,\omega_F}, \]
where $\omega_F$ stands for the usual face patch. Therefore, we readily obtain the following bound 
\begin{align*} 
\|\varphi - P_h \varphi\|_{\widehat{\Psi}(\cT_h)} & \leq C_{C\acute{e}a}( \| \varphi_1 - P_h\varphi_1\|_{0,6,\Omega} + \alpha^{\frac12}|\varphi_2-P_h\varphi_2|_{1,\cT_h})\\&\leq C_{C\acute{e}a} \max\bigl(C_2,C_3,C_5\bigr) h^{l}\bigl(|\varphi_1|_{l,6,\Omega} +  \alpha^{\frac12}|\varphi_2|_{l+1,\Omega}\bigr),
\end{align*}}
where $C_5>0$ is independent of all the parameters. This completes the proof. 
    \end{proof}

 In the next section we define a variable operator preconditioner for the Newton linearisation of  \eqref{eq:weak}. 
%%%%%%%%%%%%%%%%%%%%%%%%%%%%%%%%%
\section{Operator preconditioning}\label{sec:precond}
This section discusses preconditioners for the monolithic coupled system arising from the weighted norms introduced in the previous sections which is robust with respect to all the model parameters expect for coefficient of diffusion.

First, for \(0<s<\infty\), we denote by \(\mathcal{J}_s:\rL^s(\Omega)\to\bigl(\rL^s(\Omega)\bigr)'\)  the canonical duality map defined as 
\begin{equation}
\langle{\mathcal{J}}_s(\theta_1),\theta_2\rangle:=\int_{\Omega}|\theta_1|^{s-2}\theta_1\theta_2,\quad\forall\theta_1,\theta_2\in\rL^s(\Omega).\end{equation}
Similarly, mimicking the operator notation $-\nabla \vdiv$ that induces the div-div bilinear form, we define $-\nabla \mathcal{J}_{s}(\vdiv)$ as}
\begin{equation}\label{eq:duality2} {\langle-\nabla \mathcal{J}_{s}(\vdiv)(\bzeta_1),\bzeta_2\rangle:={\langle} \mathcal{J}_{s}(\vdiv\bzeta_1{),\vdiv\bzeta_2\rangle},\quad \forall \bzeta_1,\bzeta_2 \ \text{such that} \ \vdiv\bzeta_i \in \rL^s(\Omega).}\end{equation}

Denote by \(\mathcal{H}_1(\bu)\) the G\^ateaux derivative of the nonlinear operator $d_1(\bullet;\bullet,\bullet)$, i.e., \(|\bu|\bu\) at the state \(\bu\) in the direction of \(\delta\bu\), and \(\mathcal{H}_2(\ast)\) by the linearisation of \(d_2^{\ast}(\bullet,\bullet)\) around \(\ast\). The linearised system corresponding to \eqref{eq:strong} takes the following form
\begin{equation}\label{eq:strong.linearised}
%\left.
\begin{aligned}
[\kappa^{-1}+\rF\mathcal{H}_1(\hat\bu)]\bu + \nabla p & = \bf(\hat\varphi) \quad \text{in $\Omega$},\\
    \vdiv \bu & = 0 \quad \text{in $\Omega$},\\
    \alpha^{-1}\bzeta- \nabla \varphi + { \alpha^{-1}} \mathcal{H}_2(\hat\bu) \varphi &=\bzero \quad \text{in $\Omega$},\\
    \vdiv \bzeta - \varrho_0 \varphi & = -m \quad \text{in $\Omega$},
\end{aligned}
% \right\}
 \end{equation}
with the same boundary conditions as in \eqref{eq:strong}. 
In addition, consider the Newton--Raphson linearisation, which starting from an initial guess $\bu^0,p^0,\bzeta^0,\varphi^0$, and iterating on  $m = 1, 2, \ldots$ until convergence, seeks velocity, pressure, flux, and concentration increments $(\delta \bu, \delta p, \delta\bzeta, \delta\varphi)$ in the space \begin{align}\label{eq:linearised-spaces}\underline{\bV}\times \underline{\rQ}\times \underline{\bW}\times\underline{\Psi} &:=\bigl[\kappa^{-\frac12}\bL^2(\Omega)\cap\bH_{\Gamma_{\bu}}(\vdiv,\Omega)\cap\{\rF\mathcal{H}_1(\hat{\bu})\}^\frac13\bL^3(\Omega)\bigr]\nonumber\\ &\quad \times \bigl[\rL^2(\Omega)+\kappa^\frac12\rH_{\Gamma_p}^1(\Omega)+\{\rF{\mathcal{H}_1(\hat{\bu})}\}^{-\frac13}\rW_{\Gamma_p}^{1,\frac32}(\Omega)\bigr]\nonumber\\&\quad\times \bigl[(\mathcal{H}_2^{-1}(\hat{\bu})\alpha)^{-\frac12}\bL^2\cap\bH_{\Gamma_p}(\vdiv_{\frac65},\Omega)\bigr]\times \bigl[(\mathcal{H}^{-1}_2(\hat{\bu})\alpha)^{\frac12}\rH^1(\Omega)+\rL^6(\Omega)\bigr]
,\end{align} 
such that the following holds in $\underline{\bV}'\times \underline{\rQ}'\times\underline{\bW}'\times \underline{\Psi}'$
{\begin{equation}\label{eq:newton} 
\begin{pmatrix}[\kappa^{-1}+\rF 
\mathcal{H}({\bu}^m)]\mathcal{I} & \nabla&0&0 \\
    -\vdiv  &0 &0&0\\\alpha^{-1}\varphi^m&0&\alpha^{-1}\mathcal{I}&-\nabla+\alpha^{-1}\bu^m\\0&0&\vdiv&-\varrho_0\mathcal{I}
    \end{pmatrix}\begin{pmatrix}\delta \bu \\ \delta p\\\delta\bzeta\\\delta\varphi\end{pmatrix}= \begin{pmatrix}
      F - \{\kappa^{-1}+|\bu^m|\}({\bu}^m) - \nabla{p}^m \\ G +\vdiv{\bu}^m  \\-\alpha^{-1}\bzeta^m+\nabla\varphi^m-\alpha^{-1}\bu^m\varphi^m\\H-\vdiv\bzeta^m-\varrho_0\varphi^m
    \end{pmatrix}. \end{equation}} 
    
Here, $\mathcal{I}$ denotes the identity operator  at the continuous level, but for the discrete case it stands for the corresponding mass matrix (either for velocity or pressure or flux or concentration). Further, we note that, in weak form, the action of the operator $\rF\mathcal{H}({\bu}^m)\mathcal{I}:\underline{\bV}\to \underline{\bV}'$ reads
\[ \langle \rF \mathcal{H}({\bu}^m)(\delta \bu), \bv \rangle := \int_\Omega \rF |\bu^m|\delta\bu \cdot \bv + \int_\Omega \rF |\bu^m|^{-1}(\bu^m \cdot \delta\bu) (\bu^m \cdot \bv), \]
and \[\langle\mathcal{H}_2(\bu^m)\btau,\varphi\rangle:=\int_{\Omega}(\bu^m\cdot\btau)\varphi.\]
Then, update the solution as \((\bu^{m+1},p^{m+1},\bzeta^{m+1},\varphi^{m+1}) = (\bu^m + \delta \bu, p^m + \delta p,\bzeta^m+\delta\bzeta,\varphi^m+\delta\varphi)\). 
\noindent
Next, we follow the approach recently advanced in \cite[Section~3]{cai2026parameter} to propose a robust block-diagonal preconditioner for \eqref{eq:newton}. 
 % \rd{We need to remove a few things. Maybe I'll comment the blue part, and re-write/ keep the things we need currently.}
\begin{comment}
 \cblue
{We construct a block-diagonal preconditioner for the Darcy--Forchheimer velocity and pressure blocks following 
\cite{dhpr_sinum26}. On the other hand, and differently from \cite{cai2026parameter} and \cite{dhpr_sinum26}, we readily note that the sub-blocks associated with the ARD equations consist of a perturbed, \emph{non-symmetric} saddle-point of the form 
\begin{equation}\label{abstract-block} \begin{pmatrix} \alpha^{-1}\bbA & \bbB^* + \alpha^{-1}\bbD \\ \bbB & - \varrho_0 \bbC\end{pmatrix},\end{equation}
where we are abusing of notation to emphasise the parameter dependence of the different blocks. 
Let us first concentrate on the discrete case since it simplifies the discussion as $\bu_h^m \in \bL^{\infty}(\Omega)$ and so the Banach structure is not required. In this case, $\bbA:\bH(\vdiv,\Omega)\to \bH(\vdiv,\Omega)'$ denotes the  $\bL^2(\Omega)$ mass map in $\bH(\vdiv,\Omega)$, $\bbB:\bH(\vdiv,\Omega)\to \rL^2(\Omega)'$ is the divergence operator, $\bbD$ is the \textit{mixed} advection coupling operator ($\bbD\varphi \approx \bu^m\varphi$), and $
\bbC:\rL^2(\Omega) \to \rL^2(\Omega)'$ is the scalar mass map in $\rL^2(\Omega)$. 
Therefore a typical Riesz map preconditioner, also based on a Schur complement approximation, is motivated by \cite[Sect. 3.1]{benzi2005numerical} (see also \cite[Section 1]{baerland2020observation} for the symmetric, unperturbed case). The guiding principle is to scale the flux space with $\alpha^{-1}$, the concentration space with $\varrho_0$, and ensure that the Schur complement is spectrally equivalent to a properly weighted  operator. In this case, a block-diagonal preconditioner 
%and a (better) block-triangular version 
reads as follows 
\begin{equation}\label{abstract-prec} 
 \mathrm{diag}(\cR,\cS) = \begin{pmatrix}(\alpha^{-1}\bbA + \alpha^{-1} \bbB^*\bbB)^{-1} & 0 \\ 0 & (\varrho_0 \bbC + \alpha \bbB\bbA^{-1}\bbB^* + \bbB\bbA^{-1}\bbD)^{-1}
\end{pmatrix},
\end{equation}
% \begin{equation}
% \cgray{\mathrm{diag}(\cR,\cS) = 
%  \begin{pmatrix} \alpha \bbA^{-1} & -\alpha\bbA^{-1}(\bbB^* + \alpha^{-1}\bbD^*)(\varrho_0\bbC+ \bbB\bbA^{-1}\bbB^*)^{-1} \\ 0 & (\varrho_0 \bbC + {\alpha}\bbB\bbA^{-1}\bbB^* + {\alpha}\bbB\bbA^{-1}\bbD^*)^{-1}\end{pmatrix}.}
% \end{equation}
(and note that it is also possible to derive block-triangular versions).  For the the advection term the preconditioner part $\bbB\bbA^{-1}\bbD$ is not symmetric and it corresponds to   perturbation with respect to the diffusion term $\bbB\bbA^{-1}\bbB^*$ if $
\bu^m$ is bounded. However, for $\alpha \ll 1$,  the advection term $\bbD$ becomes dominant and the Schur complement changes from diffusion--reaction to transport--reaction and it can destroy the spectral equivalence. For a fully computable preconditioner associated with \eqref{abstract-prec}, we approximate the top diagonal block by $(\alpha^{-1} \cI - \alpha^{-1}\nabla(\vdiv) )^{-1}$, 
%\mathbf{M}_{\vdiv}^{-1}$
where $\cI - \nabla(\vdiv)$ 
%\mathbf{M}_{\vdiv}^{-1}$ 
is the Raviart--Thomas full matrix. To obtain an approximation of the Schur complement we complete the square of the positive symmetric operator 
\[ (\bbB^*+ \alpha^{-1}\bbD)^*\alpha \bbA^{-1}(\bbB^*+ \alpha^{-1}\bbD).\]
For sake of robust scaling, we use the approximation $$(\varrho_0 \bbC + \alpha \bbB\bbA^{-1}\bbB^* +  \bbB\bbA^{-1}\bbD + \rrb{\alpha^{-1}}
\bbD^*\bbA^{-1}\bbD+\rrb{\bbD^*\bbA^{-1}\bbB^*})^{-1},$$ which corresponds to a modified ADR operator of the form 
\[ (\varrho_0 \cI - {\alpha}\Delta +  \bu^m\cdot \nabla + \rrb{\alpha^{-1}}
|\bu^m|^2 \cI \,\rrb{+\,\bbD^*\bbA^{-1}\bbB^*} )^{-1},\]
where $\cI$ 
%M_\varphi$ 
is now the concentration mass matrix and $-\Delta$ and  $ \bu^m\cdot\nabla$
are  the discrete diffusion and convection operators, both  requiring stabilisation (we use symmetric interior penalty and SUPG, respectively). Consequently, this entire term is the inverse of the block % for example, SUPG 
\begin{align} \label{eq:S-block-def} %N_\varphi^{\mathrm{SUPG}} = 
\nonumber \langle\cS^{-1}\varphi_h,\psi_h \rangle & =  \int_\Omega (\varrho_0 + %\alpha^{-1}  
³|\bu_h^m|^2) \varphi_h\,\psi_h +  \alpha \sum_{K\in \cT_h} \int_K  \nabla \varphi_h \cdot \nabla \psi_h + \sum_{F\in \mathcal{F}_h}{\int_F} \frac{\alpha}{h_F} [\![\varphi_h]\!]\,[\![\psi_h]\!] \\
&  \quad - \alpha \sum_{F\in \mathcal{F}_h} \int_F  \{\!\{\nabla \varphi_h\}\!\}\cdot[\![\psi_h\bn_F]\!] - \alpha \sum_{F\in \mathcal{F}_h}{\int_F}  \{\!\{\nabla \psi_h\}\!\}\cdot[\![\varphi_h\bn_F]\!]  \nonumber \\
&  \quad 
+ 
\sum_{K\in \cT_h} \int_K (\bu_h^m\cdot \nabla \varphi_h)\,\psi_h 
- \sum_{F\in \mathcal{F}_h} \int_F (\bu_h^m \cdot \bn_F)\varphi_h^{\mathsf{upw}}\,[\![\psi_h]\!]
\nonumber \\
& \quad +  \sum_{K\in \cT_h}\tau_K \int_K ((\varrho_0+|\bu_h^m|^2)\varphi_h + \bu_h^m \cdot \nabla \varphi_h - \alpha \Delta \varphi_h)(\bu^{m}_h\cdot\nabla \psi_h), 
\end{align}
where 
$\tau_K = \bigl( \frac{2|\bu_h^m|}{h_K}+ \frac{4\alpha}{h_K^2} + \varrho_0 + |\bu_h^m|^2\bigr)^{-1} >0$ is the typical SUPG stabilisation parameter and 
$\varphi_h^{\mathsf{upw}}$ is the typical upwind concentration taking the value $\varphi_h^+$ if $\bu_h^m \cdot \bn_F \geq 0$ and $\varphi_h^-$ otherwise 
(see, e.g. \cite{franca1992stabilized}). 

Back to the continuous case in  the Banach context, we note that the mass matrix $\bbA$ is no longer an isomorphism in the correct topology. A natural choice is 
\[ \bbJ_{\bV}: \bV \to \bV', \quad \langle \bbJ_{\bv}\bzeta,\bxi\rangle := (\bzeta,\bxi) + \langle |\vdiv\bzeta|^{4}\vdiv \bzeta, \vdiv \bxi\rangle_{\rL^{6/5}(\Omega)\times \rL^6(\Omega)}, \]
that is, using the notation in \eqref{eq:duality2}, $\bbJ_{\bQ} \equiv \cI - \nabla\cJ_6(\vdiv)$. This map is nonlinear, but it is coercive, bounded, and strictly monotone. Then, the Banach counterpart of \eqref{abstract-prec} is 
\begin{equation}\label{abstract-prec-Banach} 
 \mathrm{diag}(\cR,\cS) = \begin{pmatrix}\alpha \bbJ_{\bV}^{-1} & 0 \\ 0 & (\varrho_0 \bbJ_{\bQ} + \alpha \bbB\bbJ_{\bV}^{-1}(\bbB^* +\alpha^{-1}\bbD))^{-1}
\end{pmatrix}, 
\end{equation}
where $\bbJ_{\bQ} \equiv \cJ_{6}$ is the duality mapping (cf. \eqref{eq:duality}). 
%\cgray{Furthermore,  the Banach and continuous Schur complement block contains a $p$-Laplacian structure and reads \textcolor{red}{[not sure about this one (I think the correct map should be non-local):]}
%\[ \langle\cS^{-1}\varphi,\psi\rangle = \varrho_0 \int_\Omega |\varphi|^{4}\varphi\,\psi + \alpha \int_\Omega |\nabla \varphi|^{-4/5}\nabla\varphi \cdot \nabla \psi + \int_\Omega (\bu^m \cdot \nabla \varphi)\psi.\]
It is important to remark that, while the Banach case provides the robust well-posedness, the robust boundedness of the operators, and the robustness of scaling, it is sufficient with implementing in practice the Hilbert version of the block-diagonal preconditioner \eqref{abstract-prec} with \eqref{eq:S-block-def}.  
}
\end{comment}
% \rd
{We construct block-diagonal preconditioners for the Darcy--Forchheimer velocity and pressure blocks following 
\cite{dhpr_sinum26}. On the other hand, and differently from \cite{cai2026parameter} and \cite{dhpr_sinum26}, we readily note that the sub-blocks associated with the ARD equations consist of a perturbed, \emph{non-symmetric} saddle-point of the form 
\begin{equation}\label{abstract-block} \begin{pmatrix} \alpha^{-1}\bbA & \bbB^* + \alpha^{-1}\bbD \\ \bbB & - \varrho_0 \bbC\end{pmatrix},\end{equation}
where we are abusing of notation to emphasise the parameter dependence of the different blocks. 
Let us first concentrate on the discrete case since it simplifies the discussion as $\bu_h^m \in \bL^{\infty}(\Omega)$ and so the Banach structure is not required. In this case, $\bbA:\bH(\vdiv,\Omega)\to \bH(\vdiv,\Omega)'$ denotes the  $\bL^2(\Omega)$ mass map in $\bH(\vdiv,\Omega)$, $\bbB:\bH(\vdiv,\Omega)\to \rL^2(\Omega)'$ is the divergence operator, $\bbD$ is the \textit{mixed} advection coupling operator ($\bbD\varphi \approx \bu^m\varphi$), and $
\bbC:\rL^2(\Omega) \to \rL^2(\Omega)'$ is the scalar mass map in $\rL^2(\Omega)$. 
Therefore a typical Riesz map preconditioner, also based on a Schur complement approximation, is motivated by \cite[Sect. 3.1]{benzi2005numerical} (see also \cite[Section 1]{baerland2020observation} for the symmetric, unperturbed case). The guiding principle is to scale the flux space with $\alpha^{-1}$, the concentration space with $\varrho_0$, and ensure that the Schur complement is spectrally equivalent to a properly weighted  operator.
In this case, a block-diagonal preconditioner  
reads as follows (using {a complete Schur complement structure for the (2,2)-block and an incomplete one for the (1,1)-block}) 
\begin{equation}\label{abstract-prec} 
 \mathrm{diag}(\cR,\cS) = \begin{pmatrix}(\alpha^{-1}\bbA +  \bbB^*\bbB)^{-1} & 0 \\ 0 & ([1+\varrho_0]\,\bbC)^{-1}+(\varrho_0 \bbC + \alpha \bbB\bbA^{-1}\bbB^* + \bbB\bbA^{-1}\bbD)^{-1}
\end{pmatrix},
\end{equation}
or, equivalent Schur complement structure for the (1,1)-block
(and note that it is also possible to derive block-triangular versions, see for e.g., \cite{olshanskii2007pressure}). 
For the the advection term the preconditioner part $\bbB\bbA^{-1}\bbD$ is not symmetric and it corresponds to   perturbation with respect to the diffusion term $\bbB\bbA^{-1}\bbB^*$ if $
\bu^m$ is bounded. However, for $\alpha \ll 1$,  the advection term $\bbD$ becomes dominant and the Schur complement changes from diffusion--reaction to transport--reaction and it can destroy the spectral equivalence. For a fully computable preconditioner associated with \eqref{abstract-prec}, we approximate the top diagonal block by $(\alpha^{-1} \cI - \alpha^{-1}\nabla(\vdiv) )^{-1}$, 
%\mathbf{M}_{\vdiv}^{-1}$
where $\cI - \nabla(\vdiv)$ 
%\mathbf{M}_{\vdiv}^{-1}$ 
is the Raviart--Thomas full matrix. \begin{comment}{To obtain an approximation of the Schur complement we complete the square of the positive symmetric operator 
\[ (\bbB^*+ \alpha^{-1}\bbD)^*\alpha \bbA^{-1}(\bbB^*+ \alpha^{-1}\bbD).\]
For sake of robust scaling, we use the approximation $$(\varrho_0 \bbC + \alpha \bbB\bbA^{-1}\bbB^* +  \bbB\bbA^{-1}\bbD + \rrb{\alpha^{-1}}
\bbD^*\bbA^{-1}\bbD+\rrb{\bbD^*\bbA^{-1}\bbB^*})^{-1},$$ which corresponds to a modified ADR operator of the form 
\[ (\varrho_0 \cI - {\alpha}\Delta +  \bu^m\cdot \nabla + \rrb{\alpha^{-1}}
|\bu^m|^2 \cI \,\rrb{+\,\bbD^*\bbA^{-1}\bbB^*} )^{-1},\]
where $\cI$ 
%M_\varphi$ 
is now the concentration mass matrix and $-\Delta$ and  $ \bu^m\cdot\nabla$
are  the discrete diffusion and convection operators, both  requiring stabilisation (we use symmetric interior penalty and SUPG, respectively). Consequently, this entire term is the inverse of the block
\begin{align} \label{eq:S-block-def} %N_\varphi^{\mathrm{SUPG}} = 
\nonumber \langle\cS^{-1}\varphi_h,\psi_h \rangle & =  \int_\Omega (\varrho_0 + %\alpha^{-1}  
³|\bu_h^m|^2) \varphi_h\,\psi_h +  \alpha \sum_{K\in \cT_h} \int_K  \nabla \varphi_h \cdot \nabla \psi_h + \sum_{F\in \mathcal{F}_h}{\int_F} \frac{\alpha}{h_F} [\![\varphi_h]\!]\,[\![\psi_h]\!] \\
&  \quad - \alpha \sum_{F\in \mathcal{F}_h} \int_F  \{\!\{\nabla \varphi_h\}\!\}\cdot[\![\psi_h\bn_F]\!] - \alpha \sum_{F\in \mathcal{F}_h}{\int_F}  \{\!\{\nabla \psi_h\}\!\}\cdot[\![\varphi_h\bn_F]\!]  \nonumber \\
&  \quad 
+ 
\sum_{K\in \cT_h} \int_K (\bu_h^m\cdot \nabla \varphi_h)\,\psi_h 
- \sum_{F\in \mathcal{F}_h} \int_F (\bu_h^m \cdot \bn_F)\varphi_h^{\mathsf{upw}}\,[\![\psi_h]\!]
\nonumber \\
& \quad +  \sum_{K\in \cT_h}\tau_K \int_K ((\varrho_0+|\bu_h^m|^2)\varphi_h + \bu_h^m \cdot \nabla \varphi_h - \alpha \Delta \varphi_h)(\bu^{m}_h\cdot\nabla \psi_h), 
\end{align}
where 
$\tau_K = \bigl( \frac{2|\bu_h^m|}{h_K}+ \frac{4\alpha}{h_K^2} + \varrho_0 + |\bu_h^m|^2\bigr)^{-1} >0$ is the typical SUPG stabilisation parameter and 
$\varphi_h^{\mathsf{upw}}$ is the typical upwind concentration taking the value $\varphi_h^+$ if $\bu_h^m \cdot \bn_F \geq 0$ and $\varphi_h^-$ otherwise 
(see, e.g. \cite{franca1992stabilized}).\end{comment}
The discrete diffusion and convection operators $-\Delta$ and $\bu^m\cdot\nabla$ require some stabilisation (we use symmetric interior penalty and SUPG, respectively). The operator~$\varrho_0\cI-\alpha\Delta+\bu^m\cdot\nabla$ is the inverse of the following block
\begin{align} \label{eq:S-block-def} %N_\varphi^{\mathrm{SUPG}} = 
\nonumber \langle\cS^{-1}\varphi_h,\psi_h \rangle & =  \int_\Omega \varrho_0 \varphi_h\,\psi_h +  \alpha \sum_{K\in \cT_h} \int_K  \nabla \varphi_h \cdot \nabla \psi_h + \sum_{F\in \mathcal{F}_h}{\int_F} \frac{\alpha}{h_F} [\![\varphi_h]\!]\,[\![\psi_h]\!] \\
&  \quad - \alpha \sum_{F\in \mathcal{F}_h} \int_F  \{\!\{\nabla \varphi_h\}\!\}\cdot[\![\psi_h\bn_F]\!] - \alpha \sum_{F\in \mathcal{F}_h}{\int_F}  \{\!\{\nabla \psi_h\}\!\}\cdot[\![\varphi_h\bn_F]\!]  \nonumber \\
&  \quad 
+ 
\sum_{K\in \cT_h} \int_K (\bu_h^m\cdot \nabla \varphi_h)\,\psi_h 
- \sum_{F\in \mathcal{F}_h} \int_F (\bu_h^m \cdot \bn_F)\varphi_h^{\mathsf{upw}}\,[\![\psi_h]\!]
\nonumber \\
& \quad +  \sum_{K\in \cT_h}\tau_K \int_K (\varrho_0\varphi_h + \bu_h^m \cdot \nabla \varphi_h - \alpha \Delta \varphi_h)(\bu^{m}_h\cdot\nabla \psi_h), 
\end{align}
where 
$\tau_K = \bigl( \frac{2|\bu_h^m|}{h_K}+ \frac{4\alpha}{h_K^2} + \varrho_0 \bigr)^{-1} >0$ is the typical SUPG stabilisation parameter and 
$\varphi_h^{\mathsf{upw}}$ is the typical upwind concentration taking the value $\varphi_h^+$ if $\bu_h^m \cdot \bn_F \geq 0$ and $\varphi_h^-$ otherwise 
(see, e.g. \cite{franca1992stabilized}).
Back to the continuous case in  the Banach context, we note that the mass matrix $\bbA$ is no longer an isomorphism in the correct topology. A natural choice is 
\[ \bbJ_{\bV}: \bV \to \bV', \quad \langle \bbJ_{\bv}\bzeta,\bxi\rangle := (\bzeta,\bxi) + \langle |\vdiv\bzeta|^{4}\vdiv \bzeta, \vdiv \bxi\rangle_{\rL^{6/5}(\Omega)\times \rL^6(\Omega)}, \]
that is, using the notation in \eqref{eq:duality2}, $\bbJ_{\bQ} \equiv \cI - \nabla\cJ_6(\vdiv)$. This map is nonlinear, but it is coercive, bounded, and strictly monotone. Then, the Banach counterpart of \eqref{abstract-prec} is 
\begin{equation}\label{abstract-prec-Banach} 
 \mathrm{diag}(\cR,\cS) = \begin{pmatrix}\alpha \bbJ_{\bV}^{-1} & 0 \\ 0 & (\varrho_0 \bbJ_{\bQ} + \alpha \bbB\bbJ_{\bV}^{-1}(\bbB^* +\alpha^{-1}\bbD))^{-1}
\end{pmatrix}, 
\end{equation}
where $\bbJ_{\bQ} \equiv \cJ_{6}$ is the duality mapping (cf. \eqref{eq:duality}). 
%\cgray{Furthermore,  the Banach and continuous Schur complement block contains a $p$-Laplacian structure and reads \textcolor{red}{[not sure about this one (I think the correct map should be non-local):]}
%\[ \langle\cS^{-1}\varphi,\psi\rangle = \varrho_0 \int_\Omega |\varphi|^{4}\varphi\,\psi + \alpha \int_\Omega |\nabla \varphi|^{-4/5}\nabla\varphi \cdot \nabla \psi + \int_\Omega (\bu^m \cdot \nabla \varphi)\psi.\]
It is important to remark that, while the Banach case provides the robust well-posedness, the robust boundedness of the operators, and the robustness of scaling, it is sufficient with implementing in practice the Hilbert version of the block-diagonal preconditioner \eqref{abstract-prec} with \eqref{eq:S-block-def}.  
}

% \begin{gather}

\begin{comment}{First, let us denote by \cgray{Since the nonlinear velocity map, say $A: \bL^3(\Omega)\to \bL^{3/2}(\Omega)$ is $A\bu = \rF|\bu|\bu$, we can define a scalar operator that acts as the natural duality map, say $B:\rL^{3/2}(\Omega)\subset \rW^{-1,3/2}(\Omega) \to \rW^{1,3}(\Omega)$ defined through a nonlinear elliptic problem. Then a more accurate notation for $B$ would be 
$B\phi = -\vdiv(\rF^{-3/2}|\nabla \phi|^{-1/2}\nabla\phi)$ or in operator preconditioner form $-\vdiv(\rF^{-3/2}|\nabla \bullet|^{-1/2}\nabla\bullet)^{-1}$}
$[\rF{\mathcal{H}(\hat{\bu})}]^{-1}{\Delta_{\frac32}}$  the (singular) weighted \rrb{$\frac32$}-Laplace operator associated with the space $\rW^{1,\frac{3}{2}}_{\Gamma_p}(\Omega)$, defined as  
  \[\langle [\rF{\mathcal{H}(\hat{\bu})}]^{-1}\Delta_{\frac32}\phi,\chi\rangle := \int_\Omega [\rF{\mathcal{H}(\hat{\bu})}]^{-1}|\nabla\phi|^{-\frac12}\nabla \phi \cdot \nabla \chi, \qquad \forall \chi \in \rW^{1,3}(\Omega).\] 
% and \[ \quad\rd{\langle\mathcal{J}_{\frac65}(\bzeta),\btau\rangle:=\int_{\Omega}|\vdiv\bzeta|^{-\frac45}\vdiv\bzeta\vdiv\btau,} \quad\langle \mathcal{J}_{\frac65}(\psi),\varphi\rangle:=\int_{\Omega}|\psi|^4\psi\varphi, \quad \forall\psi,\varphi\in\Psi.\]
\end{comment}
\noindent
\begin{comment}{Following \cite[Equation~7]{baerland2020observation}, an alternative type of preconditioner can be considered for the ADR equations, different from \eqref{abstract-prec}, given by
\begin{equation}\label{abstract-prec.B2} 
 \mathrm{diag}(\cR,\cS) = \begin{pmatrix}(\alpha^{-1}\bbA +  \bbB^*\bbB)^{-1} & 0 \\ 0 & ([1+\varrho_0]\,\bbC)^{-1} + (\varrho_0\bbC+\alpha\bbB\bbA^{-1}\bbB+\bbB\bbA^{-1}\bbD)^{-1}
\end{pmatrix}.
\end{equation}
\end{comment}
Finally, inspired by \eqref{abstract-prec}, {and considering another formulation inspired by the primal formulation,} we propose the following preconditioners {($\cB^{1,1}_{\mathsf{DFARD}}, \cB^{1,2}_{\mathsf{DFARD}}, \cB^{2,1}_{\mathsf{DFARD}},$ and $\cB^{2,2}_{\mathsf{DFARD}}$)} for the coupled monolithic problem~\eqref{eq:strong.linearised}, where 
\begin{align}\label{precond.linear.B2}
\cB^{i,j}_{\mathsf{DFARD}}:=\begin{pmatrix}
\mathcal{P}_i & 0 & 0 & 0\\
0 &\mathcal{Q}_i& 0 & 0\\
0 & 0 & \mathcal{R}_j & 0\\
0 & 0 & 0 & \mathcal{S}_j
\end{pmatrix},
\end{align}
with \begin{align*}
\mathcal{P}_1&:={\bigl([\kappa^{-1}+\rF\mathcal{H}(\hat{\bu})]\,\mathcal{I}-[\kappa^{-1}+\rF\mathcal{H}(\hat{\bu})]\nabla\vdiv\bigr)^{-1}},\quad \text{and}\quad\mathcal{Q}_1:=[\kappa^{-1}+\rF\mathcal{H}(\hat{\bu})]^{-1}{\mathcal{I}}^{-1}, \\
\mathcal{P}_2&:={\bigl([\kappa^{-1}+\rF\mathcal{H}(\hat{\bu})]\mathcal{I}-\nabla\vdiv\bigr)^{-1}},\quad \text{and}\quad\mathcal{Q}_2:={\mathcal{I}^{-1}+(-\kappa\Delta)^{-1}-\bigl([\rF\mathcal{H}(\hat{\bu})]^{-1}\Delta_{\frac32}\bigr)^{-1}},
\\\mathcal{R}_1&:=(\alpha^{-1}\cI +  \nabla\cdot\vdiv+\alpha^{-1}\hat\bu\nabla\cdot)^{-1},\quad \text{and}\quad\mathcal{S}_1:={(1+\varrho_0
%\bbC
)\,\mathcal{I}^{-1}+ (\varrho_0\mathcal{I}-\alpha\Delta)^{-1}},
\\\mathcal{R}_2&:=(\alpha^{-1}\cI +  \nabla\cdot\vdiv)^{-1},\quad \text{and}\quad\mathcal{S}_2:={(1+\varrho_0
%\bbC
)\,\mathcal{I}^{-1}+ (\varrho_0\mathcal{I}-\alpha\Delta+\bu^m\cdot\nabla)^{-1}}.
\end{align*}
}
\begin{comment}
As a straightforward consequence of the above discussion, the following theorem provides a robust preconditioner for the full monolithic system. \rrb{I don't think we can claim this now. There is no longer an evident relation between Riesz/duality map preconditioning and Schur-complement based preconditioners for block non-symmetric systems}
\begin{theorem}[robust preconditioner for the linearized monolithic system]
Assume that
\((\hat{\bu},\hat{p}, \hat{\bzeta}, \hat{\varphi})\in\underline{\bV}\times \underline{\rQ}\times \underline{\bW}\times \underline{\Psi}\) {is} the state around which the Newton--Raphson linearisation is performed. 
    Then the following block diagonal operator \begin{align}\label{precond.linear}
\cB:=\begin{pmatrix}
\mathcal{P} & 0 & 0 & 0\\
0 & \mathcal{Q} & 0 & 0\\
0 & 0 & \mathcal{R} & 0\\
0 & 0 & 0 & \mathcal{S}
\end{pmatrix}
\end{align}
is a parameter-robust preconditioner for  \eqref{eq:newton}, in the sense that the condition numbers of the preconditioned matrix are uniformly bounded in \(\kappa, \rF,\alpha,\) and \(h\). 
Here \begin{align*}\mathcal{P}:=\bigl(\kappa^{-1}&\mathcal{I}+\rF\,\mathcal{H}(\hat{\bu})\mathcal{I}-\nabla\vdiv\bigr)^{-1}, \quad \text{and}\quad\mathcal{Q}:=\mathcal{I}^{-1}+(-\kappa\Delta)^{-1}-\bigl([\rF\mathcal{H}(\hat{\bu})]^{-1}\Delta_{\frac32}\bigr)^{-1},\\ \mathcal{R}&=\alpha (\cI - \nabla(\vdiv) )^{-1},\qquad\text{and} \qquad{\mathcal{S}=(\varrho_0 \cI - {\alpha}\Delta +  \bu^m\cdot \nabla )^{-1}}.\end{align*} 
\end{theorem}
\end{comment}
% \end{gather}
% \rrb{I think the last one has to be a duality mapping between $L^6$ and $L^{6/5}$. Call it something different, e.g. $\mathcal{J}_{\frac65}$ as in \cite{gatica2022p}}
\begin{comment}
{\begin{proof}
The derivations of \(\mathcal{P}\) and \(\mathcal{Q}\) can be found in detail in \cite[Theorem~5.2]{dhpr_sinum26}. Let us concentrate on obtaining the blocks~\(\mathcal{R}\) and \(\mathcal{S}.\) 
{Let us remark that,
\begin{align}\label{eq:precond.full.norm}
\tnorm{(\bzeta,\varphi)}^2_{\bW\times\Psi}&=\|\bzeta\|_{\bW}^2+\|\varphi\|^2_{\Psi}+\|(\varrho_0-g)^\frac12\varphi\|_{0,\Omega}^2\eqsim\alpha^{-1}\|\bzeta\|^2_{0,\Omega}+\|\vdiv\bzeta\|_{0,\frac65,\Omega}^2\nonumber\\
&+\inf_{\xi \in \alpha^{\frac12}\rH^1_{\Gamma_{\bu}}(\Omega)}\biggl(\|[1+(\varrho_0-g)^\frac12](\varphi-\xi)\|^2_{0,6,\Omega} + \alpha\|\nabla \xi\|^2_{0,\Omega} +\|(\varrho_0-g)^\frac12\xi\|_{0,\Omega}^2\biggr),
\end{align}
% where we are using the equivalent definition of the norm \(\|\bullet\|_{\Psi}+|(\varrho_0-g)^\frac12\bullet|_c\) as follows.
% \begin{align}\label{eq:inf.norm.psi}
% \|\varphi\|^2_{\Psi}+\|(\varrho_0-g)^\frac12\varphi\|_{0,\Omega}^2 := \inf_{\xi \in \alpha^{\frac12}\rH^1_{\Gamma_{\bu}}(\Omega)}\bigl(\|\varphi-\xi\|^2_{0,6,\Omega} + \alpha\|\nabla \xi\|^2_{0,\Omega} \bigr). \end{align}
Also, note that, the above infimum (cf. \eqref{eq:precond.full.norm}) is attained by an element \(\xi\in\rH^1_{\Gamma_{\bu}}(\Omega),\) which satisfies the following system of equations for a given \(\varphi\in\Psi,\)
\begin{align*}
-\alpha\Delta \xi+ &(\varrho_0-g)\xi-\mathsf{D}(\xi)(\varrho_0-g)^6|\varphi-\xi|^4(\varphi-\xi)=\varphi\quad \text{in }
\Omega,\\ \xi&=0\quad\text{on }\Gamma_{\bu},\quad \nabla \xi\cdot \bn=0\quad \text{on } \Gamma_p,\end{align*}
where \(\mathsf{D}(\xi):=\int_{\Omega}(\varrho_0-g)^6|\varphi-\xi|^6)^{-\frac23},\) \cgray{and where the existence and uniqueness of this system  follows from the classical Lax--Milgram Theorem. 

Next, employing the solution of  \eqref{eq:inf.norm.psi} in \eqref{eq:precond.full.norm} we deduce the following relation
\begin{align*} \tnorm{(\bzeta,\varphi)}^2_{\bW\times\Psi}&=(\alpha^{-1}\bzeta,\bzeta)+\langle\mathcal{J}_{\frac65}(\vdiv\bzeta),\vdiv\bzeta\rangle+\|\varphi-\xi\|^2_{0,6,\Omega} + \alpha\|\nabla \xi\|^2_{0,\Omega}+(\varrho_0-g)\|\varphi\|_{0,\Omega}^2\\
&=\langle(\alpha^{-1}\mathcal{I}-\nabla \mathcal{J}_{\frac65}(\vdiv))\bzeta,\bzeta\rangle+\langle\varphi-\xi,\varphi-\xi\rangle+\alpha\langle\nabla \xi,\nabla \xi\rangle+([\varrho_0-g]\varphi,\varphi)\\
&=\langle(\alpha^{-1}\mathcal{I}-\nabla \mathcal{J}_{\frac65}(\vdiv))\bzeta,\bzeta\rangle+\langle-\alpha\Delta \xi,-\alpha\Delta \xi\rangle+\alpha\langle-\Delta \xi,\xi\rangle+([\varrho_0-g]\varphi,\varphi),
\end{align*}
which further implies that 
\begin{align}\label{eq:final.precond}
    \tnorm{(\bzeta,\varphi)}^2_{\bW\times\Psi}&=\langle(\alpha^{-1}\mathcal{I}-\nabla \mathcal{J}_{\frac65}(\vdiv))\bzeta,\bzeta\rangle+\langle-\alpha\Delta \xi,-\alpha\Delta \xi+\xi\rangle+([\varrho_0-g]\varphi,\varphi)\nonumber\\
    &=\langle(\alpha^{-1}\mathcal{I}-\nabla \mathcal{J}_{\frac65}(\vdiv))\bzeta,\bzeta\rangle+\langle[(-\alpha\Delta)^{-1}+\mathcal{J}_{\frac65}^{-1}+(\varrho_0-g)^{-1}\mathcal{I}^{-1}](\varphi),\varphi\rangle.
\end{align}
Let us denote \(\mathcal{R}:={\bigl(\alpha^{-1}\mathcal{I}-\nabla \mathcal{J}_{\frac65}(\vdiv)\bigr)^{-1}},\) and \(\mathcal{S}:=[(-\alpha\Delta)^{-1}+\mathcal{J}_{\frac65}^{-1}+(\varrho_0-g)^{-1}\mathcal{I}^{-1}],\) and observe that \eqref{eq:final.precond} can be re-written in the following manner \[\tnorm{(\bzeta,\varphi)}^2_{\bW\times\Psi}=\langle\mathcal{R}(\bzeta),\bzeta\rangle+\langle\mathcal{S}(\varphi),\varphi\rangle,\quad\forall(\bzeta,\varphi)\in\bW\times\Psi.\]
Therefore,  proceeding similarly as in \cite[Equation~3.27]{cai2026parameter} (see also, \cite[Section 3]{mardal2011preconditioning}, \cite[Equation~4.4]{boon2021robust}),  we can assert that \eqref{precond.linear} is a  \((\kappa,\rF,\alpha, h)\)-robust block diagonal preconditioner for  system \eqref{eq:newton}.
}}
\end{proof}
\end{comment}




\section{Numerical results}\label{sec:numer}
In this section, we illustrate the accuracy and performance of the proposed scheme and operator preconditioning through several numerical experiments. We verify that the convergence rates are robust. The numerical implementations are based on the open-source  finite element libraries \texttt{Gridap} \cite{badia22} and \texttt{FEniCS} \cite{alnaes2015fenics}.
% , and all numerical tests can be accessed from~
% \url{https://github.com/rishidas-04/RobustCoupledDarcyForchheimer-AdvectionReactionDiffusion.jl}. 

\smallskip
\noindent\textbf{Example~1 (convergence against smooth solutions in 2D).}
We first verify convergence with  unity parameters $\rho_0 = \alpha = \kappa = \alpha_0 = \alpha_1 =1.$ The domain is the unit square $\Omega = (0,1)^2$, with $\Gamma_{\bu}$ the bottom and left edges and $\Gamma_p$ the remainder of the boundary. The manufactured velocity, pressure, and concentration are 
\begin{gather*}
 \bu = \begin{pmatrix}
     \cos(\pi x)\sin(\pi y)\\-\sin(\pi x) \cos(\pi y)
 \end{pmatrix}, \quad p = \sin(\pi x) \sin (\pi y), 
 \quad \varphi = \cos(2\pi x) \cos (2\pi y).
\end{gather*}
The Darcy--Forchheimer velocity is divergence-free, and the forcing term is chosen as 
$ \bf(\varphi) = (\alpha_0 + \alpha_1 \varphi) \bg$, where $\bg = (0,1)^{\tt t}$, which clearly fulfills the condition \eqref{eq:lip.bound.f}. With these solutions we construct an exact total flux, and additional source terms {(so that the momentum balance is satisfied)} and non-homogeneous essential and natural boundary conditions required for the convergence assessment. The error history associated with the 
proposed mixed finite element  method on a sequence of successively refined partitions of the domain  are presented in Table~\ref{num:conv.rates-unity}, where we also tabulate rates of error decay computed as 
$\texttt{rate}  =\log(e_{(\bullet)}/\tilde{e}_{(\bullet)})[\log(h/\tilde{h})]^{-1}$,
where $e,\tilde{e}$ denote errors generated on two
consecutive  meshes of sizes $h$ and~$\tilde{h}$, respectively. The results indicate an optimal $O(h^{k+1})$ convergence in all fields, where we are using the natural norms, without parameter-weighting, and also the parameter-weighted norms. {Sample numerical solutions (obtained with the second-order scheme and on the finest mesh resolution) are portrayed in Figure~\ref{fig:ex01}.}




\begin{table}[!t]
\setlength{\tabcolsep}{3pt}
\centering
{\footnotesize\begin{tabular}{|rcccccccccc|}
\hline
DoF  &    $h$ & $\|\bu-\bu_h\|_{3,\vdiv,\Omega}\!$  &  \texttt{rate} &  $\|p-p_h\|_{0,\Omega}$  &  \texttt{rate} & \(\|\bzeta-\bzeta_h\|_{\vdiv_{6/5},\Omega}\)& \texttt{rate} &\(\|\varphi-\varphi_h\|_{0,6,\Omega}\)&\texttt{rate} &  \texttt{it} \\
\hline
\multicolumn{11}{|c|}{Errors and convergence rates for $k = 0$}\\
\hline
       40 & 0.7071 & 5.07e-01 & $\star$ & 2.44e-01 & $\star$ & 3.50e+01 & $\star$ & 7.09e-01 & $\star$ & 5 \\
    160 & 0.3536 & 2.92e-01 & 0.797 & 1.29e-01 & 0.921 & 1.88e+01 & 0.900 & 3.43e-01 & 1.050 & 6\\
    640 & 0.1768 & 1.51e-01 & 0.946 & 6.52e-02 & 0.982 & 9.68e+00 & 0.955 & 2.02e-01 & 0.764 & 6\\
   2560 & 0.0884 & 7.64e-02 & 0.987 & 3.27e-02 & 0.996 & 4.88e+00 & 0.988 & 1.03e-01 & 0.967 & 5\\
  10240 & 0.0442 & 3.83e-02 & 0.997 & 1.64e-02 & 0.999 & 2.45e+00 & 0.997 & 5.19e-02 & 0.992 & 5\\
  40960 & 0.0221 & 1.91e-02 & 0.999 & 8.18e-03 & 1.000 & 1.22e+00 & 0.999 & 2.60e-02 & 0.998 & 6\\
\hline
\multicolumn{11}{|c|}{Errors and convergence rates for $k = 1$}\\
\hline
     128 & 0.7071 & 1.80e-01 &$\star$& 7.62e-02 & $\star$ & 1.92e+01 &$\star$  & 3.14e-01 & $\star$ & 7 \\
    512 & 0.3536 & 4.78e-02 & 1.914 & 1.97e-02 & 1.954 & 5.20e+00 & 1.884 & 1.23e-01 & 1.356 & 7\\
   2048 & 0.1768 & 1.24e-02 & 1.943 & 4.96e-03 & 1.986 & 1.37e+00 & 1.923 & 3.33e-02 & 1.883 & 8\\
   8192 & 0.0884 & 3.15e-03 & 1.979 & 1.24e-03 & 1.997 & 3.47e-01 & 1.985 & 8.54e-03 & 1.962 & 7\\
  32768 & 0.0442 & 7.92e-04 & 1.993 & 3.11e-04 & 1.999 & 8.68e-02 & 1.997 & 2.15e-03 & 1.991 & 8\\
 131072 & 0.0221 & 1.98e-04 & 1.998 & 7.78e-05 & 2.000 & 2.17e-02 & 1.999 & 5.38e-04 & 1.998 & 8\\
 \hline
 \hline
DoF  &    $h$ & $\|\bu-\bu_h\|_{\bV}\!$  &  \texttt{rate} &  $\|p-p_h\|_{\widehat{\rQ}(\cT_h)}$  &  \texttt{rate} & \(\|\bzeta-\bzeta_h\|_{\bW}\)& \texttt{rate} &\(\|\varphi-\varphi_h\|_{\widehat{\Psi}(\cT_h)}\)&\texttt{rate} &  \texttt{it} \\
\hline
\multicolumn{11}{|c|}{Errors and convergence rates for $k = 0$}\\
\hline
    40 & 0.7071 & 9.83e-01 & $\star$ & 6.05e-02 & $\star$ & 3.50e+01 & $\star$ & 3.66e-01  & $\star$ & 6 \\
    160 & 0.3536 & 5.56e-01 & 0.824 &1.94e-02 & 1.640& 1.88e+01 & 0.900 & 1.77e-01&1.050& 7\\
    640 & 0.1768 & 2.88e-01 & 0.947 &4.96e-03 & 1.969& 9.68e+00 & 0.955 & 5.14e-02& 1.782& 7\\
   2560 & 0.0884 & 1.46e-01 & 0.986 &1.25e-03&1.990 & 4.88e+00 & 0.988 &1.33e-02 &1.948& 7\\
  10240 & 0.0442 & 7.29e-02 & 0.996 & 3.13e-04&1.996& 2.45e+00 & 0.997 & 3.36e-03&1.987 & 7\\
  40960 & 0.0221 & 3.65e-02 & 0.999 &7.84e-05 &1.998& 1.22e+00 & 0.999 &8.42e-04  & 1.997 & 7\\
\hline
\multicolumn{11}{|c|}{Errors and convergence rates for $k = 1$}\\
\hline
      128 & 0.7071 & 3.31e-01 &$\star$& 1.55e-01 & $\star$ & 1.92e+01 &$\star$  & 3.68e-01 & $\star$ & 7 \\
    512 & 0.3536 & 8.86e-02 & 1.902 & 4.68e-02  & 1.730 & 5.20e+00 & 1.884 &1.81e-01&1.019& 8\\
   2048 & 0.1768 & 2.30e-02 & 1.946 & 1.23e-02  & 1.931 & 1.37e+00 & 1.923 & 5.61e-02 &1.041 & 9\\
   8192 & 0.0884 & 5.82e-03 & 1.981 & 3.11e-03 & 1.981 & 3.47e-01 & 1.985 & 8.82e-02 & 1.938& 8\\
  32768&0.0442&1.46e-03&1.994&7.80e-04&1.994&8.68e-02&1.997&5.81e-03&1.984&8\\
  131072&0.0221&3.66e-04&1.998&1.95e-04&1.998&2.17e-02&1.999&1.46e-03&1.994&9\\
\hline
\end{tabular}}
\caption{Example 1. Error history against smooth solutions for different polynomial degrees \((k=0,1)\), and iteration count for the nonlinear Newton--Raphson solver. Here we use unit parameters, and the non-weighted norms (top rows), and the parameter-weighted norms (bottom rows) for all fields.} 
\label{num:conv.rates-unity}
\end{table}

\begin{figure}[t!]
    \centering
\includegraphics[width=0.245\linewidth]{ex01u.png}
\includegraphics[width=0.245\linewidth]{ex01p.png}
\includegraphics[width=0.245\linewidth]{ex01zeta.png}
\includegraphics[width=0.245\linewidth]{ex01phi.png}
    \caption{Example 1. Approximate velocity magnitude and line integral contours, pressure profile, flux magnitude and line integral contours, and concentration, all computed with a second-order scheme.}
    \label{fig:ex01}
\end{figure}


\smallskip
\noindent\textbf{Example 2  (convergence relative to fine-mesh reference solutions).} 
Let $\Omega=(0,2)\times(0,1)$ and $(\bf,g,m)=(\mathbf{0},0,0)$. The boundary conditions are
$\bu\cdot\bn=-0.25\,y(1-y)$ and $\varphi = 2$ on the left edge, $\bu\cdot\bn=\bzeta\cdot\bn=0$ on the top and bottom edges, $p=0$  and $\varphi = 1$ on the right edge. Following a similar test in \cite{dhpr_sinum26},   we consider the following heterogeneous  permeability
\[
\kappa(\bx)=\kappa_0 \left(1+\exp\!\left(-\tfrac{1}{2}\,[5 -4x +7y + \cos x]^2\right)\right).
\]
Convergence rates are computed against reference solutions obtained via  additional mesh refinements. Errors are reported in both the standard norms of Table~\ref{num:conv.rates-unity} and the weighted norms $\mathbf{V}, \widehat{\rQ}, \mathbf{W}, \widehat{\Psi}$. While the experimental order of convergence seems affected by the extreme parameter values when using the usual norms, the optimality is restored when using the weighted norms. Optimal rates are observed for all variables and both polynomial degrees. The Newton--Raphson iteration counts (Table~\ref{tab:conv.rates.weighted}) remain low even for extreme parameter values. We show the approximate solutions on a coarse mesh in Figure \ref{fig:ex02}. 

\begin{figure}[t!]
    \centering
\includegraphics[width=0.245\linewidth]{ex02u}
\includegraphics[width=0.245\linewidth]{ex02p}
\includegraphics[width=0.245\linewidth]{ex02zeta}
\includegraphics[width=0.245\linewidth]{ex02phi}
    \caption{Example 2. Approximate velocity magnitude and line integral contours, pressure profile, flux magnitude and line integral contours, and concentration, all computed with a lowest-order scheme on a coarse.}
    \label{fig:ex02}
\end{figure}


\begin{comment}
Recall from \cite{baerland2020observation} that in the space $\rQ_h,$ and \(\Psi_h\) the realisation of the required discrete weighted Laplacian operator that is the Riesz map for  $\kappa^{\frac12}\rH^1_{h,\Gamma_p}(\Omega)$ is as follows
\begin{align*} \langle -\vdiv(\kappa(\bx) \nabla p_h), q_h\rangle & = \sum_{K\in \cT_h} (\kappa(\bx) \nabla p_h, \nabla q_h)_K 
%- \sum_{F\in \cE_h\cup \cE_h^{\Gamma_p}}  (\{\!\!\{\kappa(\bx) \nabla p \}\!\!\}\cdot \bn_F,q)_F \\
%& \, - \sum_{F\in \cE_h\cup \cE_h^{\Gamma_p}}  (\{\!\!\{\kappa(\bx) \nabla q \}\!\!\}\cdot \bn_F,p)_F 
+ \sum_{F\in \cE_h\cup \cE_h^{\Gamma_p}} \frac{1}{h_F}(\kappa(\bx)[\![ p_h ]\!],[\![ q_h ]\!])_F,\end{align*}
where $\cE_h$ denotes the set of internal facets, $\cE_h^{\Gamma_p}$ denotes the set of facets lying on the sub-boundary $\Gamma_p$, and 
$[\![ \bullet ]\!]$ represents the % average and 
jump operator across a facet (and adopt the convention that it reduces to the identity if the face lies on the boundary).  In turn, for the space $\rF^{-\frac13}\rW^{1,\frac32}_{h,\Gamma_p}(\Omega)$  the natural operator (duality map) is the nonlinear discrete $s-$Laplacian
\begin{align*} \langle -\vdiv (\rF^{-1}|\nabla p_h|^{s-2}\nabla p_h),q_h\rangle &= \sum_{K\in \cT_h} \rF^{-1} (|\nabla p_h|^{s-2}
\nabla p_h, \nabla q_h)_K \\ &\quad + \sum_{F\in \cE_h\cup \cE_h^{\Gamma_p}} \rF^{-1} h_F^{1-s}(
|[\![ p_h ]\!]|^{s-2}
[\![ p_h ]\!],[\![ q_h ]\!])_F,  \end{align*}
with $s = \frac32$ (see, e.g., \cite{burman2008discontinuous}). 
However, since we require a linear operator to define the preconditioner-induced norm, we use instead the following linearisation around the reference pressure (cf. \cite{dhpr_sinum26})
\begin{align*}
& \sum_{K\in \cT_h} \rF^{-1} (|\nabla p_{\mathrm{ref}}|^{s-2}
\nabla p_h, \nabla q_h)_K + \sum_{F\in \cE_h\cup \cE_h^{\Gamma_p}} \rF^{-1}(s-1)h_F^{1-s}(
%\{\!\!\{|\nabla p_{\mathrm{ref}}|^{s-2}\}\!\!\}
|[\![ p_{\mathrm{ref}} ]\!]|^{s-2}
[\![ p_h ]\!],[\![ q_h ]\!])_F \\
&\quad + \sum_{K\in \cT_h} \rF^{-1}(s-2) (|\nabla p_{\mathrm{ref}}|^{s-4}
(\nabla p_{\mathrm{ref}}\cdot\nabla p_h),(\nabla p_{\mathrm{ref}} \cdot \nabla q_h))_K.  
\end{align*}
\end{comment}

\begin{table}[!t]\label{tab:conv.rates.weighted}
\setlength{\tabcolsep}{2.5pt}
\centering 
{\footnotesize \begin{tabular}{|rcccccccccc|}
\hline
DoF  &    $h$ & $\|\bu-\bu_h\|_{3,\vdiv}\!$  &  \texttt{rate} &  $\|p-p_h\|_{0,\Omega}$  &  \texttt{rate} & \(\|\bzeta-\bzeta_h\|_{\vdiv_{\frac65},\Omega}\)& \texttt{rate} &\(\|\varphi-\varphi_h\|_{0,6,\Omega}\)&\texttt{rate} &  \texttt{it} \\
% \hline
% \multicolumn{11}{|c|}{Errors and convergence rates for $k = 0$}\\
\hline
  72 & 0.7071 & 5.18e-03 & $\star$ & 1.02e+05 & $\star$ & 1.22e+00 & $\star$ & 3.71e+01 &  $\star$ & 5 \\
    168 & 0.3536 & 4.25e-03 & 0.285 & 3.86e+04 & 1.403 & 1.04e+00 & 0.222 & 2.04e+01 & 0.862&  5\\
    480 & 0.1768 & 4.56e-03 & -0.101 & 1.78e+04 & 1.122 & 1.00e+00 & 0.060 & 1.11e+01 & 0.880&  5\\
   1584 & 0.0884 & 2.77e-03 & 0.718 & 6.29e+03 & 1.497 & 5.17e-01 & 0.952 & 4.11e+00 & 1.432&  5\\
    5712 & 0.0442 & 1.53e-03 & 0.857 & 1.93e+03 & 1.707 & 2.83e-01 & 0.868 & 2.70e+00 & 0.605&  5\\
21648 & 0.0221 & 1.19e-03 & 0.383 & 1.53e+03 & 0.465 & 2.25e-01 & 0.182 & 2.14e+00 & 0.278 & 5 \\
\hline
% \multicolumn{11}{|c|}{Errors and convergence rates for $k = 1$}\\
% \hline

% \hline
% \hline
DoF  &    $h$ & $\|\bu-\bu_h\|_{\bV}\!$  &  \texttt{rate} &  $\|p-p_h\|_{\widehat{\rQ}_h}$  &  \texttt{rate} & \(\|\bzeta-\bzeta_h\|_{\bW}\)& \texttt{rate} &\(\|\varphi-\varphi_h\|_{\widehat{\Psi}_h}\)&\texttt{rate} &  \texttt{it} \\
% \hline
% \multicolumn{11}{|c|}{Errors and convergence rates for $k = 0$}\\
\hline
     72 & 0.7071 & 1.44e+01 & $\star$ & 7.54e+00 & $\star$ & 9.90e+00 & $\star$ & 6.77e+00 & $\star$ & 5 \\
    168 & 0.3536 & 6.95e+00 & 1.075 & 3.10e+00 & 1.281 & 6.71e+00 & 0.562 & 4.53e+00 & 0.579 & 5 \\
    480 & 0.1768 & 3.65e+00 & 0.905 & 1.57e+00 & 0.980 & 4.21e+00 & 0.673 & 2.70e+00 & 0.746 & 5\\
   1584 & 0.0884 & 1.64e+00 & 1.153 & 5.40e-01 & 1.542 & 1.94e+00 & 1.119 & 1.15e+00 & 1.229 & 5\\
   5712 & 0.0442 & 5.26e-01 & 1.642 & 2.93e-01 & 0.883 & 7.63e-01 & 1.343 & 6.81e-01 & 0.961 & 5\\
  21648 & 0.0221 & 2.18e-01 & 1.297 & 1.17e-01 & 1.630 & 3.17e-01 & 1.405 & 3.50e-01 & 0.985 & 5 \\
\hline
\end{tabular}}
\caption{Example 2. Error history relative to a fine-mesh reference solution: errors on successively refined grids, convergence rates, for $k=0$.  Newton--Raphson iteration counts corresponding to the nonlinear solver are also reported. Here we are using extreme values of the parameters \(\kappa_0=10^{-8},\rF=10^{3},\alpha=5\cdot 10^{-3},\rho_0=10^{-5},\) and both the usual norms (top row), the weighted norms (bottom row) for velocity, pressure, flux, and concentration, respectively.} \label{num:conv.rates}
\end{table}


\smallskip
\noindent\textbf{Example~3 (operator preconditioning).} In this example, we test the robustness of the preconditioners~{$\cB^{i,j}_{\mathsf{DFARD}},i,j=1,2$} for the fully-coupled linearised Darcy--Forchheimer and ADR system.  We stick to the domain \(\Omega=(0,1)^2,\) the lowest order
finite element family, and mixed boundary conditions. We perform a sequence of experiments in which the permeability of the material \(\kappa\), the Forchheimer coefficient   \(\rF\), the diffusivity coefficient \(\alpha,\)  the reaction coefficient \(\varrho_0,\) and the mesh size $h$ vary by several orders of magnitude.  

Note that in this context $-\Delta$ and  $ \bu^m\cdot\nabla$
are  the discrete diffusion and convection operators, and they both  require stabilisation. {To ensure robust invertibility, the classical consistency (average-times-jump) terms found in standard Symmetric Interior Penalty Galerkin (SIPG) methods \cite{arnold2002unified} are  omitted. This isolates a purely positive-definite interface penalty operator parametrised by the stabilisation penalty factor $\beta = 10$ and the local characteristic mesh sizes $h_e$ 
\begin{align*}
    \langle-\alpha\Delta\varphi_h,\psi_h\rangle=& \alpha\beta \sum_{K\in \cT_h} \int_K  \nabla \varphi_h \cdot \nabla \psi_h + \sum_{F\in \mathcal{F}_h}{\int_F} \frac{\alpha\beta}{h_F} [\![\varphi_h]\!]\,[\![\psi_h]\!]. 
 % \\& - \alpha \sum_{F\in \mathcal{F}_h} \int_F  \{\!\{\nabla \varphi_h\}\!\}\cdot[\![\psi_h\bn_F]\!] - \alpha \sum_{F\in \mathcal{F}_h}{\int_F}  \{\!\{\nabla \psi_h\}\!\}\cdot[\![\varphi_h\bn_F]\!], 
\end{align*}
Hyperbolic stability for the advection term is governed by a pointwise upwind numerical flux along the interior interfaces   and boundaries, utilizing the velocity field $\bu^m$. By defining the positive and negative part operators as $[x]_ = \max(x, 0)$ and $[x]_- = \min(x, 0)$, respectively, the corresponding block is 
\begin{align*} 
 \langle(\bu^m\cdot\nabla)\varphi_h,\psi_h \rangle  =& 
-\sum_{K\in \cT_h} \int_K (\bu_h^m\cdot \nabla \psi_h)\,\varphi_h 
+ \sum_{F\in \mathcal{F}_h} \int_F 
\bigl([(\bu^m)^+ \cdot \bn_F^+]_+ \varphi^+ + [(\bu^m)^+ \cdot \bn_F^+]_- \varphi^- \bigr) \llbracket \psi \rrbracket   \\
&\quad  + \sum_{F\in \Gamma_{\text{out}}\cup \Gamma_{\text{in}}}  \int_F [\bu^m \cdot \bn_F]_+ \varphi \psi, 
%- \sum_{F\in \mathcal{F}_h} \int_F (\bu_h^m \cdot \bn_F)\varphi_h^{\mathsf{upw}}\,[\![\psi_h]\!]\\&
% +  \sum_{K\in \cT_h}\tau_K \int_K ( \bu_h^m \cdot \nabla \varphi_h - \alpha \Delta \varphi_h)(\bu^{m}_h\cdot\nabla \psi_h), 
\end{align*}
where  $\llbracket \phi \rrbracket = \phi^+ - \phi^-$ denotes the standard jump operator across the facet $F$ associated with the outward unit normal $\bn_F^+$. This   retains the essential upwind stability and diffusion scales while keeping a diagonally dominant structure   \cite{franca1992stabilized,kanschat2008multigrid}.
Even with this, we note}  that the tested preconditioners are not robust in the advection-dominated regime (large mesh Péclet number $\mathrm{Pe}_h = \frac{|\hat{\mathbf{u}}|h}{\alpha} \gg 1$). Therefore, we restrict our attention to the diffusion-dominated regime, where $\mathrm{Pe}_h \ll 1$, in which the preconditioners exhibit the desired robustness, see Table~\ref{tab:DF-ARD-Prec}.  The  performance of the preconditioners is assessed through condition numbers of the preconditioned system and the number of MINRES iterations needed to converge, all against  mesh refinement. 

\begin{table}[t]
\centering
\footnotesize
\setlength{\tabcolsep}{5pt}
\renewcommand{\arraystretch}{1.15}

%------------------ TOP ROW -----------------11-%
{
\begin{minipage}{0.48\textwidth}
\centering
\begin{tabular}{|c |c|c| c|lll|}
\hline
$\rho_0$ & $\kappa$ & $\alpha$ & $\rF$ & \multicolumn{3}{c|}{$h$} \\
\cline{5-7}
 & & & & $2^{-2}$ & $2^{-3}$ & $2^{-4}$ \\
\hline

\multirow{3}{*}{$1$}
& $10^{-8}$ & $10^{-9}$ & $10^{6}$ & 1.8 (6) & 1.8 (6)& 1.8 (7) \\
& $10^{-4}$ & $10^{-5}$ & $1$ & 1.8 (10)& 1.8 (12) & 1.8 (14)\\
& $1$ & $10^{-1}$ & $10^{-6}$ &2.9 (27) &3.0 (28) &3.0 (28) \\
\hline

\multirow{3}{*}{$10^{-5}$}
& $10^{-8}$ & $10^{-9}$ & $10^{6}$ & 1.2 (8) &1.3 (10) &1.6 (14) \\
& $10^{-4}$ & $10^{-5}$ & $1$ &3.6 (31) &4.0 (35) &4.2 (35) \\
& $1$ & $10^{-1}$ & $10^{-6}$ &2.1 (22) &2.1 (22) &2.1 (22) \\
\hline

\multirow{3}{*}{$10^{-9}$}
& $10^{-8}$ & $10^{-9}$ & $10^{6}$ & 3.6 (24)&4.1 (26) &4.3 (28) \\
& $10^{-4}$ & $10^{-5}$ & $1$ &3.7 (32) &4.1 (35) &4.2 (35) \\
& $1$ & $10^{-1}$ & $10^{-6}$ &2.1 (22) &2.1 (22) &2.1 (22) \\
\hline
\end{tabular}
\end{minipage}}
\hfill
%%%%%%%%%%%%%%%%%%%%%%%%%%%%%%%%% 12
{\begin{minipage}{0.48\textwidth}
\centering
\begin{tabular}{|c |c|c| c|lll|}
\hline
$\rho_0$ & $\kappa$ & $\alpha$ & $\rF$ & \multicolumn{3}{c|}{$h$} \\
\cline{5-7}
 & & & & $2^{-2}$ & $2^{-3}$ & $2^{-4}$ \\
\hline

\multirow{3}{*}{$1$}
& $10^{-8}$ & $10^{-9}$ & $10^{6}$ & 1.8 (6) & 1.8 (6)& 1.8 (7) \\
& $10^{-4}$ & $10^{-5}$ & $1$ & 1.8 (10)& 1.8 (12) & 1.8 (14)\\
& $1$ & $10^{-1}$ & $10^{-6}$ &2.9 (28) &3.0 (29) &3.0 (29) \\
\hline

\multirow{3}{*}{$10^{-5}$}
& $10^{-8}$ & $10^{-9}$ & $10^{6}$ & 1.2 (8) &1.3 (10) &1.6 (14) \\
& $10^{-4}$ & $10^{-5}$ & $1$ &3.6 (32) &4.0 (36) &4.2 (36) \\
& $1$ & $10^{-1}$ & $10^{-6}$ &2.1 (23) &2.0 (23) &2.0 (23) \\
\hline

\multirow{3}{*}{$10^{-9}$}
& $10^{-8}$ & $10^{-9}$ & $10^{6}$ & 3.6 (25)&4.1 (27) &4.3 (28) \\
& $10^{-4}$ & $10^{-5}$ & $1$ &3.7 (32) &4.1 (35) &4.2 (37) \\
& $1$ & $10^{-1}$ & $10^{-6}$ &2.1 (23) &2.0 (23) &2.0 (23) \\
\hline
\end{tabular}
\end{minipage}}

\vspace{0.4cm}

%------------------ BOTTOM ROW ---------------21---%
{\begin{minipage}{0.48\textwidth}
\centering
\begin{tabular}{|c |c|c| c|lll|}
\hline
$\rho_0$ & $\kappa$ & $\alpha$ & $\rF$ & \multicolumn{3}{c|}{$h$} \\
\cline{5-7}
 & & & & $2^{-2}$ & $2^{-3}$ & $2^{-4}$ \\
\hline

\multirow{3}{*}{$1$}
& $10^{-8}$ & $10^{-9}$ & $10^{6}$ & 3.7 (21) & 4.1 (25)& 4.4 (25) \\
& $10^{-4}$ & $10^{-5}$ & $1$ & 3.6 (30)& 3.8 (32) & 3.5 (35)\\
& $1$ & $10^{-1}$ & $10^{-6}$ &2.9 (29) &3.0 (29) &3.0 (29) \\
\hline

\multirow{3}{*}{$10^{-5}$}
& $10^{-8}$ & $10^{-9}$ & $10^{6}$ & 3.7 (25) & 4.1 (29) & 4.4 (32) \\
& $10^{-4}$ & $10^{-5}$ & $1$ &3.7 (34) &4.1 (37) &4.2 (37) \\
& $1$ & $10^{-1}$ & $10^{-6}$ &1.9 (21) &1.9 (21) &1.9 (20) \\
\hline

\multirow{3}{*}{$10^{-9}$}
& $10^{-8}$ & $10^{-9}$ & $10^{6}$ & 3.7 (25)&4.1 (27) &4.4 (29) \\
& $10^{-4}$ & $10^{-5}$ & $1$ &3.7 (33) &4.1 (37) &4.2 (37) \\
& $1$ & $10^{-1}$ & $10^{-6}$ &1.9 (21) &1.9 (21) &1.9 (20) \\
\hline
\end{tabular}
\end{minipage}}
\hfill
%%%%%%%%%%%%%%%%%%%%%%%%%%%%%%%%%%%%%%%%%%%%%%%%%%
{\begin{minipage}{0.48\textwidth}
\centering
\begin{tabular}{|c |c|c| c|lll|}
\hline
$\rho_0$ & $\kappa$ & $\alpha$ & $\rF$ & \multicolumn{3}{c|}{$h$} \\
\cline{5-7}
 & & & & $2^{-2}$ & $2^{-3}$ & $2^{-4}$ \\
\hline

\multirow{3}{*}{$1$}
& $10^{-8}$ & $10^{-9}$ & $10^{6}$ & 3.7 (21) & 4.1 (25)& 4.4 (25) \\
& $10^{-4}$ & $10^{-5}$ & $1$ & 3.6 (30)& 3.8 (32) & 3.5 (35)\\
& $1$ & $10^{-1}$ & $10^{-6}$ &2.9 (29) &3.0 (30) &3.0 (31) \\
\hline

\multirow{3}{*}{$10^{-5}$}
& $10^{-8}$ & $10^{-9}$ & $10^{6}$ & 3.7 (25) &4.1 (29) &4.4 (32) \\
& $10^{-4}$ & $10^{-5}$ & $1$ &3.6 (34) &4.1 (37) &4.2 (37) \\
& $1$ & $10^{-1}$ & $10^{-6}$ &1.8 (21) &1.8 (21) & 1.8 (22) \\
\hline

\multirow{3}{*}{$10^{-9}$}
& $10^{-8}$ & $10^{-9}$ & $10^{6}$ & 3.7 (25)&4.1 (27) &4.4 (29) \\
& $10^{-4}$ & $10^{-5}$ & $1$ &3.7 (34) &4.1 (37) &4.2 (37) \\
& $1$ & $10^{-1}$ & $10^{-6}$ &1.8 (21) &1.8 (21) &1.8 (22) \\
\hline
\end{tabular}
\end{minipage}}
\vspace{0.5cm}
\caption{{Example~3. Condition numbers (and MINRES iteration counts) corresponding to the preconditioners (\eqref{precond.linear.B2})~$\mathcal{B}^{1,1}_{\mathsf{DFARD}}$ (top left), $\mathcal{B}^{1,2}_{\mathsf{DFARD}}$ (top right), $\mathcal{B}^{2,1}_{\mathsf{DFARD}}$ (bottom left), and $\mathcal{B}^{2,2}_{\mathsf{DFARD}}$ (bottom right), under different parameter regimes of the coupled Darcy--Forchheimer and ADR model.}}
\label{tab:DF-ARD-Prec}
\end{table}
% \begin{table}[t]\label{Prec:DFARD}
% \centering
% \small
% \setlength{\tabcolsep}{5pt}
% \renewcommand{\arraystretch}{1.15}

% \centering
% \begin{tabular}{|c |c|c| c|lll|}
% \hline
% $\rho_0$ & $\kappa$ & $\alpha$ & $\rF$ & \multicolumn{3}{c|}{$h$} \\
% \cline{5-7}
%  & & & & $2^{-2}$ & $2^{-3}$ & $2^{-4}$ \\
% \hline

% \multirow{3}{*}{$1$}
% & $10^{-8}$ & $10^{-9}$ & $10^{6}$ & 1.8 (6) & 1.8 (6)& 1.8 (7) \\
% & $10^{-4}$ & $10^{-5}$ & $1$ & 1.8 (10)& 1.8 (12) & 1.8 (14)\\
% & $1$ & $10^{-1}$ & $10^{-6}$ &2.9 (33) &3.0 (33) &3.0 (32) \\
% \hline

% \multirow{3}{*}{$10^{-5}$}
% & $10^{-8}$ & $10^{-9}$ & $10^{6}$ & 1.2 (8) &1.3 (10) &1.6 (14) \\
% & $10^{-4}$ & $10^{-5}$ & $1$ &3.6 (72) &4.0 (86) &4.1 (81) \\
% & $1$ & $10^{-1}$ & $10^{-6}$ &2.1 (19) &2.1 (19) &2.1 (19) \\
% \hline

% \multirow{3}{*}{$10^{-9}$}
% & $10^{-8}$ & $10^{-9}$ & $10^{6}$ & 3.6 (37)&4.0 (40) &4.3 (42) \\
% & $10^{-4}$ & $10^{-5}$ & $1$ &3.7 (25) &4.1 (30) &4.2 (32) \\
% & $1$ & $10^{-1}$ & $10^{-6}$ &2.1 (19) &2.1 (19) &2.1 (19) \\
% \hline
% \end{tabular}
% \caption{Condition numbers (and MINRES iteration counts) for the preconditioned system under different parameter regimes of the coupled Darcy--Forchheimer and ADR model corresponding to the preconditioner~$\mathcal{B}^1_{\mathsf{DFARD}}$ \eqref{precond.linear.B2}.}
% \label{tab:precond.B1.dfard}
% \end{table}
\begin{comment}
\begin{table}[t]\label{Prec:DFARD}
\centering
\small
\setlength{\tabcolsep}{5pt}
\renewcommand{\arraystretch}{1.15}

\centering
\begin{tabular}{|c |c|c| c|lll|}
\hline
$\rho_0$ & $\kappa$ & $\alpha$ & $\rF$ & \multicolumn{3}{c|}{$h$} \\
\cline{5-7}
 & & & & $2^{-2}$ & $2^{-3}$ & $2^{-4}$ \\
\hline

% \multirow{3}{*}{$1$}
% & $10^{-8}$ & 1 & $10^{-6}$ & 6.3 (39) & 6.5 (48)& 6.5 (14) \\
% & $10^{-4}$ & $10^{-4}$ & $1$ & 3.6 (28)& 3.8 (32) & 3.5 (37)\\
% & $1$ & $10^{-8}$ & $10^{6}$ &1.5 (12) &1.5 (12) &1.5 (12) \\
% \hline

% \multirow{3}{*}{$10^{-3}$}
% & $10^{-8}$ & 1 & $10^{-6}$ & 3.9 (37)&4.1 (42) &4.4 (46) \\
% & $10^{-4}$ & $10^{-4}$ & $1$ &4.1 (56) &4.0 (56) &3.7 (69) \\
% & $1$ & $10^{-8}$ & $10^{6}$ &2.1 (19) &3.0 (25) &3.7 (28) \\
% \hline

% \multirow{3}{*}{$10^{-6}$}
% & $10^{-8}$ & 1 & $10^{-6}$ & 3.9 (37)&4.1 (42) &4.4 (46) \\
% & $10^{-4}$ & $10^{-4}$ & $1$ &3.6 (39) &3.8 (44) &3.5 (48) \\
% & $1$ & $10^{-8}$ & $10^{6}$ &1.2 (11) &1.2 (11) &1.2 (11) \\
% \hline
% \multirow{3}{*}{$10^{-9}$}
% & $10^{-8}$ & 1 & $10^{-6}$ & 3.8 (37)&4.1 (42) &4.4 (46) \\
% & $10^{-4}$ & $10^{-4}$ & $1$ &4.1 (56) &4.0 (56) &3.7 (69) \\
% & $1$ & $10^{-8}$ & $10^{6}$ &4.1 (34) &4.3 (36) &4.5 (38) \\
% \hline
\multirow{3}{*}{$1$}
& $10^{-8}$ & $10^{-9}$ & $10^{6}$ & 6.3 (39) & 6.5 (48)& 6.5 (14) \\
& $10^{-4}$ & $10^{-5}$ & $1$ & 3.6 (28)& 3.8 (32) & 3.5 (37)\\
& $1$ & $10^{-1}$ & $10^{-6}$ &1.5 (12) &1.5 (12) &1.5 (12) \\
\hline

\multirow{3}{*}{$10^{-5}$}
& $10^{-8}$ & $10^{-9}$ & $10^{6}$ & 3.9 (37)&4.1 (42) &4.4 (46) \\
& $10^{-4}$ & $10^{-5}$ & $1$ &4.1 (56) &4.0 (56) &3.7 (69) \\
& $1$ & $10^{-1}$ & $10^{-6}$ &2.1 (19) &3.0 (25) &3.7 (28) \\
\hline

\multirow{3}{*}{$10^{-9}$}
& $10^{-8}$ & $10^{-9}$ & $10^{6}$ & 3.9 (37)&4.1 (42) &4.4 (46) \\
& $10^{-4}$ & $10^{-5}$ & $1$ &3.6 (39) &3.8 (44) &3.5 (48) \\
& $1$ & $10^{-1}$ & $10^{-6}$ &1.2 (11) &1.2 (11) &1.2 (11) \\
\hline
\end{tabular}
\caption{Condition numbers (and GMRES iteration counts) for the preconditioned system under different parameter regimes of the coupled Darcy--Forchheimer and ADR model corresponding to the preconditioner~$\mathcal{B}^2_{\mathrm{DFard}}$ \eqref{precond.linear.B2}.}
\label{tab:precond}
\end{table}
\end{comment}


\smallskip
\noindent\textbf{Example 4 (lid-driven cavity flow).} {Next, we present a numerical simulation following the classical lid-driven cavity problem for the Darcy--Forchheimer-transport coupling \cite{sayah2021finite}. The domain is $\Omega = (0,1)^2$ (discretised with a triangular mesh representing 82176 DoFs for the lowest-order method), the model parameters are taken as follows $\kappa = 1$, $\rF = 20$, $\alpha = 1$, $\bf(\varphi) = (10\varphi,10\varphi)^{\tt t}$, $m = 0$, $\varrho_0 = 0$, and as boundary conditions we impose $\bu \cdot \bn = 0$ everywhere on $\partial\Omega$ and so {for the sake of pressure uniqueness,  we impose $\int_\Omega p = 0$ (with a real Lagrange multiplier). We also} set $\varphi = 20$ on the top edge and $\varphi = 0$ on the remainder of the boundary. In the first two rows of Figure~\ref{fig:placeholder} we present the solution profiles, which match the results of \cite{sayah2021finite}. We also carried out an extension to 3D of the same test, taking $\Omega = (0,1)^3$ (discretised with a mesh representing 2323200 DoFs), $\bf(\varphi) = (10\varphi,10\varphi,10\varphi)^{\tt t}$, and maintained all other parameters unchanged. The obtained numerical solutions are shown in the bottom row of the same figure. For both cases we have used the lowest-order method, and the Newton--Raphson iterative scheme (using Hager--Zhang line search) took seven  iterations to reach the prescribed tolerance of $10^{-6}$. 
At least in this flow-transport regime and with the chosen parametric values, the discrepancies between 3D and 2D simulations are unnoticeable.} 


\begin{figure}[t!]
    \centering
    %\includegraphics[width=0.325\linewidth]{ex04u.png}
\includegraphics[width=0.245\linewidth]{ex04u-stream.pdf}    \includegraphics[width=0.245\linewidth]{ex04p.png}
  %  \includegraphics[width=0.325\linewidth]{ex04zeta.png}
\includegraphics[width=0.245\linewidth]{ex04zeta-stream.pdf}
\includegraphics[width=0.245\linewidth]{ex04phi.png}\\
\includegraphics[width=0.245\linewidth]{ex04u3d.png}
\includegraphics[width=0.245\linewidth]{ex04p3d.png}
\includegraphics[width=0.245\linewidth]{ex04zeta3d.png}
\includegraphics[width=0.245\linewidth]{ex04phi3d.png}
    \caption{Example 4. Numerical solutions (discharge velocity streamlines, pressure profile, total flux, and concentration) for the lid-driven cavity problem. Parameters are taken as in \cite{sayah2021finite} for the 2D case (top and middle rows), and we also show its extension to 3D (bottom row).}
    \label{fig:placeholder}
\end{figure}


\smallskip
\noindent\textbf{Example 5 (transient miscible displacement).} To conclude, we apply a time-dependent version of \eqref{eq:strong} (replacing the reaction term $\varrho_0\varphi$ by $\varrho_0\partial_t \varphi$ in the concentration equation and the diffusion $-\varepsilon \Delta \varphi$ by $-\varrho_0\varepsilon \Delta \varphi$, where $\varrho_0$ is the porosity) to the classical problem of viscous fingering in a Hele--Shaw cell. Here \(\varphi\) is interpreted as the solute concentration (in mol/m\(^3\)) of a resident fluid being displaced by a second, less viscous fluid (water). 


\begin{figure}[t!]
    \centering
    \includegraphics[width=0.245\linewidth]{ex05u-a.png}
    \includegraphics[width=0.245\linewidth]{ex05u-b.png}
   \includegraphics[width=0.245\linewidth]{ex05u-c.png}\\
   \includegraphics[width=0.245\linewidth]{ex05phi-a.png}
    \includegraphics[width=0.245\linewidth]{ex05phi-b.png}
   \includegraphics[width=0.245\linewidth]{ex05phi-c.png}
    \caption{Example 5. Snapshots of the approximate discharge flux (top) and concentration profile (bottom) at times $t=0.01$ (left), $t=6$ (centre), and $t=18$ (right).}
    \label{fig:ex05}
\end{figure}

The viscosity of the fluid mixture follows an Arrhenius law 
$\mu(\varphi) = \mu_1 \exp(R\,\frac{\varphi}{\varphi_2})$, 
where \(\mu_1\) is the viscosity of the displacing fluid (water), \(\varphi_2 = 6500\) mol/m\(^3\) is a reference concentration, and \(R = \ln r\) is the log‑mobility ratio. 
The viscosity ratio is \(r = \mu(\varphi=1)/\mu_1 = 2\), hence \(R = \ln 2\). 
The hydraulic conductivity is then \(\kappa = K/\mu(\varphi)\), with \(K\) the intrinsic permeability; the  body force \(\bf(\varphi)\) is set to zero. The computational domain is the channel \(\Omega = (0,0.03)\times(0,0.08)\) m\(^2\). 
The initial concentration field $\varphi(x,y,0)$ is a step function from 0 (water) on top of a parabolic profile located near the top boundary $y=0.0675+25(x-0.015)^2$, to $\varphi_2$ (resident fluid) below it. 
As water is injected from the top  boundary (setting a prescribed normal discharge inflow velocity $\bu\cdot\bn = -9\cdot10^{-4}$\,m/s and fixed concentration $\varphi = 0$), the resident fluid is displaced and the viscous fingering instability develops spontaneously. 
No artificial random perturbation is added; the unstructured mesh provides sufficient asymmetry to trigger the instability near the initial interface. On the vertical walls we impose \(\bu\cdot\bn = 0\) and \(\bzeta\cdot\bn = 0\). On the outlet (bottom segment) we set $p=0$ and zero diffusive flux (not total flux), this means that we need to include the remainder terms $-\varepsilon \langle \varphi, \bw\cdot \bn \rangle_{\mathrm{out}}$ and $\langle (\bzeta+\varphi\bu)\cdot \bn, \bw\cdot\bn\rangle_{\mathrm{out}}$ in equation \eqref{eq:weak,3}.  
    
The model parameters (diffusivity \(\alpha = 4\cdot10^{-8}\)\,m$^2/$s, porosity \(\varrho_0=0.5\), permeability \(K=10^{-6}\)\,m$^2$, reference concentration \(\varphi_2=6500\)\,mol/m$^3$,  the constant mobility ratio  \(R=2\), and water viscosity $\mu_1=1$\,mPa$\cdot$s) are taken from \cite{bear2013dynamics,gatica2022p}; the source term \(m\) in \eqref{eq:strong} is zero. 
The domain is discretised with an unstructured mesh of 37745 triangles. 
The time derivative is discretised using backward Euler's method with constant time step \(\Delta t = 0.01\)\,s, and the simulation runs until the final time \(t = 18\)\,s. 
The numerical results (snapshots of the approximate solutions at four different times, showing in this case only velocity and concentration)  are displayed in Figure~\ref{fig:ex05}. They reproduce the characteristic viscous fingering patterns observed in Hele--Shaw displacements. For this test the maximum number of nonlinear iterations  to convergence was four (and required at the initial time steps).



%%%%%%%%%%%%%%%%%%%%%%%%%%%%%%%%%
    \bibliographystyle{siam}
    \small 
\bibliography{refs}

%%%%%%%%%%%%%%%%%%%%%%%%%%%%%%%%%

\appendix 
\section{Preconditioners for sub-systems}
\smallskip
\noindent\textbf{Robust preconditioners for the Darcy--Forchheimer system.} 
Let us consider the Darcy--Forchheimer system alone from \eqref{eq:strong}
\begin{equation}\label{eq:strong.Darcy.Forchh}
% \left.
\begin{aligned}
    \kappa^{-1}\bu + {\rF|\bu|\bu} + \nabla p & = \bf \quad && \text{in }\Omega,\\
    \vdiv \bu & = 0 \quad && \text{in }\Omega,\\
    \bu \cdot \bn & = 0 \quad && \text{on }\Gamma_{\bu},\\
    p & = 0 \quad && \text{on }\Gamma_p.
\end{aligned}
% \right\}
\end{equation}
The following two \((\kappa, \rF, h)\)-robust preconditioners for \eqref{eq:strong.Darcy.Forchh} have been mentioned in \cite[Section~5]{dhpr_sinum26}, and the robustness of these preconditioners can be seen in Table~\ref{tab:DF.prec}.
\begin{align}\label{precond.linear}
\mathcal{B}_{\mathsf{DF}}^1:=\begin{pmatrix}
        (\kappa^{-1}\mathcal{I}+\rF 
        {\mathcal{H}({\bu})} {\mathcal{I}}-\nabla\vdiv)^{-1} & 0\\0 &\mathcal{I}^{-1}+ {\bigl(-\kappa\Delta\bigr)^{-1}-\bigl([\rF{\mathcal{H}({\bu})}]^{-1}{\Delta_{{\frac32}}}\bigr)^{-1}}
    \end{pmatrix},\end{align} 
    and \begin{align}\label{precond.linear.new}
    \mathcal{B}_{\mathsf{DF}}^2:=\begin{pmatrix}
        \bigl([\kappa^{-1}\mathcal{I}+\rF 
        {\mathcal{H}({\bu})}]\, ({\mathcal{I}}-\nabla\vdiv)\bigr)^{-1} & 0\\0 &([\kappa^{-1}\mathcal{I}+\rF 
        {\mathcal{H}({\bu})}]^{-1}\,\mathcal{I})^{-1}
    \end{pmatrix},\end{align}
    where \(\mathcal{H}(\bu)\) is the G\^ateaux derivative of the nonlinear operator \(|\bu|\bu\) in the state \(\bu\) in the direction of \(\delta\bu\), and ${\Delta_{{\frac32}}}$ is the $\frac32$-Laplace operator associated with the space {$\rW^{1,\frac32}_{\Gamma_p}(\Omega)$}, defined as  
  \[{\langle \Delta_{\frac32}\phi,\psi\rangle := \int_\Omega |\nabla\phi|^{-\frac12}\nabla \phi \cdot \nabla \psi, \qquad \forall \psi \in \rW^{1,3}(\Omega).} \] 
  
  \begin{table}[t!]
\centering
{\footnotesize \begin{minipage}{0.48\textwidth}
\centering
\begin{tabular}{|c | c | l l l|}
\hline
$\kappa$ & $\rF$ & \multicolumn{3}{c|}{$h$} \\
\cline{3-5}
 & & $2^{-2}$ & $2^{-3}$ & $2^{-4}$ \\
\hline

\multirow{3}{*}{$10^{-10}$}
& $10^{-10}$ & 3.7 (23) & 4.1 (31) & 4.4 (33) \\
& 1 & 5.7 (23) & 4.1 (31) & 4.4 (33) \\
& $10^{5}$ & 3.7 (23) & 4.1 (31) & 4.4 (33) \\

\hline

\multirow{3}{*}{$10^{-5}$}
& $10^{-10}$ & 3.7 (25) & 4.1 (31) & 4.2 (33) \\
& 1 & 3.7 (25) & 4.1 (31) & 4.2 (33) \\
& $10^{5}$ & 3.7 (27) & 4.1 (31) & 4.3 (33) \\

\hline

\multirow{3}{*}{$1$}
& $10^{-10}$ & 1.2 (11) & 1.2 (11) & 1.2 (11) \\
& 1 & 1.2 (11) & 1.2 (11) & 1.2 (11) \\
& $10^{5}$ & 4.4 (27) & 4.6 (33) & 4.5 (38) \\

\hline
\end{tabular}
% \caption{Performance of \eqref{precond.linear}}
\end{minipage}
\hfill
\begin{minipage}{0.48\textwidth}
\centering
\begin{tabular}{|c | c | l l l|}
\hline
$\kappa$ & $\rF$ & \multicolumn{3}{c|}{$h$} \\
\cline{3-5}
 & & $2^{-2}$ & $2^{-3}$ & $2^{-4}$ \\
\hline

\multirow{3}{*}{$10^{-10}$}
& $10^{-10}$ & 1.2 (1) & 1.2 (1) & 1.2 (1) \\
& $1$        & 1.2 (1) & 1.2 (1) & 1.2 (1) \\
& $10^{5}$   & 1.2 (1) & 1.2 (1) & 1.2 (1) \\

\hline

\multirow{3}{*}{$10^{-5}$}
& $10^{-10}$ & 1.2 (6) & 1.2 (6) & 1.2 (6) \\
& $1$        & 1.2 (6) & 1.2 (6) & 1.2 (6) \\
& $10^{5}$   & 1.2 (8) & 1.2 (8) & 1.2 (8) \\

\hline

\multirow{3}{*}{$1$}
& $10^{-10}$ & 1.2 (6) & 1.2 (6) & 1.2 (6) \\
& $1$        & 1.2 (8) & 1.2 (8) & 1.2 (8) \\
& $10^{5}$   & 1.4 (13) & 1.4 (13) & 1.3 (12) \\

\hline
\end{tabular}
% \caption{Performance of  \eqref{precond.linear.new}}
\end{minipage}
}

\caption{Condition numbers of the preconditioned systems (\eqref{precond.linear} on left, and \eqref{precond.linear.new} on right) with MINRES iterations in parentheses for different values of $\kappa, \rF,$ and $h$.}
\label{tab:DF.prec}
\end{table}

\smallskip
\noindent\textbf{A robust preconditioner for the reaction-diffusion system.}
Consider the mixed reaction-diffusion system 
\begin{equation}\label{eq:strong.reac.diff}
%\left.
\begin{aligned}
    \alpha^{-1}\bzeta - \nabla \varphi & = \bzero \quad && \text{in }\Omega,\\
    \vdiv \bzeta - \varrho_0 \varphi & = -m \quad && \text{in }\Omega,\\
    \varphi & = 0 \quad && \text{on }\Gamma_{\bu},\\
    \bzeta \cdot \bn & = 0 \quad && \text{on }\Gamma_p.
\end{aligned}
%\right\}
\end{equation}
The natural choice of a \((\alpha,\varphi,h)\)-robust preconditioner for \eqref{eq:strong.reac.diff} (following \cite[Equation~3.8]{boon2021robust}) is the following 
  \begin{equation}\label{abstract-prec-reac.diff} 
\mathcal{B}_{{\mathrm{RD}}}:= \begin{pmatrix} (\frac{1}{\alpha}\cI - \nabla(\vdiv) )^{-1} & 0 \\ 0 & ([1+\varrho_0]\,\mathcal{I})^{-1}+(\varrho_0\mathcal{I}- {\alpha}\Delta )^{-1}
\end{pmatrix}.
\end{equation}
The condition numbers and the MINRES iteration counts are reported in Table~\ref{tab:reac.diff}.

\begin{table}[t!]
\centering
{\footnotesize \begin{tabular}{|c | c | l l l|}
\hline
$\alpha$ & $\varrho_0$ & \multicolumn{3}{c|}{$h$} \\
\cline{3-5}
 & & $2^{-2}$ & $2^{-3}$ & $2^{-4}$ \\
\hline

\multirow{4}{*}{$10^{-9}$}
& $10^{-9}$ & 3.7 (17) & 4.1 (19) & 4.4 (21) \\
& $10^{-6}$ & 1.2 (10) & 1.7 (13) & 2.5 (18) \\
& $10^{-3}$ & 1.0 (4)  & 1.0 (4)  & 1.0 (4)  \\
& $1$       & 1.5 (3)  & 1.5 (3)  & 1.5 (3)  \\

\hline

\multirow{4}{*}{$10^{-3}$}
& $10^{-9}$ & 3.1 (24) & 3.3 (27) & 3.4 (29) \\
& $10^{-6}$ & 3.1 (24) & 3.3 (27) & 3.4 (29) \\
& $10^{-3}$ & 3.0 (23) & 3.0 (25) & 3.1 (29) \\
& $1$       & 1.9 (12) & 2.3 (15) & 2.6 (19) \\

\hline

\multirow{4}{*}{$1$}
& $10^{-9}$ & 1.2 (11) & 1.2 (11) & 1.2 (11) \\
& $10^{-6}$ & 1.2 (11) & 1.2 (11) & 1.2 (11) \\
& $10^{-3}$ & 1.2 (11) & 1.2 (11) & 1.2 (11) \\
& $1$       & 2.6 (14) & 2.6 (14) & 2.6 (14) \\

\hline
\end{tabular}}
\caption{Condition numbers of the preconditioned system (corresponding to \ref{abstract-prec-reac.diff}) with MINRES iterations in parentheses for different parameters and mesh sizes.}\label{tab:reac.diff}
\end{table}




\begin{table}[t!]\label{tab:advec.reac.diff}
\centering
{\footnotesize \begin{tabular}{|c|c|c|lll|}
\hline
$|\hat\bu|$ & $\varrho_0$ & $\alpha$ & \multicolumn{3}{c|}{$h$} \\
\cline{4-6}
& & & $2^{-2}$ & $2^{-3}$ & $2^{-4}$ \\
\hline
\multirow{12}{*}{$10^{-10}$}
& \multirow{4}{*}{$10^{-8}$}
& $10^{-9}$ & 3.3 (26) & 3.9 (31) & 4.2 (32) \\
& & $10^{-6}$ & 3.7 (19) & 4.2 (26) & 4.4 (28) \\
& & $10^{-3}$ & 3.1 (20)  & 3.3 (26)  & 3.4 (28)  \\
& & $1$       & 1.2 (11)  & 1.2 (11)  & 1.2 (11)  \\

\cline{2-6}

& \multirow{4}{*}{$10^{-4}$}
& $10^{-9}$ & 1.0 (5) & 1.0 (6) & 1.0 (7) \\
& & $10^{-6}$ & 2.2 (20) & 3.0 (29) & 3.7 (33) \\
& & $10^{-3}$ & 3.1 (20)  & 3.2 (25)  & 3.3 (28)  \\
& & $1$       & 1.2 (11)  & 1.2 (11)  & 1.2 (11)  \\

\cline{2-6}

& \multirow{4}{*}{1}
& $10^{-9}$ & 1.5 (3) & 1.5 (3) & 1.5 (3) \\
& & $10^{-6}$ & 1.5 (3) & 1.5 (3) & 1.3 (5) \\
& & $10^{-3}$ & 1.9 (11)  & 2.6 (14)  & 2.6 (18)  \\
& & $1$       & 2.6 (13)  & 2.6 (13)  & 2.6 (13)  \\

\cline{1-6}

\multirow{12}{*}{$10^{-5}$}
& \multirow{4}{*}{$10^{-8}$}
& $10^{-9}$ & 1567.6 (63) & 2232.1 (248) & 3189.9 ($\star$) \\
& & $10^{-6}$ & 64.1 (42) & 55.2(74) & 6.0 (61) \\
& & $10^{-3}$ & 3.1 (20) & 3.3 (26) & 3.3 (29) \\
& & $1$       & 1.2 (11) & 1.2 (11) & 1.2 (11) \\

\cline{2-6}

& \multirow{4}{*}{$10^{-4}$}
& $10^{-9}$ & 3.9 (32) & 64.6 (128)& 105.6 (521) \\
& & $10^{-6}$ & 3.3 (31) & 4.8 (48) & 1.5 (5) \\
& & $10^{-3}$ & 3.1(20) & 3.2 (26) & 3.3 (29) \\ 
& & $1$       & 1.2 (11) & 1.2 (11) & 1.2 (11) \\
\cline{2-6}

& \multirow{4}{*}{1}
& $10^{-9}$ & 1.5 (3) & 1.5 (4) & 1.5 (5) \\
& & $10^{-6}$ & 1.5 (3) & 1.5 (5) & 1.5 (5) \\
& & $10^{-3}$ & 1.9 (11)  & 2.3 (14)  & 2.6 (18)  \\
& & $1$       & 2.6 (13)  & 2.6 (13)  & 2.6 (13)  \\


\cline{1-6}

\multirow{12}{*}{$10^{-1}$}
& \multirow{4}{*}{$10^{-8}$}
& $10^{-9}$ & 1.4e+07 (61) & 1.1e+07 (243) & 8.6e+06 (991) \\
& & $10^{-6}$ & 22040.3 (64) & 17537.9 (250) & 12259.2 (998) \\
& & $10^{-3}$ & 21.1 (65) & 190.7 (230) & 356.6 (377) \\
& & $1$       & 1.2 (11) & 1.2 (11) & 1.2 (11) \\

\cline{2-6}

& \multirow{4}{*}{$10^{-4}$}
& $10^{-9}$ & 24970.7 (62) & 30598.7 (245) & 44066.2 (988) \\
& & $10^{-6}$ &3711.7 (64) & 4484.1 (248) & 3435.6 (996) \\
& & $10^{-3}$ & 21.0 (65) & 186.5 (230) & 345.8 (377) \\
& & $1$       & 1.2 (11) & 1.2 (11) & 1.2 (11) \\

\cline{2-6}

& \multirow{4}{*}{1}
& $10^{-9}$ & 1.9 (43) & 7.4 (143) & 9.8 (494) \\
& & $10^{-6}$ & 1.9 (43) & 7.5 (143) & 9.4 (491) \\
& & $10^{-3}$ & 1.8 (39)  & 6.0 (84)  & 18.2 (126)  \\
& & $1$       & 2.6 (14)  & 2.5 (14)  & 2.5 (13)  \\

\hline
\end{tabular}}
\caption{Condition numbers and MINRES iteration counts in parenthesis for different values of \(\alpha\) and mesh sizes $h$ corresponding to the preconditioner~ \eqref{Prec:B2.advec.diff}. The symbol $\star$ indicates that the MINRES iteration count exceeds 1000.}\label{tab:advec.reac.diff}
\end{table}


\smallskip
\noindent\textbf{Robust preconditioners for the advection-diffusion-reaction system.}
Consider now the ADR system 
\begin{equation}\label{eq:strong.advection.diff}
%\left.
\begin{aligned}
    \alpha^{-1}\bzeta - \nabla \varphi + \alpha^{-1}\hat\bu\,\varphi & = \bzero \quad && \text{in }\Omega,\\
    \vdiv \bzeta -\varrho_0\varphi& = -m \quad && \text{in }\Omega,\\
    \varphi & = 0 \quad && \text{on }\Gamma_{\bu},\\
    \bzeta \cdot \bn & = 0 \quad && \text{on }\Gamma_p.
\end{aligned}
%\right\}
\end{equation}
The preconditioners~\eqref{precond.linear.B2} for the full system already include the blocks needed for the sub-system~\eqref{eq:strong.advection.diff}.  Here, we vary the magnitude of the Darcy--Forchheimer velocity~$\hat{\bu}$ to illustrate the deterioration in performance of these preconditioners in the advection-dominated regime. Specifically, we restrict our attention to $\mathrm{diag}({\cR}_2,{\cS}_2)$ from \eqref{precond.linear.B2}, noting that the remaining preconditioners exhibit similar behaviour: 
\begin{equation}\label{Prec:B2.advec.diff}
    \mathcal{B}^{1,2}_{\mathsf{DFARD}}:=\begin{pmatrix}
        (\alpha^{-1}\cI+\nabla\vdiv)^{-1}&0\\0&{(1+\varrho_0)\,\cI^{-1}+(\varrho_0\mathcal{I}-\alpha\Delta + \hat{\bu}\cdot \nabla)^{-1}}
    \end{pmatrix}.
\end{equation}
{A conclusion from Table~\ref{tab:advec.reac.diff} is that whenever $\frac{|\hat{\mathbf{u}}|h}{\alpha} \ll 1$, both the condition numbers of the preconditioned system and the MINRES iteration counts remain well controlled. Outside this regime, however, the performance of the preconditioner deteriorates and becomes inconsistent.
Possible remedies include a better stabilisation for the advection term~$\bu^m \varphi$ (pure upwind in general will not yield robust preconditioners for extreme regimes).
It is well established \cite{ernst2000residual} that classical discretisations of convection-dominated systems give rise to highly nonnormal discrete operators. Such nonnormality is not adequately reflected by the eigenvalues of the preconditioned systems alone; instead, their behavior deteriorates markedly  by increasing  P\'eclet number and mesh refinement. As a consequence, Krylov subspace methods will struggle for this type of problems.} 

Another difficulty already mentioned in Section \ref{sec:precond} is the nonsymmetric nature of the underlying discrete operators. Operator preconditioning and the usual  Schur complement structures are not yet  systematically developed. Here we have adopted a structure motivated by symmetric elliptic operators, which do not seem to provide definitive or robust conclusions.
See \cite{wathen2015preconditioning} for a detailed review of the challenges associated with preconditioning nonsymmetric systems. 





\begin{comment}
\begin{table}[t!]\label{tab:prec.triang.0}
\centering
\begin{tabular}{|c|c|ccc|}
\hline
$\alpha$ & $\varrho_0$ & &$h$ &  \\
\cline{3-5}
 & & $2^{-2}$ & $2^{-3}$ & $2^{-4}$ \\
\hline

$10^{-9}$ & $10^{-3}$ & 21 & 82 & 281 \\
$10^{-9}$ & $10^{-2}$ & 22 & 89 & 338 \\
$10^{-9}$ & $10^{-1}$ & 20 & 54 & 165 \\
$10^{-9}$ & $1$       & 8  & 10 & 13 \\

\hline

$10^{-6}$ & $10^{-3}$ & 23 & 83 & 299 \\
$10^{-6}$ & $10^{-2}$ & 25 & 96 & 364 \\
$10^{-6}$ & $10^{-1}$ & 20 & 61 & 190 \\
$10^{-6}$ & $1$       & 9  & 11 & 16 \\

\hline

$10^{-3}$ & $10^{-3}$ & 34 & 50 & 68 \\
$10^{-3}$ & $10^{-2}$ & 32 & 44 & 63 \\
$10^{-3}$ & $10^{-1}$ & 24 & 31 & 47 \\
$10^{-3}$ & $1$       & 15 & 21 & 31 \\

\hline

$1$ & $10^{-3}$ & 35 & 57 & 84 \\
$1$ & $10^{-2}$ & 35 & 57 & 85 \\
$1$ & $10^{-1}$ & 35 & 57 & 85 \\
$1$ & $1$       & 35 & 57 & 86 \\

\hline
\end{tabular}
\caption{GMRES iterations for varying $\alpha$, $\varrho_0$, and mesh size $h$ corresponding to \eqref{Prec:triang.0}.}
\end{table}
\end{comment}

\begin{comment}
Another block-triangular preconditioner that performs better than the discussed block-diagonal variants (see Table~\ref{tab:prec.triang.0}), and does not rely on any smallness assumption on the advective velocity, is given by
\begin{align}\label{Prec:triang.0}
\mathcal{P}
:=
\begin{pmatrix}
(\alpha^{-1}\mathcal{I}-\nabla\!\vdiv) & \bbB^*+\bbD \\
0 & (1+\varrho_0)^{-1}\mathcal{I}
      + (\varrho_0\mathcal{I}-\alpha\Delta+\hat{\bu}\cdot\nabla)^{-1}
\end{pmatrix}^{-1}.
\end{align}
Several variants of this triangular preconditioner can be constructed; however, they exhibit performance similar to \eqref{Prec:triang.0} (see, for instance, \cite{olshanskii2007pressure}).
\end{comment}





\end{document}
